% L'alternateur — Physique 1ère TI — Cours signé OnBuch+
% Programme officiel MINESEC. Compiler avec : tectonic N17-alternateur.tex
\newcommand{\DOCMATIERE}{Physique}
\newcommand{\DOCNIVEAU}{1ère TI}
\newcommand{\DOCMODULE}{Module 4 : Aspects énergétiques des circuits électriques}
\newcommand{\DOCLECON}{N17}
\newcommand{\DOCTITRE}{L'alternateur}
% OnBuch+ — préambule « cours signés » ENRICHI (compilé avec tectonic / XeLaTeX)
% Système visuel « Pop » d'OnBuch+ : fond crème (jamais blanc), bordures encre
% épaisses, ombres « dures » décalées, titres Archivo Black en capitales.
% Version MATHÉMATIQUES : graphes (pgfplots/tikz), boîtes pédagogiques
% (définition, propriété, méthode, attention, le savais-tu, exercice résolu,
% exercice, TP, bilan), page de garde et signature OnBuch+ ancrée en bas.
%
% Macros attendues dans main.tex AVANT \input{preamble} :
%   \DOCMATIERE  \DOCNIVEAU  \DOCMODULE  \DOCLECON (numéro)  \DOCTITRE
\documentclass[11pt]{article}

\usepackage{fontspec}
\usepackage[french]{babel}
\usepackage[a4paper,margin=2cm,top=3.4cm,bottom=2.8cm,headheight=1.4cm,footskip=1.3cm]{geometry}
\usepackage{amsmath,amssymb,amsfonts}
\usepackage{mathtools} % \xmapsto, \coloneqq... souvent utilisés par les modèles
\usepackage{cancel} % simplifications barrées (bilans de forces)
% \abs{z} (module, valeur absolue) et \norm{u} (norme), utilisés par les modèles
\providecommand{\abs}{}\let\abs\relax\DeclarePairedDelimiter{\abs}{\lvert}{\rvert}
\providecommand{\norm}{}\let\norm\relax\DeclarePairedDelimiter{\norm}{\lVert}{\rVert}
\usepackage[table]{xcolor} % [table] : fournit aussi \rowcolors (alternance de lignes)
\usepackage{pagecolor}
\usepackage{tikz}
\usetikzlibrary{arrows.meta,positioning,calc,shapes.geometric,shapes.misc,shapes.arrows,decorations.pathmorphing,decorations.pathreplacing,decorations.markings,patterns,shadows,mindmap,trees,fit,backgrounds,matrix,intersections,angles,quotes}
\usepackage{pgfplots}
\usepackage[version=4]{mhchem} % équations nucléaires \ce{^{235}_{92}U}
\usepackage[european,siunitx]{circuitikz} % schémas électriques
\pgfplotsset{compat=1.18}
\usepgfplotslibrary{fillbetween,groupplots} % aires entre courbes (calcul intégral)
\usepackage{enumitem}
\usepackage{fancyhdr}
% [most] charge toutes les bibliothèques tcolorbox (listings, minted, raster,
% documentation...) et gonfle inutilement la mémoire TeX sur un document long
% (~300 boîtes) : on ne charge que ce qui est réellement utilisé.
\usepackage[skins,breakable]{tcolorbox}
\usepackage{graphicx}
\usepackage{colortbl}
\usepackage{booktabs}
\usepackage{tabularx}
\usepackage{diagbox} % case d'en-tête diagonale des tableaux à double entrée
% Colonnes extensibles centrée / gauche / droite (C, L, R), souvent utilisées par les modèles
\newcolumntype{C}{>{\centering\arraybackslash}X}
\newcolumntype{L}{>{\raggedright\arraybackslash}X}
\newcolumntype{R}{>{\raggedleft\arraybackslash}X}
\usepackage{array}
\usepackage{multicol}
\usepackage{siunitx}
% parse-numbers=false : accepte aussi \SI{2,5 \times 10^{-3}}{...}
\sisetup{locale=FR,output-decimal-marker={,},group-minimum-digits=4,parse-numbers=false}
\DeclareSIUnit\ohm{\ensuremath{\symup{\Omega}}} % U+2126 absent de la police
\DeclareSIUnit\division{div} % réglages d'oscilloscope (ms/div, V/div)
\DeclareSIUnit\tr{tr} % tours (vitesse de rotation, ex. \SI{1500}{\tr\per\minute})
\DeclareSIUnit\ch{ch} % cheval-vapeur (puissance moteur, unité usuelle hors SI)
\DeclareSIUnit\franccfa{FCFA} % franc CFA (coûts, contexte camerounais)
\DeclareSIUnit\dioptre{\ensuremath{\delta}} % vergence d'une lentille (dioptrie)
\usepackage{qrcode}
% \degree : commande fréquemment utilisée par les modèles pour un angle ou une
% température en dehors de \SI{}{} (non fournie par les paquets chargés).
\providecommand{\degree}{^{\circ}}
% \qed (fin de démonstration), utilisé par les modèles sans amsthm
\providecommand{\qed}{\hfill$\square$}
% \barre : double barre des valeurs interdites dans un tableau de variations (tkz-tab)
\providecommand{\barre}{\ensuremath{\|}}
% environnement proof (amsthm non chargé) utilisé par les modèles
\ifdefined\proof\else\newenvironment{proof}[1][Démonstration]{\par\noindent\textit{#1.}\ }{\qed\par}\fi
\usepackage{lastpage}
\usepackage{hyperref}
\hypersetup{hidelinks}

% ============================================================
% Palette « Pop » OnBuch+ — mêmes hex que la classe P (plus_ui.dart)
% ============================================================
\definecolor{poporange}{HTML}{F59321}
\definecolor{popdark}{HTML}{E07A0C}
\definecolor{popcream}{HTML}{FFF6E8}
\definecolor{popink}{HTML}{1C1714}
\definecolor{poppurple}{HTML}{7A5AE0}
\definecolor{popgreen}{HTML}{2E9E6A}
\definecolor{poppink}{HTML}{E0457B}
\definecolor{popblue}{HTML}{2F7DE1}
\definecolor{popgold}{HTML}{C98A12}
\definecolor{popmuted}{HTML}{8A7F76}
\definecolor{popline}{HTML}{EADFD2}
\definecolor{popgray}{HTML}{8A7F76}
\colorlet{popgrey}{popgray}
% Alias tolérants (le modèle nomme parfois les couleurs différemment)
\colorlet{popwhite}{white}
\colorlet{popblack}{popink}
\colorlet{popred}{poppink}
\colorlet{popyellow}{popgold}
% Teintes claires pour fonds de tableaux / zones
\colorlet{popblueL}{popblue!10!white}
\colorlet{poporangeL}{poporange!14!white}
\colorlet{popgreenL}{popgreen!12!white}
\colorlet{poppurpleL}{poppurple!12!white}
\colorlet{poppinkL}{poppink!12!white}
\colorlet{popgoldL}{popgold!14!white}

\pagecolor{popcream}

% ============================================================
% Typographie
% ============================================================
\defaultfontfeatures{Ligatures=TeX}
\setmainfont{PlusJakartaSans-Medium.ttf}[Path=../../../fonts/,BoldFont=PlusJakartaSans-Bold.ttf]
% Maths en sans-serif (Fira Math) avec chiffres et lettres droites de Plus
% Jakarta, pour que les formules se fondent dans le texte.
\usepackage[math-style=ISO,bold-style=ISO]{unicode-math}
\setmathfont{FiraMath-Regular.otf}[Path=../../../fonts/]
\setmathfont{PlusJakartaSans-Medium.ttf}[Path=../../../fonts/,range={up/num,up/{Latin,latin},\mathup}]
\usepackage{icomma} % virgule décimale sans espace en mode math
% Caractères mathématiques Unicode absents des polices de texte (Plus Jakarta,
% Archivo Black) : rendus par leur équivalent mathématique, en texte comme en
% formule — sinon ils s'affichent en carré vide (ex. « Arithmétique dans ℤ »).
\usepackage{newunicodechar}
\newunicodechar{ℝ}{\ensuremath{\mathbb{R}}}
\newunicodechar{ℤ}{\ensuremath{\mathbb{Z}}}
\newunicodechar{ℂ}{\ensuremath{\mathbb{C}}}
\newunicodechar{ℕ}{\ensuremath{\mathbb{N}}}
\newunicodechar{ℚ}{\ensuremath{\mathbb{Q}}}
\newunicodechar{𝔻}{\ensuremath{\mathbb{D}}}
\newunicodechar{α}{\ensuremath{\alpha}}
\newunicodechar{β}{\ensuremath{\beta}}
\newunicodechar{γ}{\ensuremath{\gamma}}
\newunicodechar{δ}{\ensuremath{\delta}}
\newunicodechar{ε}{\ensuremath{\varepsilon}}
\newunicodechar{θ}{\ensuremath{\theta}}
\newunicodechar{λ}{\ensuremath{\lambda}}
\newunicodechar{μ}{\ensuremath{\mu}}
\newunicodechar{σ}{\ensuremath{\sigma}}
\newunicodechar{τ}{\ensuremath{\tau}}
\newunicodechar{φ}{\ensuremath{\varphi}}
\newunicodechar{ψ}{\ensuremath{\psi}}
\newunicodechar{ω}{\ensuremath{\omega}}
\newunicodechar{Σ}{\ensuremath{\Sigma}}
% majuscules produites par \MakeUppercase dans les titres (α-aminés -> Α)
\newunicodechar{Α}{\ensuremath{\alpha}}
\newunicodechar{Β}{\ensuremath{\beta}}
\newunicodechar{Γ}{\ensuremath{\Gamma}}
\newunicodechar{Δ}{\ensuremath{\Delta}}
\newunicodechar{Ε}{\ensuremath{\varepsilon}}
\newunicodechar{Θ}{\ensuremath{\Theta}}
\newunicodechar{Λ}{\ensuremath{\Lambda}}
\newunicodechar{Μ}{\ensuremath{\mu}}
\newunicodechar{Τ}{\ensuremath{\tau}}
\newunicodechar{Φ}{\ensuremath{\Phi}}
\newunicodechar{Ψ}{\ensuremath{\Psi}}
\newunicodechar{Ω}{\ensuremath{\Omega}}
\newunicodechar{↦}{\ensuremath{\mapsto}}
\newunicodechar{∈}{\ensuremath{\in}}
\newunicodechar{∉}{\ensuremath{\notin}}
\newunicodechar{∧}{\ensuremath{\wedge}}
\newunicodechar{⇔}{\ensuremath{\Leftrightarrow}}
\newunicodechar{⟺}{\ensuremath{\Longleftrightarrow}}
\newunicodechar{⇒}{\ensuremath{\Rightarrow}}
\newunicodechar{⟹}{\ensuremath{\Longrightarrow}}
\newunicodechar{∘}{\ensuremath{\circ}}
\newunicodechar{∪}{\ensuremath{\cup}}
\newunicodechar{∩}{\ensuremath{\cap}}
\newunicodechar{≡}{\ensuremath{\equiv}}
\newunicodechar{⊂}{\ensuremath{\subset}}
\newunicodechar{⊥}{\ensuremath{\perp}}
\newunicodechar{∃}{\ensuremath{\exists}}
\newunicodechar{∀}{\ensuremath{\forall}}
\newunicodechar{∛}{\ensuremath{\sqrt[3]{\,}}}
\newunicodechar{∜}{\ensuremath{\sqrt[4]{\,}}}
\newunicodechar{⃗}{\ensuremath{{}^{\to}}}
\newunicodechar{ⁿ}{\ifmmode^{n}\else\textsuperscript{\ensuremath{n}}\fi}
\newunicodechar{ˣ}{\ifmmode^{x}\else\textsuperscript{\ensuremath{x}}\fi}
\newunicodechar{ᵇ}{\ifmmode^{b}\else\textsuperscript{\ensuremath{b}}\fi}
\newunicodechar{ʳ}{\ifmmode^{r}\else\textsuperscript{\ensuremath{r}}\fi}
\newunicodechar{ᵘ}{\ifmmode^{u}\else\textsuperscript{\ensuremath{u}}\fi}
\newunicodechar{ʸ}{\ifmmode^{y}\else\textsuperscript{\ensuremath{y}}\fi}
\newunicodechar{ᵃ}{\ifmmode^{a}\else\textsuperscript{\ensuremath{a}}\fi}
\newunicodechar{ᵉ}{\ifmmode^{e}\else\textsuperscript{\ensuremath{e}}\fi}
\newunicodechar{ᵏ}{\ifmmode^{k}\else\textsuperscript{\ensuremath{k}}\fi}
\newunicodechar{ᵖ}{\ifmmode^{p}\else\textsuperscript{\ensuremath{p}}\fi}
\newunicodechar{ᴺ}{\ifmmode^{N}\else\textsuperscript{\ensuremath{N}}\fi}
\newunicodechar{ᶜ}{\ifmmode^{c}\else\textsuperscript{\ensuremath{c}}\fi}
\newunicodechar{ʰ}{\ifmmode^{h}\else\textsuperscript{\ensuremath{h}}\fi}
\newunicodechar{ᵠ}{\ifmmode^{q}\else\textsuperscript{\ensuremath{q}}\fi}
\newunicodechar{ₙ}{\ifmmode_{n}\else\textsubscript{\ensuremath{n}}\fi}
\newunicodechar{ₐ}{\ifmmode_{a}\else\textsubscript{\ensuremath{a}}\fi}
\newunicodechar{ᵢ}{\ifmmode_{i}\else\textsubscript{\ensuremath{i}}\fi}
\newunicodechar{ᵣ}{\ifmmode_{r}\else\textsubscript{\ensuremath{r}}\fi}
\newunicodechar{ₖ}{\ifmmode_{k}\else\textsubscript{\ensuremath{k}}\fi}
\newunicodechar{ᵦ}{\ifmmode_{\beta}\else\textsubscript{\ensuremath{\beta}}\fi}
\newunicodechar{ₓ}{\ifmmode_{x}\else\textsubscript{\ensuremath{x}}\fi}
\newfontfamily\popdisplay{ArchivoBlack-Regular.ttf}[Path=../../../fonts/]
\newfontfamily\poplogo{SpaceGrotesk-Bold.ttf}[Path=../../../fonts/]
\newfontfamily\popsemi{PlusJakartaSans-SemiBold.ttf}[Path=../../../fonts/]
\newfontfamily\popxbold{PlusJakartaSans-ExtraBold.ttf}[Path=../../../fonts/]
\newfontfamily\popxboldls{PlusJakartaSans-ExtraBold.ttf}[Path=../../../fonts/,LetterSpace=3.0]

\setlist[enumerate]{leftmargin=1.7em,itemsep=5pt,topsep=4pt}
\setlist[itemize]{leftmargin=1.7em,itemsep=4pt,topsep=4pt,label={\color{poporange}\textbullet}}
\setlength{\parindent}{0pt}
\setlength{\parskip}{5pt}
\renewcommand{\arraystretch}{1.25}

% pgfplots : style Pop par défaut
\pgfplotsset{
  every axis/.append style={axis line style={popink,line width=1pt},tick style={popink},
    grid style={popline,line width=0.5pt},label style={font=\small},tick label style={font=\footnotesize}},
  popcurve/.style={poporange,line width=1.8pt,no markers,smooth},
}
% Même style disponible dans un \draw TikZ simple (hors pgfplots)
\tikzset{popcurve/.style={poporange,line width=1.8pt}}
\tikzset{popgrid/.style={popline,very thin}}
\colorlet{popcurve}{poporange} % le nom du style est parfois utilisé comme couleur

% ============================================================
% Pill / squiggle
% ============================================================
\newcommand{\poppill}[3]{%
  \tikz[baseline=-0.55ex]\node[draw=popink,line width=1.2pt,rounded corners=2.6pt,%
    fill=#2,inner xsep=6.5pt,inner ysep=3pt]{%
    \popxboldls\fontsize{7}{7}\selectfont\color{#3}\MakeUppercase{#1}};%
}
\newcommand{\popsquiggle}[1]{%
  \tikz[baseline=-0.4ex]{
    \draw[#1,line width=2.6pt,line cap=round]
      plot [smooth] coordinates {(0,0.09) (0.34,0.01) (0.68,0.21) (1.02,0.01) (1.36,0.10)};
  }%
}
% Mot-clé surligné (marqueur orange sous le texte)
\newcommand{\cle}[1]{{\popxbold\color{popdark}#1}}

% ============================================================
% En-tête / pied
% ============================================================
\pagestyle{fancy}
\fancyhf{}
\renewcommand{\headrulewidth}{0pt}
\renewcommand{\footrulewidth}{0pt}
\lhead{\raisebox{-0.15ex}{\poplogo\fontsize{13}{13}\selectfont\color{popink}OnBuch\color{poporange}+}}
\rhead{\poppill{\DOCMATIERE\ \textbullet\ \DOCNIVEAU\ \textbullet\ Leçon \DOCLECON}{white}{popink}}
\lfoot{\popxbold\fontsize{7.6}{7.6}\selectfont\color{popmuted}\MakeUppercase{Cours signé OnBuch+}\ \textbullet\ {\popsemi\DOCTITRE}}
\rfoot{\tikz[baseline=-0.6ex]\node[rounded corners=8.5pt,fill=popink,minimum height=17pt,minimum width=17pt,inner xsep=5pt,inner sep=0]{\popxbold\fontsize{8}{8}\selectfont\color{popcream}\thepage\,/\,\pageref*{LastPage}};}

% ============================================================
% Titres
% ============================================================
\newcommand{\chaptitre}[2]{%
  \begin{tcolorbox}[enhanced,colback=poporange,colframe=popink,boxrule=2.6pt,arc=11pt,
      drop shadow=popink,left=15pt,right=15pt,top=12pt,bottom=11pt,before skip=3pt,after skip=18pt]
    \poppill{#2}{popink}{popcream}\par\vspace{9pt}
    {\raggedright\hyphenpenalty=10000\exhyphenpenalty=10000\popdisplay\fontsize{22}{24}\selectfont\color{popcream}\MakeUppercase{#1}\par}
  \end{tcolorbox}
}
% Section numérotée : pastille orange avec le numéro + titre Archivo
\newcounter{popsec}
\newcommand{\coursec}[1]{%
  \stepcounter{popsec}\setcounter{popsub}{0}%
  \par\vspace{16pt}\noindent
  \begin{minipage}{\linewidth}
  \tikz[baseline=-0.7ex]\node[draw=popink,line width=1.6pt,fill=poporange,rounded corners=4pt,minimum size=22pt,inner sep=0]{\popdisplay\fontsize{12}{12}\selectfont\color{popcream}\thepopsec};%
  \hspace{8pt}{\popdisplay\fontsize{15}{17}\selectfont\color{popink}\MakeUppercase{#1}}\par
  \vspace{4pt}\noindent{\color{poporange}\rule{36pt}{3.6pt}}
  \end{minipage}\par\nopagebreak\vspace{8pt}
}
\newcounter{popsub}
\newcommand{\courssub}[1]{%
  \stepcounter{popsub}%
  \par\vspace{9pt}\noindent{\popxbold\fontsize{11.5}{13}\selectfont\color{popink}{\color{poporange}\thepopsec.\thepopsub}\hspace{6pt}#1}\par\nopagebreak\vspace{4pt}
}

% ============================================================
% Boîtes pédagogiques — même recette : fond blanc, bordure épaisse colorée,
% ombre dure encre, badge plein. Titre optionnel [#1].
% ============================================================
\tcbset{popbase/.style={breakable,enhanced,colback=white,boxrule=2.3pt,arc=9pt,
  shadow={3pt}{-3pt}{0pt}{popink},left=14pt,right=14pt,top=11pt,bottom=11pt,before skip=11pt,after skip=15pt,
  fontupper=\popsemi\fontsize{10.6}{14.5}\selectfont\color{popink},
  fontlower=\popsemi\fontsize{10.6}{14.5}\selectfont\color{popink}}}
% Coche dessinée (le glyphe ✓ n'existe pas dans les polices du document)
\newcommand{\popcheck}[1]{\tikz[baseline=-0.5ex]\draw[#1,line width=1.6pt,line cap=round,line join=round](-0.13,0.0)--(-0.04,-0.09)--(0.13,0.1);}
\providecommand{\coche}{\popcheck{popgreen}}
\providecommand{\resultat}[1]{\par\smallskip\noindent\fcolorbox{popgreen}{popgreenL}{\parbox{\dimexpr\linewidth-2\fboxsep-2\fboxrule\relax}{\textbf{Résultat.} #1}}\par\smallskip}
\newcommand{\popboxhead}[3]{\poppill{#1}{#2}{white}\ifx\relax#3\relax\else\hspace{6pt}{\popxbold\fontsize{10.6}{12}\selectfont\color{popink}#3}\fi\par\vspace{7pt}}

\newtcolorbox{aretenir}[1][]{popbase,colframe=popgreen,before upper={\popboxhead{À retenir}{popgreen}{#1}}}
\newtcolorbox{exemplebox}[1][]{popbase,colframe=poporange,before upper={\popboxhead{Exemple}{poporange}{#1}}}
\newtcolorbox{definition}[1][]{popbase,colframe=popblue,before upper={\popboxhead{Définition}{popblue}{#1}}}
\newtcolorbox{propriete}[1][]{popbase,colframe=poppurple,before upper={\popboxhead{Propriété}{poppurple}{#1}}}
\newtcolorbox{methode}[1][]{popbase,colframe=popgold,before upper={\popboxhead{Méthode}{popgold}{#1}}}
\newtcolorbox{attention}[1][]{popbase,colframe=poppink,before upper={\popboxhead{Attention}{poppink}{#1}}}
\newtcolorbox{experience}[1][]{popbase,colframe=popblue,colback=popblueL,before upper={\popboxhead{Expérience}{popblue}{#1}}}
% « Le savais-tu ? » : carte violette pleine (contexte, culture, Cameroun)
\newtcolorbox{savaistu}[1][]{popbase,colback=poppurple,colframe=popink,
  fontupper=\popsemi\fontsize{10.4}{14.2}\selectfont\color{white},
  before upper={\poppill{Le savais-tu\,?}{white}{poppurple}\ifx\relax#1\relax\else\hspace{6pt}{\popxbold\color{white}#1}\fi\par\vspace{7pt}}}
% Exercice résolu : énoncé en haut, \tcblower puis la solution sur fond crème
\newcounter{popexo}\newcounter{popexr}
\newtcolorbox{exoresolu}[1][]{popbase,colframe=poporange,colbacklower=popcream,
  segmentation style={popink,line width=1.2pt,dashed},
  before upper={\stepcounter{popexr}\popboxhead{Exercice résolu \thepopexr}{poporange}{#1}},
  before lower={\poppill{Solution}{popink}{popcream}\par\vspace{6pt}}}
% Exercice (non résolu en place) — la correction est dans la partie Corrigés
\newtcolorbox{exercice}[1][]{popbase,colframe=popink,
  before upper={\stepcounter{popexo}\popboxhead{Exercice \thepopexo}{popink}{#1}}}
% Corrigé
\newtcolorbox{corrige}[1][]{popbase,colframe=popgreen,colback=popgreenL,
  before upper={\popboxhead{Corrigé}{popgreen}{#1}}}
% Objectifs / prérequis / bilan
\newtcolorbox{objectifs}{popbase,colframe=popink,colback=poporangeL,
  before upper={\popboxhead{Objectifs de la leçon}{popink}{}}}
\newtcolorbox{prerequis}{popbase,colframe=popink,colback=popblueL,
  before upper={\popboxhead{Prérequis}{popblue}{}}}
\newtcolorbox{bilan}{popbase,colframe=popink,colback=white,boxrule=2.8pt,
  before upper={\popboxhead{Fiche bilan}{poporange}{L'essentiel à connaître pour l'examen}}}
% Figure encadrée avec légende
\newtcolorbox{popfigure}[1][]{enhanced,breakable=false,colback=white,colframe=popline,boxrule=1.4pt,arc=8pt,
  left=10pt,right=10pt,top=10pt,bottom=8pt,before skip=10pt,after skip=12pt,halign=center,
  lower separated=false,halign lower=center,
  fontlower=\popsemi\fontsize{9}{11.5}\selectfont\color{popmuted},
  #1}
\newcommand{\legende}[1]{\par\vspace{6pt}{\popsemi\fontsize{9}{11.5}\selectfont\color{popmuted}#1}}

% ============================================================
% Signature OnBuch+
% ============================================================
\newcommand{\poplogoinline}[1]{{\poplogo\fontsize{#1}{#1}\selectfont\color{popink}OnBuch\color{poporange}+}}
\newcommand{\signatureonbuch}{%
  \begin{tcolorbox}[enhanced,colback=popink,colframe=popink,boxrule=2.2pt,arc=10pt,
      drop shadow=poporange,left=14pt,right=14pt,top=10pt,bottom=10pt,before skip=0pt,after skip=0pt]
    \tikz[baseline=-0.7ex]\node[circle,fill=popgreen,draw=popcream,line width=1.3pt,minimum size=22pt,inner sep=0]{\popcheck{white}};%
    \hspace{9pt}\begin{minipage}[c]{0.62\linewidth}
      {\popxbold\fontsize{12}{13}\selectfont\color{popcream}Signé\ \textbullet\ {\poplogo OnBuch\color{poporange}+}}\par\vspace{2pt}
      {\raggedright\popsemi\fontsize{8}{10.5}\selectfont\color{popline}\MakeUppercase{Équipe pédagogique OnBuch+}\\\MakeUppercase{Conforme au programme officiel MINESEC}\par}
    \end{minipage}\hfill
    {\poplogo\fontsize{22}{22}\selectfont\color{popcream}OnBuch\color{poporange}+}
  \end{tcolorbox}%
}

% ============================================================
% Page de garde
% #1 titre · #2 matière · #3 niveau · #4 module · #5 n° leçon · #6 durée ·
% #7 description / objectifs (texte ou itemize)
% ============================================================
\newcommand{\pagegarde}[7]{%
  \thispagestyle{empty}
  \newgeometry{a4paper,margin=1.7cm,top=1.6cm,bottom=1.4cm}%
  \begin{tikzpicture}[remember picture,overlay]
    \fill[poporange!18] ([xshift=-3.2cm,yshift=-3.6cm]current page.north east) circle (4.6cm);
    \fill[poppurple!14] ([xshift=2.4cm,yshift=7.6cm]current page.south west) circle (3.8cm);
  \end{tikzpicture}%
  {\poplogo\fontsize{30}{30}\selectfont\color{popink}OnBuch\color{poporange}+}\hfill
  \raisebox{0.6ex}{\poppill{Cours signé \textbullet\ Premium}{popink}{popcream}}\par
  \vspace{1pt}\popsquiggle{poppurple}\par
  \vspace{8pt}
  \begin{tcolorbox}[enhanced,colback=poporange,colframe=popink,boxrule=2.8pt,arc=13pt,
      drop shadow=popink,left=18pt,right=18pt,top=12pt,bottom=13pt,after skip=10pt]
    \poppill{#2 \textbullet\ #3}{popink}{popcream}\hspace{5pt}\poppill{#4}{white}{popink}\par\vspace{8pt}
    {\popdisplay\fontsize{14}{14}\selectfont\color{popink}LEÇON #5}\par\vspace{4pt}
    {\raggedright\hyphenpenalty=10000\exhyphenpenalty=10000\popdisplay\fontsize{25}{27}\selectfont\color{popcream}\MakeUppercase{#1}\par}
    \vspace{8pt}
    {\popxbold\fontsize{9.5}{9.5}\selectfont\color{popink}\MakeUppercase{Durée conseillée : #6}}
  \end{tcolorbox}
  \begin{tcolorbox}[enhanced,colback=white,colframe=popink,boxrule=2.2pt,arc=10pt,
      drop shadow=popink,left=16pt,right=16pt,top=10pt,bottom=10pt,after skip=8pt]
    \poppill{Dans cette leçon}{poppurple}{white}\par\vspace{5pt}
    \popsemi\fontsize{9.3}{12.6}\selectfont\color{popink}#7
  \end{tcolorbox}
  \noindent\begin{minipage}[t]{0.487\linewidth}
    \begin{tcolorbox}[enhanced,colback=popgreen,colframe=popink,boxrule=2.2pt,arc=10pt,
        drop shadow=popink,left=11pt,right=11pt,top=10pt,bottom=10pt,height=3.1cm,valign=center]
      \tcbox[colback=white,colframe=popink,boxrule=1.3pt,arc=5pt,boxsep=0pt,nobeforeafter,
        left=3pt,right=3pt,top=3pt,bottom=3pt,valign=center]{\qrcode[height=1.9cm]{https://play.google.com/store/apps/details?id=com.luvvix.onbuch&hl=fr}}%
      \hspace{8pt}\begin{minipage}[c]{0.5\linewidth}
        {\popdisplay\fontsize{11}{12.5}\selectfont\color{white}\MakeUppercase{Télécharge OnBuch}}\par\vspace{4pt}
        {\raggedright\popsemi\fontsize{8.4}{11}\selectfont\color{white}Gratuit \textbullet\ Google Play\\Tuteur IA, annales, concours, résultats\par}
      \end{minipage}
    \end{tcolorbox}
  \end{minipage}\hfill
  \begin{minipage}[t]{0.487\linewidth}
    \begin{tcolorbox}[enhanced,colback=poppurple,colframe=popink,boxrule=2.2pt,arc=10pt,
        drop shadow=popink,left=11pt,right=11pt,top=10pt,bottom=10pt,height=3.1cm,valign=center]
      \tikz[baseline=-0.7ex]\node[circle,fill=white,draw=popink,line width=1.3pt,minimum size=1.6cm,inner sep=0]{\popdisplay\fontsize{24}{24}\selectfont\color{poppurple}+};%
      \hspace{8pt}\begin{minipage}[c]{0.6\linewidth}
        {\popdisplay\fontsize{11}{12.5}\selectfont\color{white}\MakeUppercase{Passe à OnBuch+}}\par\vspace{4pt}
        {\raggedright\popsemi\fontsize{8.4}{11}\selectfont\color{white}Tous les cours signés, Tuteur IA illimité, fiches PDF — onglet OnBuch+ de l'app\par}
      \end{minipage}
    \end{tcolorbox}
  \end{minipage}
  \vspace{8pt}
  \popcontenu
  \vfill
  \signatureonbuch
  \restoregeometry
  \newpage
}

% Rangée « Ce cours contient » (page de garde)
\newcommand{\poptile}[3]{% #1 couleur #2 gros texte #3 légende
  \begin{tcolorbox}[enhanced,colback=#1,colframe=popink,boxrule=2pt,arc=9pt,shadow={2.5pt}{-2.5pt}{0pt}{popink},
      left=6pt,right=6pt,top=9pt,bottom=9pt,halign=center,valign=center,height=1.75cm,width=\linewidth,nobeforeafter]
    {\popdisplay\fontsize{12.5}{13}\selectfont\color{white}\MakeUppercase{#2}}\par\vspace{3pt}
    {\centering\popsemi\fontsize{7.8}{9.5}\selectfont\color{white}#3\par}
  \end{tcolorbox}}
\newcommand{\popcontenu}{%
  \par\noindent{\popxboldls\fontsize{7.5}{7.5}\selectfont\color{popmuted}CE COURS SIGNÉ CONTIENT}\par\vspace{7pt}
  \noindent\begin{minipage}[t]{0.235\linewidth}\poptile{popblue}{Cours}{complet, illustré, pas à pas}\end{minipage}\hfill
  \begin{minipage}[t]{0.235\linewidth}\poptile{poporange}{Méthodes}{et exercices résolus}\end{minipage}\hfill
  \begin{minipage}[t]{0.235\linewidth}\poptile{poppink}{Exercices}{type examen + corrigés}\end{minipage}\hfill
  \begin{minipage}[t]{0.235\linewidth}\poptile{popgreen}{Fiche bilan}{l'essentiel à retenir}\end{minipage}\par}

% Sommaire
\newtcolorbox{sommaire}{enhanced,colback=white,colframe=popink,boxrule=2.2pt,arc=10pt,
  drop shadow=popink,left=18pt,right=18pt,top=13pt,bottom=13pt,before skip=6pt,after skip=20pt,
  before upper={\poppill{Sommaire}{poppurple}{white}\par\vspace{9pt}\popsemi\fontsize{10.6}{14.5}\selectfont\color{popink}}}

% Signature de fin de cours — poussée tout en bas de la dernière page
\newcommand{\findecours}{%
  \par\vfill\nopagebreak
  \begin{center}\popsquiggle{poporange}\end{center}\vspace{4pt}
  \signatureonbuch
}
\AtBeginDocument{\def\llbracket{\mathopen{[\mkern-3.5mu[}}\def\rrbracket{\mathclose{]\mkern-3.5mu]}}} % intervalles d'entiers (glyphe ⟦ absent de la police)
% Glyphes absents de Fira Math : redéfinis à partir de glyphes disponibles
\AtBeginDocument{%
  \def\setminus{\mathbin{\backslash}}\let\smallsetminus\setminus
  \def\vdots{\mathord{\vbox{\baselineskip3.2pt\lineskiplimit0pt\kern4pt\hbox{.}\hbox{.}\hbox{.}}}}%
  \def\ddots{\mathinner{\mkern1mu\raise6.5pt\hbox{.}\mkern2mu\raise3.6pt\hbox{.}\mkern2mu\raise0.7pt\hbox{.}\mkern1mu}}%
  \def\triangle{\mathord{\scalebox{0.9}{$\symup{\Delta}$}}}%
  \def\nabla{\mathord{\raisebox{\depth}{\scalebox{1}[-1]{$\symup{\Delta}$}}}}%
  \def\checkmark{\popcheck{popdark}}%
  \def\star{\mathbin{\ast}}\def\bigstar{\mathord{\ast}}%
  \def\diamond{\mathbin{\scalebox{0.6}{$\Diamond$}}}%
  \def\bigcup{\mathop{\mathchoice{\raisebox{-0.3em}{\scalebox{1.7}{$\cup$}}}{\scalebox{1.25}{$\cup$}}{\cup}{\cup}}\displaylimits}%
  \def\bigcap{\mathop{\mathchoice{\raisebox{-0.3em}{\scalebox{1.7}{$\cap$}}}{\scalebox{1.25}{$\cap$}}{\cap}{\cap}}\displaylimits}%
  \def\lesseqqgtr{\mathrel{\lessgtr}}\def\bigoplus{\mathop{\oplus}\displaylimits}%
}
\providecommand{\ssi}{si et seulement si} % abréviation fréquente
\AtBeginDocument{\providecommand{\wideparen}{\overparen}} % arc de cercle

%% FIGFIX-BEGIN v5 — correctifs globaux des figures (docs/onbuch-generation/tools/figfix)
\usepackage{adjustbox}\usepackage{etoolbox}\tcbuselibrary{raster}
% toute figure TikZ/pgfplots plus large que la ligne est réduite à la largeur disponible
\BeforeBeginEnvironment{tikzpicture}{\begin{adjustbox}{max width=\linewidth}}
\AfterEndEnvironment{tikzpicture}{\end{adjustbox}}
% légendes pgfplots lisibles par-dessus les courbes
\pgfplotsset{every axis/.append style={legend style={fill=white,fill opacity=0.92,text opacity=1,draw=black!15,inner sep=3pt}}}
% étiquettes d'arêtes (syntaxe edge["..."]) avec fond blanc
\tikzset{every edge quotes/.append style={fill=white,inner sep=1.2pt}}
%% FIGFIX-END
\begin{document}

\pagegarde{L'alternateur}{Physique}{1ère TI}{Module 4 : Aspects énergétiques des circuits électriques}{N17}{4 h}{Cette leçon présente l'alternateur, machine électromécanique qui convertit l'énergie mécanique en énergie électrique sous forme de courant alternatif. De la dynamo de vélo aux centrales hydroélectriques du Cameroun, l'élève découvre son principe de fonctionnement, ses parties essentielles et la chaîne de transformations d'énergie qui aboutit à la production d'électricité.
\vspace{4pt}
\begin{multicols}{2}\small
\begin{itemize}
\item À la fin de cette leçon, je sais définir un alternateur et donner son rôle dans la production d'énergie électrique
\item À la fin de cette leçon, je sais expliquer le principe de fonctionnement d'un alternateur à partir de l'induction électromagnétique
\item À la fin de cette leçon, je sais identifier sur un schéma ou un dispositif réel le rotor, le stator et les autres parties essentielles d'un alternateur
\item À la fin de cette leçon, je sais donner le rôle de chaque partie essentielle d'un alternateur
\item À la fin de cette leçon, je sais faire le schéma de principe d'un alternateur
\item À la fin de cette leçon, je sais décrire la chaîne des transformations d'énergie lors de la production d'énergie électrique
\end{itemize}\end{multicols}}

\begin{sommaire}
\begin{multicols}{2}
\begin{enumerate}
\item Rôle de l'alternateur et mise en évidence expérimentale
\item Principe de fonctionnement
\item Parties essentielles de l'alternateur
\item Chaîne de transformations d'énergie
\item Exemples d'alternateurs et applications
\item Réalisation d'un alternateur simple et synthèse
\item Activité d'intégration
\item Méthodes et exercices résolus
\item Exercices
\item Corrigés des exercices
\item Fiche bilan
\end{enumerate}
\end{multicols}
\end{sommaire}

\begin{objectifs}
\begin{itemize}
\item À la fin de cette leçon, je sais définir un alternateur et donner son rôle dans la production d'énergie électrique
\item À la fin de cette leçon, je sais expliquer le principe de fonctionnement d'un alternateur à partir de l'induction électromagnétique
\item À la fin de cette leçon, je sais identifier sur un schéma ou un dispositif réel le rotor, le stator et les autres parties essentielles d'un alternateur
\item À la fin de cette leçon, je sais donner le rôle de chaque partie essentielle d'un alternateur
\item À la fin de cette leçon, je sais faire le schéma de principe d'un alternateur
\item À la fin de cette leçon, je sais décrire la chaîne des transformations d'énergie lors de la production d'énergie électrique
\item À la fin de cette leçon, je sais citer des exemples d'alternateurs utilisés dans la vie courante et dans les centrales électriques
\item À la fin de cette leçon, je sais réaliser un dispositif simple (bobine + aimant en rotation) produisant un courant alternatif
\end{itemize}
\end{objectifs}

\begin{prerequis}
\begin{itemize}
\item Le champ magnétique créé par un aimant ou un courant (notions de Seconde)
\item Le flux magnétique à travers une bobine et le phénomène d'induction électromagnétique (leçons précédentes du module)
\item La tension et le courant alternatifs sinusoïdaux : période, fréquence, valeur efficace (leçon sur le courant alternatif)
\item Les formes d'énergie et les conversions d'énergie (énergie mécanique, énergie électrique)
\item La puissance électrique P = U×I et le rendement d'une conversion
\end{itemize}
\end{prerequis}

\newpage

\coursec{Rôle de l'alternateur et mise en évidence expérimentale}

À Yaoundé comme à Ngaoundéré, quand ENEO coupe le courant, les quartiers s'éclairent grâce aux groupes électrogènes ; au barrage de Songloulou, d'immenses machines tournent jour et nuit pour alimenter tout le Sud-Cameroun. Pourtant, un simple vélo équipé d'une dynamo produit aussi de l'électricité pour éclairer sa lampe. Quel point commun entre une dynamo de vélo, un groupe électrogène de quartier et les turbines du barrage de la Sanaga ? Tous contiennent un \cle{alternateur}, machine capable de transformer un mouvement de rotation en courant électrique alternatif. Comment fonctionne-t-il et d'où vient l'énergie électrique qu'il fournit ? C'est à ces questions que répond cette leçon, en commençant par une mise en évidence expérimentale simple que tu pourras reproduire au laboratoire.

\courssub{De la dynamo de vélo à la centrale hydroélectrique : une même machine}

\textbf{Observation quotidienne.} Pédale sur un vélo équipé d'une « dynamo » : la lampe s'allume dès que la roue tourne, brille davantage quand tu pédales plus vite, et s'éteint à l'arrêt. Aucune pile, aucune batterie : le seul « carburant » est le mouvement. De même, un groupe électrogène ne produit du courant que lorsque son moteur à essence tourne ; au barrage de Songloulou, c'est la chute d'eau de la Sanaga qui fait tourner les turbines. Dans les trois cas, un \emph{mouvement de rotation} est transformé en \emph{énergie électrique}.

\textbf{Analyse du point commun.} Ces trois dispositifs, malgré leurs tailles très différentes, reposent sur le même organe : une machine électromécanique tournante, appelée \cle{alternateur}. Seule change la source d'énergie mécanique qui l'entraîne :

\begin{center}
\begin{tabularx}{\linewidth}{|l|X|X|}
\hline
\rowcolor{popblueL} \textbf{Dispositif} & \textbf{Source d'énergie mécanique} & \textbf{Ordre de grandeur de la puissance} \\
\hline
Dynamo de vélo & Muscles du cycliste & quelques watts (\SI{3}{W}) \\
\hline
Groupe électrogène & Moteur à essence ou gasoil & de \SI{1}{kW} à quelques \SI{100}{kW} \\
\hline
Alternateur de barrage & Turbine entraînée par l'eau & plusieurs dizaines de mégawatts \\
\hline
\end{tabularx}
\end{center}

\begin{definition}[L'alternateur]
Un \cle{alternateur} est une \textbf{machine électromécanique tournante} qui convertit l'\textbf{énergie mécanique de rotation} qu'on lui fournit en \textbf{énergie électrique}, délivrée sous forme d'un \textbf{courant alternatif} (tension et courant qui changent périodiquement de sens).
\end{definition}

\begin{aretenir}[Rôle de l'alternateur]
L'alternateur est le \textbf{générateur de courant alternatif} par excellence : c'est lui qui produit l'électricité dans \emph{toutes} les centrales électriques (hydrauliques comme Songloulou, Edéa ou Nachtigal, thermiques, éoliennes, nucléaires) et dans tous les groupes électrogènes. Sans alternateur, pas de réseau électrique national.
\end{aretenir}

\begin{attention}[L'aimant ne « fournit » pas d'électricité]
Erreur fréquente : croire que l'aimant de l'alternateur est une sorte de « réserve d'électricité » qui s'épuiserait. C'est faux. Un aimant immobile devant une bobine ne produit \emph{rien}. C'est le \textbf{mouvement relatif} entre l'aimant et la bobine qui est à l'origine de la tension : l'énergie électrique provient de l'\textbf{énergie mécanique} dépensée pour faire tourner la machine (muscles, moteur, chute d'eau). L'aimant permanent ne s'use pas et ne se « décharge » jamais.
\end{attention}

\courssub{Mise en évidence expérimentale de la production d'un courant alternatif}

\begin{experience}[Un aimant tournant devant une bobine]
\textbf{Matériel :} une bobine plate de plusieurs centaines de spires de fil de cuivre, un aimant droit (ou un aimant en U) monté sur un axe de rotation (manivelle ou petit moteur), une lampe à incandescence de faible puissance (\SI{1.5}{V}) ou une DEL avec résistance de protection, des fils de connexion, et pour l'observation fine, un oscilloscope.

\textbf{Protocole :}
\begin{enumerate}
\item Relier les deux bornes de la bobine à la lampe. Aimant au repos devant la bobine : observer.
\item Faire tourner lentement l'aimant devant la bobine, à vitesse à peu près constante : observer la lampe.
\item Augmenter la vitesse de rotation : observer de nouveau.
\item Remplacer la lampe par un oscilloscope branché aux bornes de la bobine et enregistrer la tension $u(t)$.
\end{enumerate}

\textbf{Observations :}
\begin{itemize}
\item Aimant immobile : la lampe reste éteinte, la tension est nulle.
\item Aimant en rotation : la lampe s'allume (elle scintille si la rotation est lente) ; l'oscilloscope affiche une tension qui prend alternativement des valeurs positives et négatives.
\item Plus l'aimant tourne vite, plus la lampe brille, plus l'amplitude de la tension est grande \emph{et} plus les oscillations sont rapprochées (fréquence plus élevée).
\end{itemize}

\textbf{Interprétation :} la rotation de l'aimant fait varier, à travers la bobine, le champ magnétique qui la traverse. Cette variation engendre dans la bobine une tension électrique : c'est le phénomène d'\cle{induction électromagnétique}. Comme le champ vu par la bobine change de sens à chaque demi-tour de l'aimant, la tension produite change de sens périodiquement : elle est \textbf{alternative}.
\end{experience}

\begin{popfigure}
\begin{center}
\begin{tikzpicture}[scale=1.0, >=Stealth]
% Bobine (vue en coupe, spires)
\foreach \y in {-0.9,-0.45,0,0.45,0.9}{
  \draw[popblue, thick] (3.2,\y) circle (0.16);
  \fill[popblue] (3.2,\y) circle (0.03);
}
\draw[popblue, thick] (2.9,-1.15) rectangle (3.5,1.15);
\node[popblue] at (3.2,1.55) {\small Bobine ($N$ spires)};
% Aimant en rotation
\fill[popink] (-1.1,-0.35) rectangle (0,-0.02);
\fill[popmuted] (-1.1,0.02) rectangle (0,0.35);
\draw[thick] (-1.1,-0.35) rectangle (0,0.35);
\node[white] at (-0.85,0.18) {\small N};
\node[white] at (-0.25,-0.18) {\small S};
% Axe de rotation
\draw[thick, popdark] (-0.55,0) -- (-0.55,-1.3);
\fill[popdark] (-0.55,-1.3) circle (0.07);
\node[popdark] at (-0.55,-1.65) {\small axe (manivelle)};
% Flèche de rotation
\draw[->, poporange, very thick] (0.35,0.75) arc (60:120:1.15);
\node[poporange] at (-0.55,1.35) {\small rotation $\omega$};
% Champ magnétique vers la bobine
\draw[->, popgreen, thick] (0.25,0.55) -- (2.6,0.55) node[midway, above, popgreen]{\small $\vec{B}$ variable};
\draw[->, popgreen, thick] (0.25,-0.55) -- (2.6,-0.55);
% Fils vers la lampe
\draw[thick] (3.5,0.9) -- (5.3,0.9) -- (5.3,0.35);
\draw[thick] (3.5,-0.9) -- (5.3,-0.9) -- (5.3,-0.35);
% Lampe
\draw[thick] (5.3,0) circle (0.35);
\draw[thick] (5.12,-0.18) -- (5.48,0.18);
\draw[thick] (5.12,0.18) -- (5.48,-0.18);
\node at (6.25,0) {\small lampe};
\node at (5.3,-0.75) {\small (ou oscilloscope)};
\end{tikzpicture}
\end{center}
\legende{Dispositif de mise en évidence : un aimant mis en rotation devant une bobine reliée à une lampe. La lampe s'allume alors que le circuit ne contient ni pile ni batterie.}
\end{popfigure}

\begin{methode}[Mettre en évidence la production d'un courant alternatif]
\begin{enumerate}
\item Réaliser un circuit fermé contenant \textbf{uniquement} une bobine et un récepteur témoin (lampe, DEL, oscilloscope) : aucune source d'énergie électrique dans le circuit.
\item Approcher un aimant et le mettre en \textbf{mouvement relatif} par rapport à la bobine (rotation de l'aimant, ou déplacement de la bobine, ou rotation de la bobine devant l'aimant fixe).
\item Observer le récepteur : si la lampe brille ou si l'oscilloscope affiche une tension, de l'énergie électrique est produite.
\item Vérifier le caractère \textbf{alternatif} : la tension change de signe périodiquement (courbe sinusoïdale à l'oscilloscope, ou DEL qui ne s'allume qu'une alternance sur deux).
\item Conclure : c'est le \textbf{mouvement} (énergie mécanique) qui a été converti en énergie électrique.
\end{enumerate}
\end{methode}

\courssub{Analyse des observations : tension produite et fréquence}

\textbf{Forme de la tension.} Lorsque l'aimant tourne à vitesse angulaire constante $\omega$ devant la bobine, le champ magnétique vu par la bobine varie de façon périodique et régulière. La tension recueillie aux bornes de la bobine est alors pratiquement \textbf{sinusoïdale} :
\[ u(t) = U_{\max}\,\sin(\omega t + \varphi) \]
où $U_{\max}$ est l'\textbf{amplitude} (valeur maximale, en volts) et $\omega$ la \textbf{pulsation} (en $\mathrm{rad\cdot s^{-1}}$), liée à la fréquence $f$ par $\omega = 2\pi f$. Vérifions l'homogénéité : $\omega t$ est en $\mathrm{rad\cdot s^{-1}}\times\mathrm{s}$, c'est-à-dire sans dimension (un angle), donc l'argument du sinus est cohérent et $u$ s'exprime bien en volts comme $U_{\max}$.

\begin{popfigure}
\begin{center}
\begin{tikzpicture}
\begin{axis}[
  width=0.85\linewidth, height=5.2cm,
  axis lines=middle,
  xlabel={$t$ (ms)}, ylabel={$u$ (V)},
  xlabel style={at={(ticklabel* cs:1)}, anchor=west},
  ylabel style={at={(ticklabel* cs:1)}, anchor=south},
  xmin=0, xmax=42, ymin=-3.2, ymax=3.2,
  xtick={0,10,20,30,40},
  ytick={-2.5,0,2.5},
  yticklabels={$-U_{\max}$,$0$,$U_{\max}$},
  grid=both, grid style={popline!40},
]
\addplot[popcurve, domain=0:40, samples=300] {2.5*sin(360*0.025*x)};
\draw[dashed, popmuted] (axis cs:0,2.5) -- (axis cs:40,2.5);
\draw[<->, poporange, thick] (axis cs:10,0) -- (axis cs:20,0) node[midway, below, poporange]{\small $T$};
\end{axis}
\end{tikzpicture}
\end{center}
\legende{Tension observée à l'oscilloscope aux bornes de la bobine : tension alternative sinusoïdale de période $T$, avec $f = 1/T$. Ici $T = \SI{10}{ms}$, soit $f = \SI{100}{Hz}$.}
\end{popfigure}

\textbf{Influence de la vitesse de rotation.} L'expérience montre deux faits essentiels, qu'il faut retenir :
\begin{itemize}
\item \textbf{Amplitude :} plus l'aimant tourne vite, plus le champ magnétique varie rapidement dans la bobine, et plus la tension produite est grande : $U_{\max}$ \emph{augmente} avec la vitesse de rotation.
\item \textbf{Fréquence :} chaque tour de l'aimant correspond à un cycle complet de la tension. Si l'aimant effectue $n$ tours par seconde, la tension effectue $n$ périodes par seconde : la fréquence $f$ de la tension est \emph{proportionnelle} à la vitesse de rotation.
\end{itemize}

\begin{exemplebox}[Fréquence d'une dynamo de vélo]
La « dynamo » d'un vélo est un petit alternateur dont l'aimant tourne à $n = 20$ tours par seconde quand le vélo roule à allure modérée. La tension produite a pour fréquence :
\[ f = n = \SI{20}{Hz} \qquad \text{soit une période } T = \dfrac{1}{f} = \dfrac{1}{20} = \SI{0.05}{s} = \SI{50}{ms}. \]
Si le cycliste double sa vitesse, $n$ passe à $40$ tours par seconde : la fréquence devient $f' = \SI{40}{Hz}$ \emph{et} la lampe brille davantage car l'amplitude de la tension augmente. C'est exactement ce que tu observes quand tu pédales plus vite.
\end{exemplebox}

\courssub{Conclusion : le mouvement relatif, cause de la production d'électricité}

\begin{propriete}[Condition de production d'une tension par induction]
Une tension apparaît aux bornes d'une bobine \textbf{uniquement si le champ magnétique qui la traverse varie} au cours du temps. Cette variation est obtenue par le \textbf{mouvement relatif} entre l'aimant (ou l'électroaimant) et la bobine. Ce phénomène porte le nom d'\cle{induction électromagnétique} ; il constitue le principe de fonctionnement de tous les alternateurs.
\end{propriete}

Retenons la logique énergétique de l'expérience : le circuit ne contient aucune pile, pourtant la lampe s'allume. D'où vient l'énergie ? De l'opérateur qui tourne la manivelle : en tournant, il doit vaincre des actions électromagnétiques qui résistent au mouvement (tu le sens concrètement : la manivelle devient plus « dure » à tourner quand la lampe est branchée !). Le travail mécanique fourni est converti en énergie électrique. C'est le sens précis de la définition de l'alternateur.

\begin{exoresolu}[Exploiter les observations de l'expérience]
Un élève fait tourner un aimant devant une bobine de $500$ spires reliée à un oscilloscope. À la vitesse de $5$ tours par seconde, il mesure une tension d'amplitude $U_{\max} = \SI{1.5}{V}$. (a) Quelle est la fréquence de la tension obtenue ? (b) Que deviennent qualitativement l'amplitude et la fréquence s'il tourne deux fois plus vite ? (c) L'élève affirme : « c'est l'aimant qui fournit l'électricité à la bobine ». Corriger cette affirmation.
\tcblower
\textbf{Solution.} (a) À chaque tour de l'aimant correspond une période de la tension, donc $f = n = \SI{5}{Hz}$ (période $T = 1/f = \SI{0.20}{s}$). (b) Si la vitesse double ($n' = 10$ tours/s), la fréquence double : $f' = \SI{10}{Hz}$, et l'amplitude \emph{augmente} car le champ magnétique varie plus rapidement dans la bobine ($U'_{\max} > \SI{1.5}{V}$). (c) L'affirmation est fausse : l'aimant immobile ne produit aucune tension. C'est le \textbf{mouvement} de l'aimant, c'est-à-dire l'\textbf{énergie mécanique} fournie par l'élève, qui est convertie en énergie électrique. L'aimant ne se consomme pas ; seule l'énergie mécanique est dépensée.
\end{exoresolu}

\begin{savaistu}[La « dynamo » de vélo est en réalité un alternateur !]
Le langage courant appelle « dynamo » le petit appareil qui éclaire les vélos. Or, en physique, une \emph{dynamo} est une machine qui produit un courant \textbf{continu}, grâce à un collecteur qui redresse la tension. Le dispositif de vélo, lui, produit un courant \textbf{alternatif} : c'est donc un véritable petit \textbf{alternateur}, composé d'un aimant tournant devant une bobine fixe. La confusion vient de l'histoire : les premières machines à courant continu ont été appelées « dynamos », et le mot est resté dans le langage populaire. Toi qui étudies la physique, tu sais maintenant faire la différence !
\end{savaistu}

\coursec{Principe de fonctionnement : la loi de Faraday}

\courssub{Flux magnétique et force électromotrice induite}

L'expérience qualitative montre qu'un mouvement relatif aimant/bobine crée une tension. Pour le quantifier, on introduit le \textbf{flux magnétique} $\Phi$ traversant une spire de surface $S$ et de normale $\vec{n}$ :
\[ \Phi = \vec{B} \cdot \vec{S} = B S \cos\theta \]
où $\theta$ est l'angle entre $\vec{B}$ et la normale à la spire. Pour une bobine de $N$ spires identiques, le flux total (ou flux enchainé) est $\Phi_{\text{tot}} = N \Phi$.

\begin{propriete}[Loi de Faraday (ou loi de l'induction)]
La force électromotrice (f.e.m.) $e$ induite aux bornes d'une bobine est égale à l'opposé de la dérivée temporelle du flux magnétique total qui la traverse :
\[ e = -\frac{d\Phi_{\text{tot}}}{dt} = -N \frac{d\Phi}{dt} \]
L'unité SI est le volt ($\mathrm{V} = \mathrm{Wb\cdot s^{-1}}$). Le signe moins traduit la \textbf{loi de Lenz} : le courant induit (s'il circule) crée un champ magnétique qui s'oppose à la variation du flux qui l'a produit.
\end{propriete}

\courssub{Application à l'alternateur : tension sinusoïdale}

Considérons une bobine de $N$ spires, de surface $S$, tournant à vitesse angulaire constante $\omega$ dans un champ magnétique uniforme $\vec{B}$. L'angle entre $\vec{B}$ et la normale à la bobine vaut $\theta = \omega t$ (choix de l'origine des temps). Le flux est :
\[ \Phi(t) = B S \cos(\omega t) \]
La f.e.m. induite est donc :
\[ e(t) = -N \frac{d}{dt}\left[ B S \cos(\omega t) \right] = N B S \omega \sin(\omega t) \]
On retrouve une tension \textbf{sinusoïdale} d'amplitude $U_{\max} = N B S \omega$ et de pulsation $\omega$.

\begin{aretenir}[Liens fondamentaux pour un alternateur]
\begin{itemize}
\item \textbf{Fréquence :} $f = \dfrac{\omega}{2\pi} = n$ (nombre de tours par seconde pour un alternateur bipolaire).
\item \textbf{Amplitude :} $U_{\max} = N B S \omega = 2\pi N B S f$. Elle est proportionnelle à la fréquence (donc à la vitesse de rotation), au nombre de spires, au champ magnétique et à la surface de la bobine.
\end{itemize}
\end{aretenir}

\begin{attention}[Sens de la f.e.m. et tension de sortie]
La f.e.m. $e$ est la tension que l'alternateur \emph{produit} intérieurement. La tension $u$ aux bornes de l'alternateur (tension de charge) vaut $u = e - r i$ (régime continu) ou $\underline{U} = \underline{E} - \underline{I} \underline{Z}_{\text{int}}$ (régime sinusoïdal), où $r$ (ou $Z_{\text{int}}$) est la résistance (impédance) interne. Au Probatoire, on confond souvent $u$ et $e$ pour un alternateur \emph{idéal} (à vide, $i=0$).
\end{attention}

\coursec{Parties essentielles de l'alternateur}

\courssub{Architecture générale : stator et rotor}

Un alternateur industriel (centrale, groupe électrogène) est une machine \textbf{synchrone}. Il comporte deux parties principales :

\begin{center}
\begin{tabularx}{\linewidth}{|l|X|X|}
\hline
\rowcolor{popblueL} \textbf{Partie} & \textbf{Description} & \textbf{Rôle} \\
\hline
\textbf{Stator} (partie fixe) & Carcasse en acier magnétique portant un enroulement \textbf{induit} (généralement 3 phases décalées de $120^\circ$) logé dans des rainures. & C'est dans l'induit que la tension alternative est \textbf{produite} et récupérée. Il supporte les forts courants de charge. \\
\hline
\textbf{Rotor} (partie tournante) & \textbf{Inducteur} : électroaimant alimenté en courant \textbf{continu} (excitation) via des anneaux collecteurs et balais, ou aimants permanents (petites puissances). & Crée le champ magnétique $\vec{B}$ tournant. C'est la source du flux magnétique variable dans l'induit. \\
\hline
\end{tabularx}
\end{center}

\begin{popfigure}
\begin{center}
\begin{tikzpicture}[scale=0.9, >=Stealth]
% Stator (carcasse externe)
\draw[popdark, thick, fill=popblue!5] (0,0) circle (3.5);
\draw[popdark, thick] (0,0) circle (2.8);
% Enroulements stator (3 phases symboliques)
\foreach \angle/\color in {30/popblue, 150/popgreen, 270/poporange}{
  \draw[\color, thick, fill=\color!20] (\angle:2.9) rectangle ++(0.6,0.6);
  \draw[\color, thick, fill=\color!20] (\angle+180:2.9) rectangle ++(0.6,0.6);
}
% Rotor (pôles saillants bipolaire)
\draw[popink, thick, fill=popink!20] (-1.2,-0.5) rectangle (1.2,0.5);
\node[popink] at (0,0) {\textbf{N}};
\draw[popmuted, thick, fill=popmuted!20] (-1.2,-0.5) rectangle (1.2,0.5);
\node[popmuted] at (0,-1.2) {\textbf{S}}; % Just for visual separation, actually N/S on same rotor body
% Correction: Rotor is one body with N and S poles.
\draw[popink, thick, fill=popink!20] (-1.2,0.5) -- (0,1.2) -- (1.2,0.5) -- (1.2,-0.5) -- (0,-1.2) -- (-1.2,-0.5) -- cycle;
\node[white] at (0,0.3) {\large N};
\node[white] at (0,-0.3) {\large S};
% Axe
\draw[popdark, thick] (0,0) -- (0,-2.5);
\draw[popdark, thick] (0,0) -- (0,2.5);
% Flèches rotation
\draw[->, poporange, very thick] (2.5,2.5) arc (45:135:3.5);
\node[poporange] at (0,3.5) {$\omega$};
% Légendes
\node[popblue] at (0,4) {\textbf{Stator (Induit - 3 phases)}};
\node[popink] at (0,-3.5) {\textbf{Rotor (Inducteur - Excitation DC)}};
% Balais / Anneaux (schématique)
\draw[popdark, thick] (2.5,0) -- (3.5,0);
\draw[popdark, thick, fill=popgold] (3.5,0.15) rectangle (3.8,0.35);
\draw[popdark, thick, fill=popgold] (3.5,-0.35) rectangle (3.8,-0.15);
\node at (4.5,0) {\small Balais + Anneaux (Excitation)};
\draw[->, poppurple] (3.65,0.25) -- (3.65,0.8) node[above, poppurple] {$I_{\text{exc}}$};
\end{tikzpicture}
\end{center}
\legende{Schéma de principe d'un alternateur synchrone bipolaire : le rotor (inducteur) tourne à l'intérieur du stator (induit). L'excitation du rotor se fait par courant continu via anneaux et balais.}
\end{popfigure}

\courssub{Le circuit d'excitation}

Pour que le rotor produise un champ magnétique intense et réglable, on utilise un \textbf{électroaimant} alimenté en courant continu : c'est le \textbf{circuit d'excitation}.
\begin{itemize}
\item \textbf{Source d'excitation :} petite machine auxiliaire (excitatrice) montée sur le même arbre, ou redresseur statique (thyristors) alimenté par le réseau.
\item \textbf{Anneaux collecteurs (bagues) et balais :} assurent le contact électrique tournant entre le circuit fixe (excitatrice) et le circuit tournant (bobine du rotor). Contrairement au collecteur d'une dynamo (courant continu), ces anneaux sont \textbf{lis} : ils ne redressent pas le courant, ils le transmettent tel quel (continu).
\end{itemize}

\begin{propriete}[Régulation de la tension]
En faisant varier le courant d'excitation $I_{\text{exc}}$, on fait varier le champ magnétique $B$ du rotor, donc l'amplitude $U_{\max}$ de la tension de sortie ($U_{\max} \propto B \propto I_{\text{exc}}$). C'est le principe du \textbf{régulateur de tension} (AVR - Automatic Voltage Regulator) qui maintient la tension du réseau stable (ex: \SI{230}{V} ou \SI{400}{V}) malgré les variations de charge.
\end{propriete}

\courssub{Schéma de principe à retenir pour le Probatoire}

\begin{methode}[Schéma de principe d'un alternateur synchrone]
\begin{enumerate}
\item Dessiner le \textbf{stator} : un cercle extérieur avec 3 enroulements (ou un seul pour le monophasé) notés $S_1, S_2, S_3$ (ou $A, B, C$).
\item Dessiner le \textbf{rotor} : un cercle intérieur avec un enroulement unique (bobine d'excitation) noté $R$ (ou $E_+$, $E_-$), muni de deux anneaux collecteurs et balais.
\item Indiquer le \textbf{sens de rotation} $\omega$ (flèche circulaire).
\item Indiquer le \textbf{courant d'excitation} $I_{\text{exc}}$ entrant/sortant par les balais.
\item Indiquer les \textbf{tensions de sortie} $u_1, u_2, u_3$ (déphasées de $120^\circ$) aux bornes de l'induit.
\end{enumerate}
\end{methode}

\coursec{Chaîne de transformations d'énergie}

\courssub{Bilan énergétique global}

L'alternateur est un \textbf{convertisseur d'énergie}. Il ne \emph{crée} pas d'énergie, il la transforme.

\begin{center}
\begin{tikzpicture}[node distance=1.5cm, >=Stealth, thick]
\node[draw, rounded corners, fill=popgreen!20, minimum width=3cm, align=center] (source){Source d'énergie primaire\\(Chute d'eau, Vapeur, Vent,\\Moteur thermique, Muscles)};
\node[draw, rounded corners, fill=poporange!20, minimum width=3cm, right=of source, align=center] (meca){Énergie mécanique\\(Arbre de rotation\\Couple $C$, $\omega$)};
\node[draw, rounded corners, fill=popblue!20, minimum width=3cm, right=of meca, align=center] (elec){Énergie électrique\\(Tension $u$, Courant $i$,\\Fréquence $f$)};
\node[draw, rounded corners, fill=poppink!20, minimum width=2.5cm, below=of meca, align=center] (pertes){Pertes\\(Joule $R_i I^2$, Frottements,\\Hystérésis, Foucault)};
\draw[->, popgreen] (source) -- node[above] {Transformation} (meca);
\draw[->, poporange] (meca) -- node[above, align=center]{Conversion\\électromécanique} (elec);
\draw[->, popmuted] (meca) -- node[left] {Pertes mécaniques} (pertes);
\draw[->, popmuted] (elec) -- node[right] {Pertes électriques} (pertes);
\end{tikzpicture}
\end{center}

\begin{definition}[Rendement de l'alternateur]
Le rendement $\eta$ (eta) est le rapport de l'énergie électrique utile fournie à l'énergie mécanique reçue :
\[ \eta = \frac{P_{\text{élec utile}}}{P_{\text{méca reçue}}} = \frac{P_{\text{élec utile}}}{P_{\text{élec utile}} + P_{\text{pertes}}} \]
Pour les gros alternateurs de centrale (Songloulou, Nachtigal), $\eta$ dépasse \SI{98}{\%}. Pour un petit groupe électrogène, $\eta \approx \SI{85}{\%}$ à \SI{90}{\%}.
\end{definition}

\courssub{Action électromagnétique de freinage (Lenz)}

Quand l'alternateur délivre un courant $i$ à une charge, les conducteurs de l'induit (stator) parcourus par ce courant subissent des forces de Laplace $\vec{F} = i \vec{l} \wedge \vec{B}$. Ces forces créent un \textbf{couple électromagnétique de freinage} $C_{\text{em}}$ qui s'oppose à la rotation (loi de Lenz).

\begin{propriete}[Équilibre mécanique de l'alternateur en charge]
Le moteur premier (turbine, moteur diesel) doit fournir un couple moteur $C_{\text{mot}}$ tel que :
\[ C_{\text{mot}} = C_{\text{em}} + C_{\text{pertes méca}} \]
Plus la charge électrique est forte (grand $i$), plus $C_{\text{em}}$ est grand, plus le moteur doit fournir d'effort, plus il consomme de carburant (ou d'eau). C'est la traduction mécanique de la conservation de l'énergie.
\end{propriete}

\begin{exemplebox}[Groupe électrogène de quartier]
Un groupe de \SI{10}{kVA} tourne à \SI{1500}{tr/min} ($f=\SI{50}{Hz}$ pour un 4 pôles). À vide, le moteur diesel tourne facilement. Quand on branche des lampes, un frigo, une TV (charge totale \SI{6}{kW}), le moteur « tousse », son régime tend à baisser, le régulateur de vitesse (gouvernor) ouvre l'arrivée de gazole pour maintenir \SI{1500}{tr/min}. La consommation de gazole augmente proportionnellement à la puissance électrique fournie (plus pertes).
\end{exemplebox}

\coursec{Exemples d'alternateurs et applications au Cameroun}

\courssub{Typologie selon la source d'énergie mécanique}

\begin{center}
\begin{tabularx}{\linewidth}{|l|X|X|}
\hline
\rowcolor{popblueL} \textbf{Type d'alternateur} & \textbf{Source mécanique (Moteur premier)} & \textbf{Exemples au Cameroun} \\
\hline
\textbf{Hydro-alternateur} & Turbine hydraulique (Pelton, Francis, Kaplan) & Barrages de \textbf{Songloulou} (\SI{388}{MW}), \textbf{Edéa} (\SI{263}{MW}), \textbf{Nachtigal} (\SI{420}{MW}), \textbf{Lom Pangar} (\SI{30}{MW}), \textbf{Memve'ele} (\SI{211}{MW}). \\
\hline
\textbf{Turbo-alternateur} & Turbine à vapeur (centrales thermiques) ou turbine à gaz & Centrale thermique de \textbf{Kribi} (gaz, \SI{216}{MW}), \textbf{Dibamba} (fioul lourd, \SI{88}{MW}), \textbf{Limbe}. \\
\hline
\textbf{Alternateur éolien} & Éolienne (hélice) & Parc éolien projeté au Nord (Moyen-Chari), petits sites isolés. \\
\hline
\textbf{Alternateur de groupe électrogène} & Moteur diesel / essence & Quartiers non raccordés, hôpitaux (groupe de secours), chantiers, télécoms (sites MTN/Orange/CAMTEL). \\
\hline
\textbf{Alternateur de véhicule} & Moteur thermique (via courroie) & Tous les véhicules (voitures, motos, camions) : recharge batterie + alimentation éclairage, injection, clim. \\
\hline
\end{tabularx}
\end{center}

\courssub{Particularités constructives}

\begin{itemize}
\item \textbf{Nombre de pôles :} Pour obtenir $f = \SI{50}{Hz}$ avec une vitesse de turbine faible (ex: \SI{150}{tr/min} = \SI{2.5}{tr/s} pour une turbine Kaplan), on augmente le nombre de pôles $p$ : $f = p \times n$. Un alternateur de Songloulou a $p = 20$ pôles ($f = 20 \times 2.5 = \SI{50}{Hz}$).
\item \textbf{Tension de sortie :} \SI{10}{kV} à \SI{15}{kV} en centrale (pour limiter les courants et pertes lignes), puis transformée en \SI{225}{kV} / \SI{400}{kV} pour le transport (SONATREL), puis \SI{30}{kV} / \SI{15}{kV} distribution, \SI{400}{V} / \SI{230}{V} utilisation.
\item \textbf{Refroidissement :} Air (petits), Hydrogène ou Eau (gros alternateurs centraux) pour évacuer les pertes Joule massives.
\end{itemize}

\begin{savaistu}[Le barrage de Nachtigal : un exemple majeur]
Mis en service progressivement depuis 2023, le barrage de Nachtigal sur la Sanaga (près de Batchenga) abrite \textbf{7 alternateurs} de \SI{60}{MW} chacun (total \SI{420}{MW}). Ce sont des alternateurs synchrones à rotor à pôles saillants, entraînés par des turbines Francis à vitesse constante (\SI{166.7}{tr/min} pour $f=\SI{50}{Hz}$ avec $p=18$ pôles). Leur tension de sortie est de \SI{11}{kV}. Ils fournissent environ \SI{30}{\%} de l'électricité du réseau interconnecté Sud (RIS) géré par ENEO/SONATREL.
\end{savaistu}

\coursec{Réalisation d'un alternateur simple et synthèse}

\courssub{Activité : Construire un alternateur de laboratoire}

\begin{experience}[Alternateur à aimant tournant (ou bobine tournante) - TP réalisable au lycée]
\textbf{Objectif :} Produire une tension alternative mesurable à l'oscilloscope et vérifier $f \propto n$ et $U_{\max} \propto n$.

\textbf{Matériel :}
\begin{itemize}
\item Aimants néodyme puissants (disques \SI{20}{mm} x \SI{5}{mm}, 2 à 4 pièces).
\item Fil émaillé (\SI{0.3}{mm} à \SI{0.5}{mm}), bobine de \SI{500}{à 1000} spires sur formeur cylindrique (diam \SI{3}{cm}).
\item Moteur DC jouet (\SI{3}{à 12}{V}) avec manivelle ou alimentation variable pour vitesse contrôlée.
\item Support imprimé 3D ou bois pour fixer aimant sur axe moteur et bobine fixe face à aiment (entrefer \SI{1}{à 2}{mm}).
\item Oscilloscope (ou carte son + logiciel Audacity/Scilab en entrée micro), multimètre RMS.
\end{itemize}

\textbf{Protocole :}
\begin{enumerate}
\item Monter l'aimant sur l'axe du moteur (rotor). Fixer la bobine (stator) face à l'aimant.
\item Brancher l'oscilloscope aux bornes de la bobine (couplage AC).
\item Faire tourner le moteur à vitesse faible, moyenne, rapide (varier tension moteur).
\item Mesurer pour chaque vitesse : $f$ (curseurs oscilloscope), $U_{\max}$ (amplitude crête).
\item Tracer $f(n)$ et $U_{\max}(n)$ : vérifier proportionnalité.
\end{enumerate}

\textbf{Résultats attendus :}
\begin{itemize}
\item Courbe $u(t)$ sinusoïdale (plus ou moins déformée si aimant non parfait).
\item $f = k \times n$ (droite passant par l'origine, $k =$ nombre de paires de pôles = 1 ici).
\item $U_{\max} = k' \times n$ (droite passant par l'origine).
\end{itemize}
\end{experience}

\courssub{Synthèse des compétences (Check-list Probatoire)}

\begin{aretenir}[Ce que tu dois savoir faire pour le Probatoire]
\begin{enumerate}
\item \textbf{Définir} l'alternateur : machine tournante, conversion énergie mécanique $\to$ électrique, sortie alternative.
\item \textbf{Citer} le rôle : production de l'électricité dans toutes les centrales et groupes.
\item \textbf{Identifier et nommer} sur un schéma : Stator (induit), Rotor (inducteur), Anneaux/Balais (excitation), Axe, Couple moteur, Tensions de sortie.
\item \textbf{Expliquer le principe} : Induction électromagnétique (Faraday), variation flux $\Phi = B S \cos(\omega t) \to e = N B S \omega \sin(\omega t)$.
\item \textbf{Établir les relations} : $f = p \cdot n$ ($p$ paires de pôles), $U_{\max} \propto N B S f$.
\item \textbf{Dresser le bilan énergétique} : Source $\to$ Mécanique $\to$ Électrique + Pertes. Expliquer l'augmentation de consommation carburant avec la charge (couple freinant Lenz).
\item \textbf{Donner des exemples} camerounais : Songloulou, Nachtigal, Edéa (hydro), Kribi, Dibamba (thermique), groupes quartiers, alternateur auto.
\item \textbf{Réaliser / Schématiser} un dispositif simple produisant du courant alternatif (aimant + bobine + rotation).
\end{enumerate}
\end{aretenir}

\begin{exercice}[Entraînement Probatoire - Alternateur de centrale]
Un alternateur bipolaire ($p=1$) de centrale hydroélectrique est entraîné par une turbine Kaplan tournant à $n = \SI{150}{tr/min}$.
\begin{enumerate}
\item Calculer la fréquence $f$ de la tension produite. Est-elle compatible avec le réseau ENEO ?
\item L'alternateur a un induit (stator) triphasé à $N = 200$ spires par phase, surface par spire $S = \SI{0.5}{m^2}$, champ magnétique dans l'entrefer $B = \SI{0.8}{T}$. Calculer l'amplitude $U_{\max}$ de la tension simple (phase-neutre) à vide.
\item On raccorde cet alternateur au réseau \SI{400}{V} (tension composée efficace) via un transformateur élévateur. Quelle doit être la tension efficace phase-neutre à la sortie de l'alternateur ? En déduire le rapport de transformation $K$ du transformateur (secondaire/primaire) si le secondaire est en \SI{225}{kV}.
\end{enumerate}
\end{exercice}

\begin{corrige}[Exercice 2]
\begin{enumerate}
\item $n = \SI{150}{tr/min} = \frac{150}{60} = \SI{2.5}{tr/s} = \SI{2.5}{Hz}$. Pour un bipolaire ($p=1$), $f = p \cdot n = \SI{2.5}{Hz}$. \textbf{Non compatible} avec le réseau \SI{50}{Hz}. Il faut un alternateur multipolaire ($p = 20$ pôles = 10 paires) pour avoir $f = 10 \times 2.5 = \SI{50}{Hz}$.
\item $\omega = 2\pi f = 2\pi \times 50 = \SI{100\pi}{rad/s} \approx \SI{314}{rad/s}$.
$U_{\max} = N B S \omega = 200 \times 0.8 \times 0.5 \times 314 = \SI{25120}{V} \approx \SI{25.1}{kV}$ (crête).
Tension efficace phase-neutre : $U = \frac{U_{\max}}{\sqrt{2}} \approx \SI{17.8}{kV}$.
\item Réseau \SI{400}{V} composé efficace $\Rightarrow$ Tension phase-neutre efficace $U_{\text{réseau}} = \frac{400}{\sqrt{3}} \approx \SI{231}{V}$.
L'alternateur doit délivrer cette tension (après transformation). Mais ici l'alternateur sort du \SI{17.8}{kV}. Le transformateur élévateur est donc \textbf{abaisseur} pour raccorder au réseau HT ? Non, l'alternateur sort \SI{10-15}{kV}, on élève en \SI{225}{kV}.
Reprenons : L'exercice dit "raccordé au réseau \SI{400}{V} via transformateur élévateur". C'est contradictoire (élévateur monte la tension). Supposons : Alternateur $\to$ Transformateur Élévateur $\to$ Réseau HT \SI{225}{kV}.
Tension primaire (alternateur) $U_1 = \SI{17.8}{kV}$ (efficace). Tension secondaire $U_2 = \SI{225}{kV}$.
Rapport $K = \frac{U_2}{U_1} = \frac{225}{17.8} \approx 12.6$.
\end{enumerate}
\end{corrige}

\begin{exercice}[QCM - Pièges classiques]
Choisir la bonne réponse.
\begin{enumerate}
\item Dans un alternateur synchrone, le courant d'excitation circule dans : (a) l'induit (stator) (b) l'inducteur (rotor) (c) les deux.
\item La fréquence de la tension de sortie dépend : (a) seulement de la vitesse de rotation (b) seulement du nombre de pôles (c) de la vitesse de rotation ET du nombre de pôles.
\item Quand on augmente la charge (puissance active) sur un groupe électrogène tournant à vitesse constante : (a) la tension augmente (b) la tension baisse si pas de régulation AVR (c) la fréquence augmente.
\item L'énergie électrique fournie par l'alternateur provient : (a) de l'aimant (b) du carburant (énergie mécanique) (c) du stator.
\end{enumerate}
\end{exercice}

\begin{corrige}[Exercice 3 - QCM]
\begin{enumerate}
\item \textbf{(b)} L'excitation (courant continu) alimente l'inducteur (rotor) pour créer le champ $B$. L'induit (stator) délivre le courant alternatif de puissance.
\item \textbf{(c)} $f = p \cdot n$. Les deux paramètres entrent en jeu.
\item \textbf{(b)} La charge crée un courant $I$ dans l'induit $\to$ réaction d'induit (champ déaimantant) $\to$ chute de tension interne $I Z_s$. Sans régulateur de tension (AVR) agissant sur $I_{\text{exc}}$, la tension aux bornes baisse.
\item \textbf{(b)} Conservation de l'énergie. L'aimant/électroaimant est seulement le médiateur du champ.
\end{enumerate}
\end{corrige}

\coursec{Principe de fonctionnement}

Dans la section précédente, l'expérience de la bobine reliée à un galvanomètre nous a livré un fait capital : lorsqu'un aimant tourne devant une bobine fixe, un courant apparaît dans le circuit, et ce courant \emph{change de sens} périodiquement. L'alternateur produit donc une tension \cle{alternative}. Reste à comprendre \emph{pourquoi}. C'est l'objet de cette section : nous allons partir du phénomène d'induction électromagnétique, suivre pas à pas la rotation de l'aimant, et établir que la tension produite est sinusoïdale, de fréquence égale à la fréquence de rotation. C'est ce principe qui, du barrage de Nachtigal au groupe électrogène du quartier, alimente tout le Cameroun en courant alternatif de fréquence \SI{50}{Hz}.

\courssub{Rappel : l'induction électromagnétique}

\begin{definition}[Flux magnétique à travers une bobine]
Le \cle{flux magnétique} $\Phi$ à travers une bobine de $N$ spires, de surface $S$ chacune, placée dans un champ magnétique uniforme $\vec{B}$ faisant un angle $\theta$ avec la normale aux spires, est :
\[ \Phi = N\,B\,S\cos\theta \]
Le flux s'exprime en webers (Wb), avec $B$ en teslas (T) et $S$ en mètres carrés (\si{m^2}).
\end{definition}

L'observation fondamentale, due à Faraday (1831), est la suivante : ce n'est pas le champ magnétique lui-même qui crée une tension dans la bobine, mais la \textbf{variation} du flux. Une bobine plongée dans un champ intense mais constant ne produit rien ; en revanche, dès que le flux à travers elle varie — parce que l'aimant s'approche, s'éloigne ou tourne — une \cle{force électromotrice induite} (f.é.m.), notée $e$, apparaît à ses bornes.

\begin{aretenir}[L'idée directrice]
Toute \textbf{variation de flux magnétique} à travers un circuit y induit une f.é.m. Si le circuit est fermé, cette f.é.m. y fait circuler un \textbf{courant induit}. Dans l'alternateur, c'est la \textbf{rotation} de l'aimant (ou de l'électroaimant) qui crée cette variation de flux.
\end{aretenir}

\courssub{Rotation de l'aimant devant la bobine : analyse pas à pas}

Plaçons-nous dans la configuration la plus simple : un aimant droit (pôle Nord, pôle Sud) monté sur un axe, tournant à vitesse constante devant une bobine fixe dont le circuit est fermé sur une lampe. Suivons le flux $\Phi$ à travers la bobine au cours d'un tour complet, en repérant quatre positions remarquables.

\begin{itemize}
\item \textbf{Position 1 ($\theta = 0$)} : le pôle Nord de l'aimant fait face à la bobine. Le champ $\vec{B}$ traverse les spires perpendiculairement : le flux est \textbf{maximal}. Mais à cet instant précis, il cesse de croître pour commencer à décroître : sa \emph{variation} est nulle, donc la f.é.m. induite est nulle.
\item \textbf{Position 2 ($\theta = 90°$)} : l'aimant est horizontal, parallèle au plan des spires. Le flux est \textbf{nul}, mais il varie alors le plus vite (il passe du maximum à zéro) : la f.é.m. induite est \textbf{maximale}, et le courant circule dans un sens.
\item \textbf{Position 3 ($\theta = 180°$)} : c'est le pôle Sud qui fait face à la bobine. Le flux est à nouveau maximal en valeur absolue, mais de \textbf{sens opposé} ; sa variation est nulle : la f.é.m. repasse par zéro.
\item \textbf{Position 4 ($\theta = 270°$)} : le flux est nul et varie le plus vite en sens inverse : la f.é.m. est maximale de \textbf{signe opposé}, le courant circule dans l'autre sens.
\end{itemize}

Puis l'aimant revient à la position 1 : le cycle recommence. À chaque tour de l'aimant correspond donc \textbf{une alternance complète} de la tension : c'est le cœur du principe de l'alternateur.

\begin{popfigure}
\begin{center}
\begin{tikzpicture}[scale=0.92, every node/.style={font=\small}]
% ---------- Position 1 : theta = 0 ----------
\begin{scope}[shift={(0,0)}]
\draw[popblue, thick] (0,-0.9) ellipse (0.28 and 0.9);
\draw[popblue, thick] (0.35,-0.9) ellipse (0.28 and 0.9);
\node[popblue, below] at (0.2,-1.95) {bobine fixe};
\fill[poporange] (-2.3,-0.55) rectangle (-1.3,0.55);
\fill[popgreen] (-3.3,-0.55) rectangle (-2.3,0.55);
\node[white] at (-1.8,0) {\textbf{N}};
\node[white] at (-2.8,0) {\textbf{S}};
\draw[-{Stealth}, popdark, thick] (-2.3,0.75) arc (150:30:0.6);
\node[popdark, above] at (-2.3,0.95) {$\vec{\omega}$};
\node[align=center] at (-1.5,-2.6) {\textbf{1.} $\theta=0$ : $\Phi$ maximal\\ $e=0$};
\end{scope}
% ---------- Position 2 : theta = 90 ----------
\begin{scope}[shift={(6.2,0)}]
\draw[popblue, thick] (0,-0.9) ellipse (0.28 and 0.9);
\draw[popblue, thick] (0.35,-0.9) ellipse (0.28 and 0.9);
\fill[poporange] (-2.55,-1.35) rectangle (-1.55,-0.35);
\fill[popgreen] (-2.55,0.35) rectangle (-1.55,1.35);
\node[white] at (-2.05,-0.85) {\textbf{N}};
\node[white] at (-2.05,0.85) {\textbf{S}};
\draw[-{Stealth}, popdark, thick] (-2.05,1.55) arc (150:30:0.6);
\node[align=center] at (-1.5,-2.6) {\textbf{2.} $\theta=90°$ : $\Phi=0$\\ $e$ maximale, courant $\nearrow$};
\draw[-{Stealth}, poppink, very thick] (0.6,0) -- (1.5,0) node[right]{$i$};
\end{scope}
% ---------- Position 3 : theta = 180 ----------
\begin{scope}[shift={(0,-4.6)}]
\draw[popblue, thick] (0,-0.9) ellipse (0.28 and 0.9);
\draw[popblue, thick] (0.35,-0.9) ellipse (0.28 and 0.9);
\fill[popgreen] (-2.3,-0.55) rectangle (-1.3,0.55);
\fill[poporange] (-3.3,-0.55) rectangle (-2.3,0.55);
\node[white] at (-1.8,0) {\textbf{S}};
\node[white] at (-2.8,0) {\textbf{N}};
\draw[-{Stealth}, popdark, thick] (-2.3,0.75) arc (150:30:0.6);
\node[align=center] at (-1.5,-2.6) {\textbf{3.} $\theta=180°$ : $\Phi$ minimal\\ $e=0$};
\end{scope}
% ---------- Position 4 : theta = 270 ----------
\begin{scope}[shift={(6.2,-4.6)}]
\draw[popblue, thick] (0,-0.9) ellipse (0.28 and 0.9);
\draw[popblue, thick] (0.35,-0.9) ellipse (0.28 and 0.9);
\fill[popgreen] (-2.55,-1.35) rectangle (-1.55,-0.35);
\fill[poporange] (-2.55,0.35) rectangle (-1.55,1.35);
\node[white] at (-2.05,-0.85) {\textbf{S}};
\node[white] at (-2.05,0.85) {\textbf{N}};
\draw[-{Stealth}, popdark, thick] (-2.05,1.55) arc (150:30:0.6);
\node[align=center] at (-1.5,-2.6) {\textbf{4.} $\theta=270°$ : $\Phi=0$\\ $e$ minimale, courant $\searrow$};
\draw[-{Stealth}, poppink, very thick] (1.5,0) -- (0.6,0) node[left]{$i$};
\end{scope}
\end{tikzpicture}
\end{center}
\legende{Quatre positions successives de l'aimant tournant devant la bobine fixe. Le flux $\Phi$ est maximal quand l'aimant fait face à la bobine (positions 1 et 3), nul quand il lui est parallèle (positions 2 et 4). La f.é.m. induite, liée à la \emph{vitesse de variation} du flux, est nulle aux positions 1 et 3, extrémale aux positions 2 et 4, avec un courant qui change de sens à chaque demi-tour.}
\end{popfigure}

\begin{propriete}[F.é.m. induite par rotation uniforme — admise]
La rotation \textbf{uniforme} d'un aimant (ou d'un électroaimant) devant une bobine fixe y induit une f.é.m. \cle{sinusoïdale} :
\[ e(t) = E_{\max}\,\sin(\omega t + \varphi) \]
dont la fréquence est \textbf{égale à la fréquence de rotation} de l'aimant. Ce résultat est admis en Première ; il se démontre en Terminale à partir de la loi de Faraday (voir le complément ci-dessous).
\end{propriete}

\courssub{De quoi dépend la f.é.m. induite ?}

Reprenons l'expérience de la section 1 en faisant varier un paramètre à la fois (démarche expérimentale rigoureuse !). Les observations sont sans ambiguïté :

\begin{itemize}
\item on remplace la bobine par une bobine de \textbf{plus grand nombre de spires} $N$ : l'éclat de la lampe augmente ;
\item on utilise un aimant \textbf{plus puissant} (champ $B$ plus intense) : l'éclat augmente ;
\item on fait tourner l'aimant \textbf{plus vite} : l'éclat augmente (et le clignotement s'accélère).
\end{itemize}

\begin{propriete}[Facteurs influençant la f.é.m. induite]
La valeur maximale $E_{\max}$ de la f.é.m. induite dans la bobine est d'autant plus grande que :
\begin{itemize}
\item le \textbf{nombre de spires} $N$ de la bobine est grand ;
\item l'\textbf{intensité du champ magnétique} $B$ de l'inducteur est grande ;
\item la \textbf{vitesse de rotation} $\omega$ de l'aimant est grande.
\end{itemize}
Pour une bobine de $N$ spires de surface $S$ dans le champ $B$ d'un aimant tournant à la pulsation $\omega$, on montre en Terminale que $E_{\max} = N\,B\,S\,\omega$.
\end{propriete}

\begin{attention}[Deux effets de la vitesse à ne pas confondre]
Augmenter la vitesse de rotation a \textbf{deux conséquences distinctes} : la f.é.m. maximale augmente ($E_{\max} = NBS\omega$) \emph{et} la fréquence de la tension augmente ($f = n$). Dans un alternateur de centrale, on ne peut donc pas « tourner plus vite pour avoir plus de tension » sans changer la fréquence du réseau : c'est pourquoi ENEO maintient la vitesse de rotation de ses alternateurs rigoureusement constante.
\end{attention}

\courssub{Fréquence de la tension produite : la relation $f = n$}

L'analyse des quatre positions ci-dessus nous a montré qu'\textbf{un tour complet de l'aimant} produit exactement \textbf{une période complète} de la tension (une alternance positive puis une alternance négative). Cette correspondance un tour $\leftrightarrow$ une période est la clé de la relation fondamentale de l'alternateur.

\begin{propriete}[Fréquence de la tension d'un alternateur à une paire de pôles]
Pour un alternateur dont le rotor comporte \textbf{une paire de pôles} (un Nord, un Sud), la fréquence $f$ de la tension produite est égale à la fréquence de rotation $n$ du rotor :
\[ \boxed{\;f = n\;} \]
avec $f$ en hertz (Hz) et $n$ en \textbf{tours par seconde} (tr/s).
\end{propriete}

\begin{experience}[Démonstration guidée : pourquoi $f = n$ ?]
Raisonnons sur les durées, en quatre étapes.
\begin{enumerate}
\item \textbf{Un tour = une période.} L'aimant n'a qu'une paire de pôles : au cours d'un tour, le pôle Nord puis le pôle Sud passent une fois devant la bobine. Le flux effectue donc exactement un cycle complet (maximal, nul, minimal, nul, maximal). La f.é.m. induite effectue donc exactement \textbf{une période} $T$ pendant que le rotor effectue \textbf{un tour}.
\item \textbf{Durée d'un tour.} Si le rotor tourne à $n$ tours par seconde, la durée d'un tour est $\tau = \dfrac{1}{n}$ (en secondes).
\item \textbf{Égalité des durées.} D'après l'étape 1, la période de la tension est $T = \tau = \dfrac{1}{n}$.
\item \textbf{Conclusion.} Par définition de la fréquence, $f = \dfrac{1}{T} = \dfrac{1}{1/n} = n$. \hfill$\square$
\end{enumerate}
\textit{Remarque} : si le rotor comporte $p$ paires de pôles, chaque tour produit $p$ périodes, et la relation devient $f = p\,n$. Les alternateurs des barrages hydroélectriques, qui tournent lentement, possèdent de nombreuses paires de pôles précisément pour atteindre \SI{50}{Hz} malgré une rotation lente.
\end{experience}

\begin{exemplebox}[La fréquence du réseau ENEO]
Le réseau électrique camerounais distribue une tension alternative de fréquence $f = \SI{50}{Hz}$. Pour un alternateur à une paire de pôles, cela impose une rotation à $n = f = \SI{50}{tr/s}$, soit $50 \times 60 = \SI{3000}{tr/min}$. C'est exactement la vitesse de rotation des groupes turbo-alternateurs des centrales thermiques (comme celles de Kribi ou de Limbé) : la turbine à vapeur ou à gaz entraîne le rotor à \SI{3000}{tr/min}, et le réseau « bat » à \SI{50}{Hz}.
\end{exemplebox}

\begin{exoresolu}[Calcul de la fréquence produite par un alternateur]
\textbf{Énoncé.} Le rotor d'un alternateur de centrale thermique, comportant une seule paire de pôles, tourne à la vitesse de $n = \SI{3000}{tr/min}$. 1) Convertir cette fréquence de rotation en tours par seconde. 2) En déduire la fréquence $f$ de la tension produite. 3) Calculer la période $T$ de cette tension.
\tcblower
\textbf{Solution détaillée.}
\begin{enumerate}
\item \textbf{Conversion.} Une minute contient 60 secondes, donc :
\[ n = \frac{3000 \text{ tours}}{60 \text{ s}} = \SI{50}{tr/s} \]
\item \textbf{Fréquence de la tension.} Le rotor possède une paire de pôles, donc :
\[ f = n = \SI{50}{Hz} \]
On retrouve la fréquence du réseau ENEO : c'est rassurant, c'est bien ainsi que le courant de nos prises est produit !
\item \textbf{Période.} Par définition :
\[ T = \frac{1}{f} = \frac{1}{50} = \SI{0,020}{s} = \SI{20}{ms} \]
La tension aux bornes d'une prise camerounaise effectue donc 50 oscillations complètes chaque seconde, chacune durant \SI{20}{ms}.
\end{enumerate}
\end{exoresolu}

\begin{attention}[Erreur fréquente au Probatoire : tr/min contre Hz]
L'erreur la plus classique consiste à écrire directement $f = 3000$ Hz parce que l'énoncé donne $n = \SI{3000}{tr/min}$. La relation $f = n$ n'est valable que si $n$ est exprimée en \textbf{tours par seconde}. Réflexe systématique : \textbf{diviser par 60} toute vitesse donnée en tr/min avant d'appliquer $f = n$. Vérifie aussi la cohérence du résultat : les fréquences industrielles sont de l'ordre de \SI{50}{Hz}, jamais de plusieurs milliers de hertz !
\end{attention}

\courssub{Complément Probatoire : la loi de Faraday et la forme sinusoïdale}

Ce complément, \textbf{hors programme strict de Première mais valorisé}, prépare le cours de Terminale et justifie rigoureusement ce que nous avons admis.

\begin{propriete}[Loi de Faraday — admise]
La f.é.m. induite dans un circuit est égale à l'opposé de la dérivée par rapport au temps du flux magnétique à travers ce circuit :
\[ e = -\frac{d\Phi}{dt} \]
Le signe « $-$ » traduit la loi de Lenz : les effets du courant induit s'opposent toujours à la variation de flux qui leur donne naissance.
\end{propriete}

Appliquons cette loi à notre alternateur. Le rotor tourne uniformément à la pulsation $\omega$ : l'angle entre le champ $\vec{B}$ et la normale aux spires est $\theta(t) = \omega t$. Le flux à travers la bobine de $N$ spires de surface $S$ s'écrit donc :
\[ \Phi(t) = NBS\cos(\omega t) = \Phi_{\max}\cos(\omega t) \]
La loi de Faraday donne alors :
\begin{align*}
e(t) &= -\frac{d\Phi}{dt} \\
&= -\frac{d}{dt}\left[\Phi_{\max}\cos(\omega t)\right] \\
&= \Phi_{\max}\,\omega\,\sin(\omega t) = NBS\omega\sin(\omega t)
\end{align*}
La f.é.m. est bien \textbf{sinusoïdale}, de pulsation $\omega$ égale à celle de la rotation, donc de fréquence $f = \dfrac{\omega}{2\pi} = n$. On retrouve au passage $E_{\max} = NBS\omega$, qui confirme l'influence de $N$, $B$ et de la vitesse de rotation. Remarque essentielle : $e$ est proportionnelle à la \emph{dérivée} de $\Phi$ ; or la dérivée d'un cosinus est un sinus : \textbf{flux et f.é.m. sont déphasés d'un quart de période}. Quand le flux est extrémal, la f.é.m. est nulle ; quand le flux est nul, la f.é.m. est extrémale — exactement ce que l'analyse des quatre positions nous avait fait deviner.

\begin{popfigure}
\begin{center}
\begin{tikzpicture}
\begin{axis}[
width=0.95\linewidth, height=6.2cm,
xlabel={$t$ (ms)}, ylabel={},
xmin=0, xmax=40, ymin=-1.5, ymax=1.5,
xtick={0,10,20,30,40},
ytick={-1,0,1},
grid=both, grid style={popline!40},
axis lines=left,
legend style={at={(0.98,0.98)}, anchor=north east, font=\small}
]
\addplot[popblue, very thick, domain=0:40, samples=200] {cos(360*x/20)};
\addlegendentry{$\Phi(t)$ (flux)}
\addplot[poporange, very thick, domain=0:40, samples=200] {1.3*sin(360*x/20)};
\addlegendentry{$e(t)$ (f.é.m.)}
\addplot[popmuted, dashed, domain=0:40, samples=2] {0};
\draw[popmuted, dashed] (axis cs:5,-1.5) -- (axis cs:5,1.5);
\draw[popmuted, dashed] (axis cs:10,-1.5) -- (axis cs:10,1.5);
\node[popmuted, font=\footnotesize, anchor=south] at (axis cs:7.5,1.35) {$T/4$};
\end{axis}
\end{tikzpicture}
\end{center}
\legende{Flux $\Phi(t)$ (bleu) et f.é.m. $e(t)$ (orange) pour un alternateur tournant à $n = \SI{50}{tr/s}$ ($T = \SI{20}{ms}$). Les deux grandeurs sont sinusoïdales de même fréquence, mais \textbf{déphasées d'un quart de période} : $e$ est nulle quand $\Phi$ est extrémal (positions 1 et 3 de l'aimant) et extrémale quand $\Phi$ s'annule (positions 2 et 4).}
\end{popfigure}

\begin{methode}[Exploiter le principe de l'alternateur dans un exercice]
\begin{enumerate}
\item \textbf{Identifier} le nombre de paires de pôles du rotor (souvent une seule en Première).
\item \textbf{Convertir} la vitesse de rotation en tours par seconde si elle est donnée en tr/min ($n = \dfrac{n_{\text{tr/min}}}{60}$).
\item \textbf{Appliquer} $f = n$ (une paire de pôles) ou $f = p\,n$ ($p$ paires de pôles).
\item \textbf{En déduire} la période $T = \dfrac{1}{f}$ si elle est demandée.
\item \textbf{Conclure} en précisant que la tension est sinusoïdale et en citant les paramètres ($N$, $B$, $\omega$) qui fixent son amplitude.
\end{enumerate}
\end{methode}

\begin{savaistu}[Pourquoi tout le réseau camerounais « bat-il » à la même fréquence ?]
Tous les alternateurs interconnectés au réseau SONATREL — ceux d'Édéa, de Song-Loulou, de Nachtigal, les centrales thermiques — doivent tourner \emph{en synchronisme} : leurs rotors, bien que distants de centaines de kilomètres, produisent des tensions qui oscillent exactement ensemble à \SI{50}{Hz}. Si la demande d'électricité augmente brusquement (le soir, quand toute la ville s'éclaire), les rotors ont tendance à ralentir : la fréquence chute légèrement, et les régulateurs des centrales injectent aussitôt plus de puissance mécanique pour la ramener à \SI{50}{Hz}. Les délestages que connaît parfois le réseau sont précisément des situations où la puissance disponible ne suffit plus à maintenir ce synchronisme.
\end{savaistu}

\begin{aretenir}[L'essentiel de la section]
\begin{itemize}
\item La rotation d'un aimant devant une bobine fait \textbf{varier périodiquement le flux} magnétique à travers elle : une \textbf{f.é.m. induite alternative} apparaît.
\item Pour une rotation uniforme, cette f.é.m. est \textbf{sinusoïdale}, de fréquence égale à la fréquence de rotation : $f = n$ (une paire de pôles), avec $n$ en tr/s.
\item L'amplitude de la f.é.m. augmente avec le \textbf{nombre de spires} $N$, le \textbf{champ} $B$ et la \textbf{vitesse de rotation} $\omega$ : $E_{\max} = NBS\omega$.
\item Flux et f.é.m. sont \textbf{déphasés d'un quart de période} (loi de Faraday $e = -\dfrac{d\Phi}{dt}$, admise).
\item Piège à éviter : toujours \textbf{convertir les tr/min en tr/s} (diviser par 60) avant d'appliquer $f = n$.
\end{itemize}
\end{aretenir}

\coursec{Parties essentielles de l'alternateur}

Dans la section précédente, nous avons mis en évidence le principe fondamental : la rotation d'un aimant devant une bobine (ou l'inverse) crée une force électromotrice alternative par induction électromagnétique. Pour transformer ce principe en une machine industrielle fiable, capable de délivrer de l'énergie électrique à grande échelle (barrage de \textbf{Lom Pangar}, centrale thermique de \textbf{Kribi}) ou à l'échelle d'un véhicule (alternateur de taxi, dynamo de vélo), il faut une architecture mécanique et électromagnétique précise. C'est l'objet de cette section : disséquer l'alternateur pour identifier chaque organe, comprendre son rôle et savoir le représenter sur un schéma de principe normalisé.

\courssub{Le rotor : la partie tournante, source du champ magnétique}

Le \cle{rotor} (ou \emph{inducteur}) est l'organe mobile de la machine. Son rôle exclusif est de créer un champ magnétique \textbf{tournant} dans l'entrefer (l'espace d'air entre rotor et stator). C'est la variation de ce flux magnétique à travers les bobines fixes qui induira la f.é.m.

\begin{definition}[Rotor (Inducteur)]
Le \textbf{rotor} est la partie tournante de l'alternateur. Il porte le système aimanté (aimants permanents ou électroaimant alimenté en courant continu) qui génère le champ magnétique $\vec{B}$ variable dans le référentiel du stator. Il est solidaire de l'arbre de rotation entraîné par la source d'énergie mécanique (turbine hydraulique, turbine à gaz, moteur diesel, pédalier de vélo).
\end{definition}

Il existe deux technologies majeures pour le rotor, qui définissent le type d'alternateur :

\begin{itemize}
    \item \textbf{Rotor à aimants permanents :} Des aimants (souvent en néodyme-fer-bore ou ferrite) sont fixés sur le pourtour du rotor. Le champ magnétique est constant dans le référentiel du rotor. Structure simple, robuste, sans entretien (pas de balais), mais tension de sortie non réglable et champ magnétique non modulable. C'est le cas de la \textbf{dynamo de vélo} et des petits groupes électrogènes portatifs.
    \item \textbf{Rotor à électroaimant (bobinage d'excitation) :} Le rotor porte un bobinage (enroulement d'excitation) traversé par un courant continu $I_e$ (courant d'excitation). Ce courant crée le champ magnétique $\vec{B} \propto I_e$. L'avantage majeur : en faisant varier $I_e$ (via un régulateur de tension), on module l'amplitude de la f.é.m. induite et on stabilise la tension de sortie quelle que soit la charge. C'est la technologie de l'\textbf{alternateur de voiture} et des \textbf{alternateurs de centrales électriques (ENEO/SONATREL)}.
\end{itemize}

\begin{popfigure}
\begin{tikzpicture}[scale=0.9, every node/.style={font=\small}]
    % Rotor à aimants permanents (gauche)
    \begin{scope}[xshift=-4cm]
        \draw[thick, fill=popgray!20] (0,0) circle (1.5cm); % Rotor body
        \foreach \angle in {0,90,180,270} {
            \draw[fill=popdark] (\angle:1.0) rectangle ++(0.8,0.6);
            \node at (\angle:1.3) {\tiny N/S};
        }
        \draw[->, thick, poporange] (0,-2) arc (-90:-270:2);
        \node[poporange] at (0,-2.5) {Rotation $\omega$};
        \node[align=center] at (0,2.5) {\textbf{Rotor à aimants}\\\textbf{permanents}};
        \node[align=center] at (0,0) {Noyau ferromagnétique\\+ Aimants collés};
    \end{scope}
    
    % Rotor à électroaimant (droite)
    \begin{scope}[xshift=4cm]
        \draw[thick, fill=popgray!20] (0,0) circle (1.5cm); % Rotor body
        % Poles
        \foreach \angle/\col in {0/popblue, 180/popred} {
            \draw[fill=\col!30, thick] (\angle:0.3) -- ++(\angle:1.0) arc (\angle:\angle+60:1.0) -- ++(\angle+60:-1.0) arc (\angle+60:\angle:0.3);
            \node at (\angle+30:1.2) {\textbf{\col Pôle}};
        }
        % Excitation winding representation
        \draw[poppurple, thick, decorate, decoration={snake, amplitude=1pt, segment length=4pt}] (0,0) circle (0.6cm);
        \node[poppurple] at (0,0) {$I_e$};
        \draw[->, thick, poporange] (0,-2) arc (-90:-270:2);
        \node[poporange] at (0,-2.5) {Rotation $\omega$};
        \node[align=center] at (0,2.5) {\textbf{Rotor à électroaimant}\\\textbf{(bobinage d'excitation)}};
        \node[align=center] at (0,0) {Noyau à pôles saillants\\+ Bobinage $I_e$};
    \end{scope}
\end{tikzpicture}
\legende{Deux technologies de rotor : à aimants permanents (simplicité, vélo) et à électroaimant (réglage tension, voiture, centrales). Le courant d'excitation $I_e$ est continu.}
\end{popfigure}

\courssub{Le stator : la partie fixe, siège de l'induction}

Le \cle{stator} (ou \emph{induit}) est l'organe fixe. Il reçoit le flux magnétique tournant créé par le rotor. C'est dans ses bobines que naît la force électromotrice alternative $e(t) = E_{\max} \sin(\omega t + \varphi)$ qui sera délivrée au réseau ou à la charge.

\begin{definition}[Stator (Induit)]
Le \textbf{stator} est la partie fixe de l'alternateur. Il est constitué d'un noyau magnétique (paquet de tôles isolées pour limiter les courants de Foucault) traversé par des enroulements (bobines) dans lesquels le flux magnétique tournant $\Phi(t)$ induit une f.é.m. alternative. C'est lui qui est raccordé au circuit extérieur (réseau ENEO, batterie de voiture via redresseur, phare de vélo).
\end{definition}

\begin{propriete}[Architecture du stator triphasé (Centrales, Industrie)]
Dans les alternateurs de puissance (barrages, centrales thermiques), le stator porte \textbf{trois enroulements identiques} décalés de $120^\circ$ électriques (décalage spatial de $2\pi/3$ rad). Ils produisent trois f.é.m. déphasées de $120^\circ$ : c'est la production \textbf{triphasée}, base du transport d'énergie électrique sur le réseau \textbf{SONATREL}.
\end{propriete}

\begin{attention}[Erreur fréquente : Inversion Rotor/Stator]
\textbf{Piège de l'examen :} « Le stator tourne et le rotor est fixe. » \textbf{FAUX.}
\emph{Mnémotechnique :} \textbf{R}otor = \textbf{R}otation (il \textbf{R}oule). \textbf{S}tator = \textbf{S}tationnaire (il \textbf{S}tagne).
Dans l'immense majorité des alternateurs (voiture, centrales, vélo), \textbf{l'aimant tourne (rotor) et la bobine est fixe (stator)}. Cela évite de devoir récupérer le courant de puissance (fort ampérage) sur des pièces mobiles (balais), ce qui serait source d'usure, d'étincelles et de pertes.
\end{attention}

\courssub{L'arbre de rotation et les paliers : la liaison mécanique}

L'arbre (ou \cle{arbre de rotation}) est l'axe cylindrique en acier trempé qui supporte le rotor et transmet le couple mécanique $C_m$ (en $\si{N.m}$) depuis la machine motrice (turbine Pelton à \textbf{Lom Pangar}, moteur thermique de groupe électrogène, pédalier) vers le rotor.

\begin{itemize}
    \item \textbf{Paliers (roulements) :} Ils assurent le guidage radial et axial de l'arbre en rotation (vitesse $\omega$ jusqu'à $3000\ \si{tr/min}$ pour $50\ \si{Hz}$ en bipolaire, $1500\ \si{tr/min}$ en quadripolaire). Ils supportent le poids du rotor, les efforts radiaux (courroie, engrenage) et l'effort axial (poussée magnétique). Lubrification par graisse ou huile sous pression (grands alternateurs).
    \item \textbf{Accouplement :} Liaison élastique ou rigide entre l'arbre de la machine motrice et l'arbre de l'alternateur. Alignement précis indispensable pour éviter vibrations et fatigue.
\end{itemize}

\courssub{Les balais et les bagues (collecteur) : l'alimentation de l'excitation}

\textbf{Uniquement pour les alternateurs à rotor bobiné (électroaimant).} Le rotor tournant a besoin d'un courant continu $I_e$ pour s'aimanter. On ne peut pas souder des fils sur une pièce tournante.

\begin{definition}[Balais et bagues collectrices]
Les \textbf{bagues} (ou anneaux collecteurs) sont des anneaux en cuivre isolés entre eux et de l'arbre, solidaires du rotor. Les \textbf{balais} (charbon/graphite, parfois cuivre-graphite) sont des pièces fixes, pressées par des ressorts sur les bagues. Ils assurent le contact électrique glissant pour injecter le courant continu d'excitation $I_e$ (faible intensité, quelques ampères) vers le bobinage du rotor.
\end{definition}

\begin{aretenir}[Pourquoi des bagues et non un collecteur à segments (comme en machine continue) ?]
Le courant d'excitation est \textbf{continu} dans le rotor. Il n'y a pas de commutation (inversion de courant) à réaliser. Des anneaux lisses (bagues) suffisent. C'est plus simple, plus fiable, moins d'étincelles, moins d'usure que le collecteur à segments d'une dynamo/machine continue. \textbf{Attention :} Les balais de l'alternateur de voiture ne récupèrent \textbf{PAS} le courant de charge (courant fort) ; ils ne servent qu'à l'excitation (courant faible). Le courant de charge sort par les bornes du stator.
\end{aretenir}

\courssub{La carcasse (bâti) : protection et circuit magnétique de retour}

La \cle{carcasse} (ou bâti, châssis) est l'enveloppe métallique (fonte, acier soudé, aluminium) qui remplit trois fonctions :
\begin{enumerate}
    \item \textbf{Support mécanique :} Elle maintient le stator (serrage des tôles), les paliers d'arbre, les porte-balais, le carter d'étanchéité.
    \item \textbf{Protection :} Indice de protection IP (ex. IP23, IP54) contre les corps solides, l'eau, la poussière (crucial pour les alternateurs de centrale en bord de rivière ou de groupe électrogène de chantier au Cameroun).
    \item \textbf{Ferme le circuit magnétique :} Elle assure le retour du flux magnétique entre les pôles du rotor (fermeture du circuit magnétique stator-entrefer-rotor-entrefer-stator). Elle doit être en matériau ferromagnétique (faible réluctance).
\end{enumerate}

\courssub{Distinction : Alternateur à aimant permanent vs Alternateur à électroaimant}

\begin{center}
\begin{tabularx}{\linewidth}{|l|X|X|}
\hline
\rowcolor{popblueL}
\textbf{Caractéristique} & \textbf{Alternateur à aimants permanents} & \textbf{Alternateur à électroaimant (bobiné)} \\
\hline
\textbf{Exemples} & Dynamo de vélo, petit groupe portable, éolienne domestique, alternateur de moto & Alternateur de voiture (camion, taxi), Alternateurs de centrales (Lom Pangar, Nachtigal, Kribi), Groupes électrogènes industriels \\
\hline
\textbf{Rotor} & Aimants collés/insérés (NdFeB, Ferrite) & Noyau à pôles + Bobinage d'excitation $I_e$ \\
\hline
\textbf{Balais / Bagues} & \textbf{Aucun} (si sortie directe stator) & \textbf{Oui} (alimentation $I_e$ rotor) \\
\hline
\textbf{Régulation tension} & Impossible (dépend seulement $\omega$, charge) & \textbf{Oui} (variation $I_e$ via régulateur AVR) \\
\hline
\textbf{Rendement} & Élevé (pas de pertes excitation) & Légerement inférieur (pertes Joule rotor $R_e I_e^2$) \\
\hline
\textbf{Coût / Maintenance} & Faible / Nulle & Plus élevé / Maintenance balais (rare) \\
\hline
\end{tabularx}
\end{center}

\courssub{Schéma de principe normalisé de l'alternateur}

À l'examen (Probatoire), on te demande souvent de « faire le schéma de principe d'un alternateur ». Il ne s'agit pas d'un dessin artistique, mais d'un \textbf{schéma normalisé} (symboles normalisés NF EN 60617) montrant les trois éléments fonctionnels essentiels.

\begin{methode}[Tracer le schéma de principe d'un alternateur (3 éléments)]
\begin{enumerate}
    \item \textbf{Le Rotor (Inducteur) :} Dessine un \textbf{cercle} (représentant la section du rotor). À l'intérieur, trace le symbole de l'aimant ou de l'électroaimant :
    \begin{itemize}
        \item Aimant permanent : Un rectangle avec \textbf{N} (Nord) et \textbf{S} (Sud) aux extrémités.
        \item Électroaimant : Un bobinage (spire rectangle) traversé par un courant continu $I_e$ (flèche $\rightarrow$) alimenté par une source continue $E_e$ (batterie, redresseur) via deux bagues et balais (souvent omis au profit de la simplicité : on note juste « Excitation $I_e$ »).
    \end{itemize}
    Ajoute une flèche circulaire $\omega$ (ou $n$) pour indiquer la rotation.
    \item \textbf{Le Stator (Induit) :} Dessine \textbf{un (monophasé) ou trois (triphasé) bobinages} (symboles de self/inductance : demi-cercles ou rectangles à traits pointillés) \textbf{autour} du rotor (ou à côté avec flèches de flux). Pour le triphasé, note les bornes \textbf{U, V, W} (et N pour le neutre).
    \item \textbf{La sortie (Circuit extérieur) :} Relie les bornes du stator à une charge $Z$ (récepteur) ou au symbole du réseau.
\end{enumerate}
\end{methode}

\begin{popfigure}
\begin{tikzpicture}[scale=1.1, every node/.style={font=\small}]
    % Stator windings (3 phases)
    \foreach \ang/\col/\lab in {90/popblue/U, 210/popred/V, 330/popgreen/W} {
        \draw[\col, thick, decorate, decoration={snake, amplitude=1.5pt, segment length=5pt}] (\ang:2.5) arc (\ang:\ang+40:2.5);
        \node[\col] at (\ang+20:3.0) {\textbf{\lab}};
    }
    % Neutral point
    \node[draw, circle, fill=popcream, inner sep=1pt] (N) at (0,0) {N};
    \foreach \ang in {90, 210, 330} {
        \draw[thin, dashed] (\ang:2.5) -- (0,0);
    }
    
    % Rotor (Electromagnet)
    \draw[thick, fill=popgray!10] (0,0) circle (1.2cm);
    % Poles N/S
    \draw[thick, fill=popblue!30] (60:0.3) -- ++(60:0.7) arc (60:120:0.7) -- ++(120:-0.7);
    \draw[thick, fill=popred!30] (240:0.3) -- ++(240:0.7) arc (240:300:0.7) -- ++(300:-0.7);
    \node[popblue] at (90:0.7) {\textbf{N}};
    \node[popred] at (270:0.7) {\textbf{S}};
    % Excitation winding symbol
    \draw[poppurple, thick] (-0.6,0) ellipse (0.6 and 0.3);
    \node[poppurple] at (0,0) {$I_e$};
    
    % Rotation arrow
    \draw[->, thick, poporange] (1.8,0) arc (0:180:1.8);
    \node[poporange] at (0,2.0) {$\omega$ (Rotation)};
    
    % Labels
    \node[align=center, anchor=south] at (0, -2.8) {\textbf{Stator (Induit) : 3 enroulements décalés de 120°}};
    \node[align=center, anchor=north] at (0, 2.3) {\textbf{Rotor (Inducteur) : Électroaimant excité par $I_e$}};
    
    % Output terminals
    \foreach \ang/\lab in {90/U, 210/V, 330/W} {
        \draw[thick] (\ang:3.2) -- ++(\ang:0.8) node[anchor=\ang+180] {\textbf{\lab}};
    }
    \draw[thick] (0,-0.15) -- (0,-1.0) node[below] {\textbf{N}};
\end{tikzpicture}
\legende{Schéma de principe normalisé d'un alternateur triphasé à rotor bobiné (type centrale/voiture). Le rotor (aimant tournant N/S) crée le flux $\Phi$ traversant les 3 enroulements statoriques (U, V, W). Sortie triphasée + Neutre.}
\end{popfigure}

\begin{exemplebox}[Identification sur photographie / coupe réelle]
Sur une photo de l'alternateur de la voiture de ton oncle (taximan à Yaoundé) ou d'un démontage en TP :
\begin{enumerate}
    \item \textbf{La poulie / le ventilateur} visible à l'avant $\rightarrow$ Côté \textbf{arbre d'entraînement} (moteur thermique).
    \item \textbf{Le carter arrière} avec des ailettes de refroidissement et \textbf{3 gros fils} qui en sortent (vers le redresseur/batterie) $\rightarrow$ C'est la sortie \textbf{Stator} (courant fort, alternatif avant redressement).
    \item \textbf{Deux petits fils} (souvent un connecteur 2 broches) allant vers un petit boîtier (régulateur) puis vers l'intérieur du carter $\rightarrow$ Alimentation \textbf{Excitation Rotor} (balais/bagues). Courant faible ($\sim 2\ \si{A}$), continu.
    \item Si tu ouvres : le \textbf{rotor à griffes} (pôles N/S alternés) tourne \textbf{dans} le stator (bobines dans les rainures du carter).
\end{enumerate}
\end{exemplebox}

\courssub{Chaîne énergétique locale (rappel et ancrage)}

Avant d'aborder la section 4 (Transformations d'énergie), ancrons les organes dans la chaîne :

\[
\textbf{Source Mécanique} \xrightarrow{\text{Couple } C_m, \text{ Arbre}} \fbox{\textbf{Rotor (Champ } $\vec{B}$\text{)}} \xrightarrow{\text{Induction } e = -\frac{d\Phi}{dt}} \fbox{\textbf{Stator (F.é.m. } e(t)\text{)}} \xrightarrow{\text{Balais/Bagues (si exc.)}} \textbf{Circuit Extérieur}
\]

\begin{savaistu}[L'alternateur de voiture : une « machine inversée » ?]
Dans l'alternateur de ta voiture (ou du taxi « \textit{Clando} »), c'est bien le \textbf{rotor (bobinage d'excitation) qui tourne} et le \textbf{stator (3 phases) qui est fixe}.
Pourquoi ? Parce qu'il est facile de faire passer $\sim 2\ \si{A}$ continu sur des bagues (balais), mais très difficile de récupérer $\sim 50\text{--}100\ \si{A}$ alternatif (courant de charge) sur des bagues sans destruction rapide.
Le courant alternatif triphasé « fort » sort directement du stator vers un \textbf{redresseur à 6 diodes} (pont de Graetz) intégré dans le carter arrière, qui le convertit en \textbf{courant continu} pour charger la batterie $12\ \si{V}$ et alimenter les consommateurs (phares, démarreur, autoradio MTN/Orange). Le régulateur de tension (souvent intégré aux balais) module $I_e$ pour maintenir $U_{\text{batterie}} \approx 14,4\ \si{V}$.
\end{savaistu}

\begin{exoresolu}[Analyse d'un alternateur de centrale - Lom Pangar]
\textbf{Énoncé :} Un alternateur du barrage de Lom Pangar a les caractéristiques nominales suivantes : Puissance apparente $S_n = 30\ \si{MVA}$ ; Tension stator $U_n = 10,5\ \si{kV}$ (triphasé) ; Vitesse de rotation $n = 150\ \si{tr/min}$ ; Fréquence $f = 50\ \si{Hz}$. Le rotor est à pôles saillants bobinés. Le courant d'excitation nominal est $I_{e,n} = 800\ \si{A}$ sous $U_{e,n} = 350\ \si{V}$ continu.
\begin{enumerate}
    \item Déterminer le nombre de pôles $p$ du rotor.
    \item Calculer l'intensité nominale par phase $I_n$ du stator.
    \item Calculer la puissance absorbée par le circuit d'excitation (pertes Joule rotor). Commenter.
    \item Sur une coupe schématique, indiquer le sens du couple électromagnétique $C_{em}$ s'opposant au couple moteur $C_m$ de la turbine.
\end{enumerate}
\tcblower
\textbf{Solution détaillée :}
\begin{enumerate}
    \item Relation fréquence / vitesse / paires de pôles : $f = p \cdot \frac{n}{60}$.
    \[ p = \frac{f \cdot 60}{n} = \frac{50 \times 60}{150} = 20 \text{ paires de pôles} \quad \Rightarrow \quad \textbf{40 pôles (20 Nord, 20 Sud)}. \]
    \textit{Analyse dimensionnelle :} $[f]=\si{s^{-1}}$, $[n]=\si{min^{-1}}$, $[60]=\si{s/min}$ $\rightarrow$ $p$ sans dimension. Correct.
    \item Puissance triphasée : $S_n = \sqrt{3} \cdot U_n \cdot I_n$.
    \[ I_n = \frac{S_n}{\sqrt{3} \cdot U_n} = \frac{30 \times 10^6}{\sqrt{3} \times 10,5 \times 10^3} \approx \frac{30 \times 10^3}{1,732 \times 10,5} \approx \textbf{1 653 A}. \]
    (Courant fort justifiant l'absence de balais sur le stator).
    \item Pertes Joule rotor (excitation) : $P_{\text{exc}} = U_{e,n} \cdot I_{e,n} = 350 \times 800 = 280\,000\ \si{W} = \textbf{280 kW}$.
    \textit{Commentaire :} $280\ \si{kW}$ pour $30\ \si{MVA}$ produit $\rightarrow$ ratio $\approx 0,9\%$. Faible mais non négligeable ; nécessite un système de refroidissement du rotor (ventilation forcée, parfois hydrogène sous pression dans les très grosses machines).
    \item \textbf{Loi de Lenz / Action-Réaction :} La turbine entraîne le rotor dans le sens de rotation (couple moteur $C_m$). Le stator produit un courant qui crée son propre champ magnétique. L'interaction des deux champs génère un \textbf{couple électromagnétique $C_{em}$ de sens opposé à la rotation} (couple freinant, couple de charge). L'équilibre mécanique (vitesse constante) impose $C_m = C_{em} + C_{\text{pertes mécaniques}}$. C'est la conversion $\text{Mécanique} \to \text{Électrique}$.
\end{enumerate}
\end{exoresolu}

\begin{exercice}[Reconnaissance et Schéma]
\begin{enumerate}
    \item Sur la figure ci-dessous (photo d'un alternateur de groupe électrogène diesel ouvert), identifier : (1) Stator, (2) Rotor, (3) Balais, (4) Bagues, (5) Arbre, (6) Carter/Carcasse, (7) Ventilateur de refroidissement.
    \item Tracer le schéma de principe normalisé correspondant (monophasé pour simplifier), en précisant le sens de rotation $\omega$, les pôles N/S du rotor, la bobine du stator et la charge $R$.
    \item Expliquer pourquoi les balais ne chauffent pas excessivement malgré le courant qui les traverse.
\end{enumerate}
\begin{center}
    \textit{(Insérer ici une photo ou un schéma TikZ simplifié d'alternateur ouvert pour l'exercice)}
\end{center}
\end{exercice}

\begin{corrige}[Exercice n°17-3]
\begin{enumerate}
    \item (1) Stator : masse de tôles en couronne avec fils cuivre dans rainures (fixe). (2) Rotor : corps cylindrique à pôles saillants bobinés (tournant). (3) Balais : petits blocs carbone sur porte-balais fixes. (4) Bagues : deux anneaux cuivrés sur l'arbre (tournants). (5) Arbre : axe central. (6) Carter : enveloppe extérieure. (7) Ventilateur : ailettes sur rotor ou poulie.
    \item Schéma : Cercle rotor (N/S + flèche $\omega$) + Bobine stator (self) à côté + Flèches flux $\Phi$ + Charge $R$ aux bornes stator.
    \item Les balais ne portent que le courant d'excitation $I_e$ (continu, faible, $\sim \text{quelques A à centaines d'A}$), pas le courant de charge $I_n$ (alternatif, fort, $\sim \text{kA}$). Pertes Joule $R_{\text{contact}} I_e^2$ faibles.
\end{enumerate}
\end{corrige}

\courssub{Synthèse des parties essentielles}

\begin{center}
\begin{tabularx}{\linewidth}{|l|X|X|}
\hline
\rowcolor{popblueL}
\textbf{Organe} & \textbf{Rôle Physique} & \textbf{Particularité / Point de vigilance} \\
\hline
\textbf{Rotor (Inducteur)} & Créer le champ magnétique tournant $\vec{B}(t)$ & Aimants permanents (vélo) \textbf{OU} Électroaimant $I_e$ (voiture, centrale). \textbf{Tourne}. \\
\hline
\textbf{Stator (Induit)} & Recevoir $\vec{B}(t)$, induire f.é.m. $e(t)$, débiter courant & Bobinages dans tôles (anti-Foucault). \textbf{Fixe}. Sortie puissance (fort $I$). \\
\hline
\textbf{Arbre + Paliers} & Transmettre couple $C_m$, guider rotation & Alignement précis. Vitesse $n = 60f/p$. \\
\hline
\textbf{Balais + Bagues} & Alimenter rotor en $I_e$ (continu) & \textbf{Seulement si rotor bobiné}. Courant faible. Usure lente. \\
\hline
\textbf{Carcasse} & Support, protection (IP), retour flux magnétique & Ferme circuit magnétique. Refroidissement (ailettes, eau, H$_2$). \\
\hline
\end{tabularx}
\end{center}

\begin{aretenir}[Check-list pour l'oral / l'écrit (Probatoire)]
\begin{itemize}
    \item \textbf{Citer} les 5 parties essentielles : Rotor, Stator, Arbre/Paliers, Balais/Bagues (si exc.), Carcasse.
    \item \textbf{Différencier} Rotor (tourne, champ B) et Stator (fixe, f.é.m., courant fort). \emph{Astuce : R = Rotation, S = Stationnaire}.
    \item \textbf{Distinguer} les deux types : Aimants permanents (pas de balais, tension non réglée) vs Électroaimant (balais/bagues pour $I_e$, tension réglée par AVR).
    \item \textbf{Savoir tracer} le schéma de principe : Rotor (aimant N/S + $\omega$) + Stator (bobine(s) U,V,W) + Sortie (Charge/Réseau).
    \item \textbf{Expliquer} pourquoi le courant de puissance ne passe pas par les balais (usure, étincelles, isolation) $\rightarrow$ Stator fixe = sortie directe.
\end{itemize}
\end{aretenir}

Cette connaissance anatomique de l'alternateur est indispensable pour aborder sereinement la section suivante : la \textbf{chaîne des transformations d'énergie} (rendement, pertes, bilan de puissance) et la \textbf{réalisation pratique} d'un modèle réduit.

\coursec{Chaîne de transformations d'énergie}

Après avoir identifié les \textbf{parties essentielles de l'alternateur} (stator, rotor, collecteur/anneaux, balais) et compris comment la rotation de l'inducteur devant l'induit crée une force électromotrice alternative par induction électromagnétique, nous devons maintenant nous poser la question fondamentale de tout ingénieur énergéticien : \emph{d'où vient l'énergie électrique fournie ?} Un alternateur n'est pas une source d'énergie, c'est un \textbf{convertisseur}. Il respecte scrupuleusement le principe de conservation de l'énergie.

\courssub{Bilan énergétique global : convertisseur méca-électrique}

Rappelons le rôle principal établi en Section 1 : l'alternateur transforme l'énergie mécanique de rotation fournie à son arbre (par une turbine, un moteur thermique, une éolienne...) en énergie électrique sous forme de courant alternatif aux bornes du stator.

\begin{propriete}[Conversion d'énergie dans l'alternateur]
L'alternateur est un convertisseur d'énergie \textbf{réversible} (en théorie) qui assure la transformation :
\[
\textbf{Énergie mécanique de rotation } (E_{\text{mec}}) \longrightarrow \textbf{Énergie électrique alternative } (E_{\text{élec}})
\]
Au cours de cette conversion, une partie de l'énergie mécanique reçue est inévitablement dissipée sous forme d'énergie thermique (échauffement) : ce sont les \textbf{pertes}.
\end{propriete}

\courssub{Chaînes complètes selon la source primaire}

Selon l'origine de l'énergie mécanique entraînant l'alternateur, la chaîne de transformations énergétiques complète diffère. Les trois cas majeurs au Cameroun et dans le monde sont illustrés ci-dessous.

\begin{popfigure}
\begin{tikzpicture}[
    node distance=1.2cm and 1.5cm,
    block/.style={rectangle, draw=popdark, thick, fill=popblueL, text width=2.8cm, text centered, rounded corners, minimum height=1.2cm, font=\small},
    arrow/.style={-Stealth, thick, popdark},
    loss/.style={rectangle, draw=poporange, thick, fill=poporangeL, text width=2.5cm, text centered, rounded corners, minimum height=1cm, font=\small\itshape},
]
% Hydro
\node[block, align=center] (hydro1){Énergie potentielle\\de l'eau (barrage)};
\node[block, right=of hydro1, align=center] (hydro2){Énergie cinétique\\de l'eau (tuyau)};
\node[block, right=of hydro2, align=center] (hydro3){Énergie mécanique\\de rotation (turbine)};
\node[block, right=of hydro3, align=center] (hydro4){Énergie électrique\\alternative (alternateur)};
\node[loss, above=of hydro3, align=center] (lossH){Pertes hydrauliques\\+ frottements turbine};

\draw[arrow] (hydro1) -- (hydro2) node[midway, above] {Écoulement};
\draw[arrow] (hydro2) -- (hydro3) node[midway, above] {Turbine};
\draw[arrow] (hydro3) -- (hydro4) node[midway, above] {Alternateur};
\draw[arrow] (hydro3) -- (lossH);

% Thermique (below hydro)
\node[block, below=2.5cm of hydro1, align=center] (therm1){Énergie chimique\\carburant (fioul, gaz)};
\node[block, right=of therm1, align=center] (therm2){Énergie thermique\\combustion};
\node[block, right=of therm2, align=center] (therm3){Énergie mécanique\\moteur/turbine à gaz};
\node[block, right=of therm3, align=center] (therm4){Énergie électrique\\alternative (alternateur)};
\node[loss, below=of therm3, align=center] (lossT){Pertes thermiques\\+ frottements moteur};

\draw[arrow] (therm1) -- (therm2) node[midway, above] {Combustion};
\draw[arrow] (therm2) -- (therm3) node[midway, above] {Moteur thermique};
\draw[arrow] (therm3) -- (therm4) node[midway, above] {Alternateur};
\draw[arrow] (therm3) -- (lossT);

% Eolien (below therm)
\node[block, below=2.5cm of therm1, align=center] (wind1){Énergie cinétique\\du vent};
\node[block, right=of wind1, align=center] (wind2){Énergie mécanique\\rotation pale/rotor};
\node[block, right=of wind2, align=center] (wind3){Énergie électrique\\alternative (alternateur)};
\node[loss, below=of wind2, align=center] (lossW){Pertes aérodynamiques\\+ frottements};

\draw[arrow] (wind1) -- (wind2) node[midway, above] {Éolienne};
\draw[arrow] (wind2) -- (wind3) node[midway, above] {Alternateur};
\draw[arrow] (wind2) -- (lossW);

% Labels
\node[font=\bfseries\color{popblue}] at (hydro1.north) [above=0.2cm] {Centrale hydroélectrique (ex: Songloulou, Nachtigal)};
\node[font=\bfseries\color{poporange}] at (therm1.north) [above=0.2cm] {Centrale thermique / Groupe électrogène (ex: Kribi, Dibamba)};
\node[font=\bfseries\color{popgreen}] at (wind1.north) [above=0.2cm] {Parc éolien (ex: futur parc Nord)};
\end{tikzpicture}
\legende{Chaînes de transformations d'énergie aboutissant à l'alternateur selon la source primaire. Les pertes (orange) surviennent à chaque étape.}
\end{popfigure}

\begin{savaistu}[Le Cameroun, pays de l'hydroélectricité]
Le Cameroun possède un potentiel hydroélectrique énorme (2\textsuperscript{e} d'Afrique après la RDC). Les barrages de \textbf{Songloulou} (\SI{388}{MW}), \textbf{Nachtigal} (\SI{420}{MW}, mis en service récemment), \textbf{Mékin} (\SI{15}{MW}) et le futur \textbf{Grand Eweng} transforment l'énergie potentielle de l'eau des rivières (Sanaga, Ntem) en électricité. L'alternateur est le cœur de cette valorisation : c'est lui qui permet d'injecter sur le réseau SONATREL/ENEO l'énergie qui alimente Yaoundé, Douala, les usines (Alucam) et les foyers. Les grands alternateurs de ces barrages (alternateurs synchrones à pôles saillants) atteignent des rendements supérieurs à \textbf{97\%}.
\end{savaistu}

\courssub{Pertes d'énergie et rendement de l'alternateur}

Toute conversion réelle s'accompagne de pertes. Dans un alternateur, l'énergie mécanique fournie $E_{\text{mec}}$ (ou la puissance mécanique $P_{\text{mec}}$) se partage en énergie électrique utile $E_{\text{élec}}$ ($P_{\text{élec}}$) et en pertes $E_{\text{pertes}}$ ($P_{\text{pertes}}$) qui chauffent la machine.

\begin{definition}[Rendement d'un alternateur]
Le rendement $\eta$ (lettre grecque \emph{êta}) d'un alternateur est le rapport de la puissance électrique utile fournie $P_{\text{élec}}$ (puissance active aux bornes du stator) à la puissance mécanique reçue $P_{\text{mec}}$ sur l'arbre d'entraînement :
\[
\eta = \frac{P_{\text{élec}}}{P_{\text{mec}}}
\]
C'est un nombre sans dimension, toujours \textbf{inférieur à 1} (ou $100\%$). On l'exprime souvent en pourcentage : $\eta_{\%} = \eta \times 100$.
\end{definition}

\begin{attention}[Piège majeur : l'alternateur ne « crée » pas l'énergie !]
\textbf{Erreur fréquente au Probatoire :} écrire « L'alternateur produit de l'énergie » ou « L'alternateur fournit 10 kW d'énergie ».
\begin{itemize}
    \item \textbf{Correct :} « L'alternateur \emph{convertit} de l'énergie mécanique en énergie électrique » ou « L'alternateur \emph{fournit une puissance électrique} de \SI{10}{kW} ».
    \item L'énergie est conservée : $P_{\text{mec}} = P_{\text{élec}} + P_{\text{pertes}}$.
    \item Si l'alternateur est débranché (vide), il tourne presque librement ($P_{\text{mec}} \approx P_{\text{pertes}}$). S'il est chargé, le couple électromagnétique freine le rotor : il faut augmenter le couple moteur (ouvrir les vannes, accélérer le moteur thermique) pour maintenir la vitesse. \textbf{L'énergie électrique fournie provient intégralement de l'énergie mécanique fournie par le premier moteur.}
\end{itemize}
\end{attention}

\paragraph{Origine des pertes (bilan des puissances).}
On distingue trois grandes familles de pertes dans un alternateur réel :
\begin{enumerate}
    \item \textbf{Pertes mécaniques (ou pertes à vide mécaniques) :} Frottements dans les roulements (palier), frottement de l'air sur le rotor (ventilation), frottements balais/anneaux (si machine à bagues). Elles dépendent principalement de la vitesse de rotation $n$.
    \item \textbf{Pertes fer (ou pertes magnétiques) :} Pertes par hystérésis et courants de Foucault dans le noyau magnétique du stator (et du rotor pour les pôles pleins). Elles dépendent de la fréquence $f$ et de l'induction $B$ (donc de la tension $U$). Elles existent dès que la machine tourne sous tension, même à vide.
    \item \textbf{Pertes Joule (ou pertes cuivre) :} Effet Joule dans les bobinages du stator ($3 R I^2$ en triphasé) et du rotor ($R_{\text{ex}} I_{\text{ex}}^2$ pour l'excitation). Elles dépendent du \textbf{carré du courant} (donc de la charge).
\end{enumerate}

\begin{popfigure}
\begin{tikzpicture}[
    scale=0.9,
    every node/.style={font=\small},
    arrow/.style={-Stealth, thick},
    lossarrow/.style={-Stealth, thick, poporange},
]
% Puissance mécanique entrante (grosse flèche verticale)
\draw[arrow, popblue, line width=4pt] (0,-1) -- (0,0) node[midway, left=0.2cm, black] {$P_{\text{mec}}$};
\node[above=0.1cm] at (0,0) [align=center, font=\bfseries] {Alternateur};

% Boîte centrale représentant l'alternateur
\draw[thick, draw=popdark, fill=popcream, rounded corners] (-2.5,0) rectangle (2.5,3.5);

% Pertes mécaniques (sortie gauche bas)
\draw[lossarrow, line width=2pt] (0,0.5) -- (-2.5,0.5) -- (-3,0.5) node[left] {$P_{\text{mec, pertes}}$};
\node[left=0.1cm, rotate=90, anchor=south, poporange] at (-2.5,0.5) {Frottements, ventilation};

% Pertes fer (sortie gauche milieu)
\draw[lossarrow, line width=2pt] (0,1.75) -- (-2.5,1.75) -- (-3,1.75) node[left] {$P_{\text{fer}}$};
\node[left=0.1cm, rotate=90, anchor=south, poporange] at (-2.5,1.75) {Hystérésis + Foucault};

% Pertes Joule stator (sortie gauche haut)
\draw[lossarrow, line width=3pt] (0,3) -- (-2.5,3) -- (-3,3) node[left] {$P_{\text{J, stator}} = 3RI^2$};
\node[left=0.1cm, rotate=90, anchor=south, poporange] at (-2.5,3) {Effet Joule stator};

% Pertes Joule rotor (excitation) - sortie gauche aussi pour cohérence
\draw[lossarrow, line width=1.5pt] (0,2.5) -- (-2.5,2.5) -- (-3,2.5) node[left] {$P_{\text{J, rotor}}$};
\node[left=0.1cm, rotate=90, anchor=south, poporange] at (-2.5,2.5) {Excitation (pertes)};

% Puissance électrique utile (sortie droite, GROSSE flèche)
\draw[arrow, popgreen, line width=3.5pt] (2.5,1.75) -- (4,1.75) node[right] {$P_{\text{élec}} = \sqrt{3}UI\cos\varphi$};
\node[above=0.2cm, popgreen, font=\bfseries] at (3.2,1.75) {Puissance utile};

% Légende épaisseurs
\node[align=left, font=\tiny, text width=3cm] at (-4.5, 2) {Épaisseur $\propto$\\puissance (qualitatif)};
\end{tikzpicture}
\legende{Diagramme de puissances (bilan énergétique) d'un alternateur en charge. L'épaisseur des flèches est proportionnelle à l'ordre de grandeur des puissances. $P_{\text{mec}} = P_{\text{élec}} + \sum P_{\text{pertes}}$.}
\end{popfigure}

\courssub{Valeurs typiques du rendement}

Le rendement dépend de la puissance nominale et de la technologie.
\begin{itemize}
    \item \textbf{Petits alternateurs} (groupes électrogènes portatifs, $< \SI{10}{kW}$) : $\eta \approx 75\% \text{ à } 85\%$.
    \item \textbf{Alternateurs industriels moyens} (\SI{100}{kW} -- \SI{1}{MW}) : $\eta \approx 90\% \text{ à } 94\%$.
    \item \textbf{Gros alternateurs de centrale} ($> \SI{50}{MW}$, ex: Songloulou, Nachtigal) : $\eta \approx 97\% \text{ à } 98,5\%$.
\end{itemize}
Le rendement est maximal au voisinage de la \textbf{charge nominale} (souvent entre $80\%$ et $100\%$ de $P_n$). À vide, $\eta = 0\%$ (toute la puissance mécanique compense les pertes). En surcharge, les pertes Joule ($RI^2$) explosent et le rendement chute.

\courssub{Exemples chiffrés et exercices résolus}

\begin{exemplebox}[Calcul de la puissance électrique fournie]
Un alternateur de groupe électrogène d'hôpital reçoit sur son arbre une puissance mécanique de $P_{\text{mec}} = \SI{10}{kW}$ fournie par le moteur diesel. Son rendement à cette charge est $\eta = 0,95$ ($95\%$). Quelle est la puissance électrique active $P_{\text{élec}}$ disponible pour alimenter les services (bloc opératoire, réfrigération vaccins, éclairage) ?
\end{exemplebox}
\begin{align*}
\text{Données : } & P_{\text{mec}} = \SI{10}{kW} = \SI{10\,000}{W},\quad \eta = {0,95}. \\
\text{Formule : } & \eta = \frac{P_{\text{élec}}}{P_{\text{mec}}} \implies P_{\text{élec}} = \eta \times P_{\text{mec}}. \\
\text{Calcul : } & P_{\text{élec}} = {0,95} \times \SI{10\,000}{W} = \SI{9\,500}{W} = \SI{9,5}{kW}. \\
\text{Pertes : } & P_{\text{pertes}} = P_{\text{mec}} - P_{\text{élec}} = \SI{0,5}{kW} = \SI{500}{W} \text{ (échauffement de la machine)}.
\end{align*}
\begin{aretenir}
L'alternateur fournit \textbf{\SI{9,5}{kW}} de puissance électrique active. Les \textbf{\SI{500}{W}} restants chauffent l'alternateur (ventilation nécessaire).
\end{aretenir}

\begin{exoresolu}[Dimensionnement inverse : puissance mécanique nécessaire]
Un village rural doit être électrifié par un groupe électrogène. La puissance électrique active maximale estimée (pic du soir : éclairage, moulin, pompe, téléphones) est de $P_{\text{élec}} = \SI{50}{kW}$. L'alternateur disponible a un rendement $\eta = 0,90$ ($90\%$) à cette charge. Quelle puissance mécanique minimale $P_{\text{mec}}$ le moteur thermique (diesel) doit-il fournir sur l'arbre de l'alternateur ? Quelle puissance le moteur doit-il développer en tenant compte d'une marge de sécurité de $10\%$ ?
\tcblower
\textbf{Solution détaillée.}
\begin{enumerate}
    \item \textbf{Puissance mécanique strictement nécessaire à l'alternateur :}
    \[
    P_{\text{mec}} = \frac{P_{\text{élec}}}{\eta} = \frac{\SI{50}{kW}}{{0,90}} \approx \SI{55,56}{kW}.
    \]
    \item \textbf{Marge de sécurité ($10\%$) :} Le moteur ne doit pas tourner à son couple maximal en permanence. On choisit un moteur de puissance nominale $P_{\text{moteur}} \ge 1,1 \times P_{\text{mec}}$.
    \[
    P_{\text{moteur}} \ge 1,1 \times \SI{55,56}{kW} \approx \SI{61,1}{kW}.
    \]
    \item \textbf{Choix standardisé :} On sélectionnera un moteur diesel de \textbf{\SI{65}{kW}} (ou \SI{70}{kW}) pour assurer la tenue en charge et l'altitude/chaleur.
\end{enumerate}
\end{exoresolu}

\begin{aretenir}
Pour fournir \SI{50}{kW} électriques avec $\eta = 90\%$, il faut \textbf{$\approx \SI{55,6}{kW}$} mécaniques sur l'arbre. Le moteur thermique doit être dimensionné à \textbf{$\ge \SI{61}{kW}$} (choix standard \SI{65}{kW}).
\end{aretenir}

\courssub{Synthèse : de la source à la prise}

\begin{aretenir}[Bilan énergétique global à retenir]
\begin{itemize}
    \item L'alternateur est un \textbf{convertisseur} : $E_{\text{mec}} \to E_{\text{élec}}$. Il ne crée pas l'énergie.
    \item Chaîne type \textbf{barrage} : $E_{\text{pot, eau}} \to E_{\text{cin, eau}} \to E_{\text{mec, turbine}} \xrightarrow{\text{alternateur}} E_{\text{élec}}$.
    \item Chaîne type \textbf{thermique} : $E_{\text{chim, carburant}} \xrightarrow{\text{combustion}} E_{\text{therm}} \to E_{\text{mec, moteur}} \xrightarrow{\text{alternateur}} E_{\text{élec}}$.
    \item Chaîne type \textbf{éolien} : $E_{\text{cin, vent}} \to E_{\text{mec, pale}} \xrightarrow{\text{alternateur}} E_{\text{élec}}$.
    \item Rendement $\eta = P_{\text{élec}} / P_{\text{mec}} < 1$. Pertes = mécaniques + fer + Joule $\to$ chaleur.
    \item $\eta$ typique : $75\%$ (petit) à $> 97\%$ (gros barrage).
    \item Relation fondamentale : $P_{\text{mec}} = P_{\text{élec}} + P_{\text{pertes}}$.
\end{itemize}
\end{aretenir}

\begin{methode}[Résoudre un problème de bilan de puissance sur alternateur]
\begin{enumerate}
    \item \textbf{Identifier} les puissances connues : $P_{\text{mec}}$ (moteur/turbine) ou $P_{\text{élec}}$ (charge), et le rendement $\eta$.
    \item \textbf{Écrire} la définition : $\eta = P_{\text{élec}} / P_{\text{mec}}$.
    \item \textbf{Isoler} l'inconnue : $P_{\text{élec}} = \eta P_{\text{mec}}$ ou $P_{\text{mec}} = P_{\text{élec}} / \eta$.
    \item \textbf{Calculer} les pertes si demandé : $P_{\text{pertes}} = P_{\text{mec}} - P_{\text{élec}} = P_{\text{mec}}(1-\eta)$.
    \item \textbf{Vérifier} l'homogénéité (Watts) et l'ordre de grandeur ($\eta < 1 \implies P_{\text{élec}} < P_{\text{mec}}$).
    \item \textbf{Conclure} par une phrase : « La puissance électrique fournie est de ... » / « Le moteur doit fournir au moins ... ».
\end{enumerate}
\end{methode}

\begin{exercice}[Application : Rendement et pertes d'un alternateur de barrage]
L'alternateur n°3 du barrage de \textbf{Songloulou} a une puissance nominale apparente $S_n = \SI{100}{MVA}$ et un facteur de puissance $\cos\varphi = 0,9$. Son rendement à la charge nominale est $\eta = 0,975$.
\begin{enumerate}
    \item Calculer la puissance électrique active nominale $P_n$ (en MW).
    \item Calculer la puissance mécanique $P_{\text{mec}}$ que la turbine Francis doit fournir sur l'arbre à la charge nominale (en MW).
    \item Calculer la puissance totale des pertes $P_{\text{pertes}}$ dans l'alternateur à cette charge (en MW). Commenter l'ordre de grandeur.
\end{enumerate}
\end{exercice}

\begin{corrige}[Exercice : Rendement et pertes d'un alternateur de barrage]
\begin{enumerate}
    \item $P_n = S_n \times \cos\varphi = \SI{100}{MVA} \times 0,9 = \SI{90}{MW}$.
    \item $P_{\text{mec}} = \dfrac{P_n}{\eta} = \dfrac{\SI{90}{MW}}{{0,975}} \approx \SI{92,31}{MW}$.
    \item $P_{\text{pertes}} = P_{\text{mec}} - P_n = \SI{92,31}{MW} - \SI{90}{MW} = \SI{2,31}{MW}$.
    \textbf{Commentaire :} Bien que le rendement soit excellent ($97,5\%$), les pertes atteignent \textbf{\SI{2,31}{MW}} soit l'équivalent de la consommation de plusieurs milliers de foyers ! Cette énergie est dissipée en chaleur : il faut un système de refroidissement performant (air/eau/hydrogène) pour évacuer ces \SI{2,31}{MW} thermiques sans détruire l'isolation des bobinages.
\end{enumerate}
\end{corrige}

\coursec{Exemples d'alternateurs et applications}

Dans la section précédente, nous avons établi la \textbf{chaîne de transformations d'énergie} universelle à tout alternateur : énergie primaire $\rightarrow$ énergie mécanique (rotation) $\rightarrow$ énergie électrique alternative. Nous allons maintenant décliner ce schéma universel en examinant les principaux types d'alternateurs rencontrés, du plus modeste au plus puissant, en portant une attention particulière au contexte énergétique camerounais.

\courssub{De la dynamo de vélo aux centrales géantes : une même physique, des échelles différentes}

Rappelons le principe fondamental : un alternateur est une machine tournante qui convertit l'énergie mécanique fournie à son arbre (le rotor) en énergie électrique aux bornes de ses enroulements (le stator), grâce à l'induction électromagnétique ($e = - \frac{\mathrm{d}\Phi}{\mathrm{d}t}$). La nature de la \textbf{source d'énergie mécanique} (musculaire, thermique, hydraulique, éolienne) et la \textbf{puissance} en jeu sont les deux critères qui distinguent les applications.

\begin{popfigure}
\centering
\begin{tikzpicture}
\begin{axis}[
    width=\linewidth,
    height=0.55\linewidth,
    xmode=log,
    log basis x=10,
    xmin=1, xmax=1e9,
    ytick={1,2,3,4,5,6},
    yticklabels={%
        {Dynamo vélo},
        {Alternateur auto},
        {Groupe électrogène},
        {Éolienne terrestre},
        {Centrale thermique},
        {Barrage hydroélectrique}
    },
    xlabel={Puissance nominale (W) -- échelle logarithmique},
    grid=major,
    major grid style={popmuted},
    axis y line*=left,
    axis x line*=bottom,
    tick label style={font=\small},
    label style={font=\small},
]
% Barres de puissance (plage min-max)
\addplot[popblue!60, fill=popblue!20, very thick, error bars/.cd, x dir=both, x explicit, error bar style={popblue!80, thick}] coordinates {
    (10,1) +- (9,11)       % Dynamo 1-100 W
    (500,2) +- (400,600)   % Auto 100-1000 W
    (5000,3) +- (1000,20000) % Groupe 1-20 kW
    (2e6,4) +- (5e5,5e6)   % Éolienne 0.5-5 MW
    (2e8,5) +- (5e7,5e8)   % Thermique 50-500 MW
    (3e8,6) +- (1e8,1e9)   % Hydro 100-1000 MW
};
% Points nominaux
\addplot[only marks, mark=*, mark size=3, popblue] coordinates {
    (10,1) (500,2) (5000,3) (2e6,4) (2e8,5) (3e8,6)
};
% Lignes de référence
\draw[poporange, dashed, thick] ({1e3},0.5) -- ({1e3},6.5) node[above, font=\tiny, poporange] {1 kW};
\draw[poporange, dashed, thick] ({1e6},0.5) -- ({1e6},6.5) node[above, font=\tiny, poporange] {1 MW};
\draw[poporange, dashed, thick] ({1e9},0.5) -- ({1e9},6.5) node[above, font=\tiny, poporange] {1 GW};
\end{axis}
\end{tikzpicture}
\legende{Échelle des puissances des principaux types d'alternateurs (échelle logarithmique). Notez l'écart de 8 ordres de grandeur entre la dynamo et les grands barrages.}
\end{popfigure}

\courssub{Alternateurs de faible puissance : usages individuels et de proximité}

\courssub{La dynamo de vélo : l'alternateur à aimant permanent}

C'est l'exemple le plus accessible. Sur un vélo, une petite roue (le rotor) frottant sur le pneu (ou intégrée au moyeu) tourne. À l'intérieur, un \textbf{aimant permanent} tourne face à un enroulement fixe (stator). Pas de balais, pas d'excitation séparée : le champ magnétique est constant, fourni par l'aimant.

\begin{propriete}[Dynamo de vélo (alternateur à aimant permanent)]
\begin{itemize}
    \item \textbf{Source mécanique :} force musculaire du cycliste (vitesse angulaire $\omega$ proportionnelle à la vitesse du vélo).
    \item \textbf{Champ magnétique :} aimants permanents (néodyme ou ferrite) sur le rotor.
    \item \textbf{Tension de sortie :} alternative, amplitude $U_{\text{max}} \propto \omega$. Typiquement \SI{6}{\volt} efficaces à vitesse de croisière ($\approx \SI{15}{\kilo\meter\per\hour}$).
    \item \textbf{Puissance :} quelques watts (\SIrange{3}{10}{\watt}), suffisante pour l'éclairage (LED modernes).
\end{itemize}
\end{propriete}

\begin{attention}[Dynamo $\neq$ Alternateur de voiture]
Ne confondez pas la \textbf{dynamo de vélo} (aimant permanent, tension $\propto$ vitesse, faible puissance) et l'\textbf{alternateur de voiture} (électroaimant excité, tension régulée à \SI{14,4}{\volt} quelle que soit la vitesse moteur, puissance $\sim \SI{1}{\kilo\watt}$). La dynamo est un \emph{alternateur} (courant alternatif redressé en interne ou non), mais le terme « dynamo » désigne historiquement une machine à courant continu (avec collecteur). Au Probatoire, on parle d'\textbf{alternateur} pour toute machine produisant du courant alternatif.
\end{attention}

\courssub{L'alternateur de voiture : régulation de tension et forte puissance}

L'alternateur automobile est un alternateur \textbf{à excitation séparée} (électroaimant au rotor). Il est entraîné par le moteur thermique via une courroie trapézoïdale (rapport de démultiplication $\approx 1:3$ à $1:4$ : l'alternateur tourne plus vite que le moteur).

\begin{experience}[Démontage / observation d'un alternateur automobile (TP ou vidéo)]
\textbf{Matériel :} alternateur de récupération, multimètre, lampe témoin \SI{12}{\volt}/\SI{21}{\watt}.
\textbf{Observations :}
\begin{enumerate}
    \item \textbf{Rotor (induit) :} bobine traversante alimentée par deux bagues collectrices et balais (courant d'excitation $I_{\text{exc}} \sim \SIrange{2}{3}{\ampere}$). C'est un électroaimant.
    \item \textbf{Stator (induit) :} trois enroulements décalés de $120^\circ$ (trois phases) $\rightarrow$ courant triphasé.
    \item \textbf{Pont de diodes (redresseur) :} 6 diodes intégrées $\rightarrow$ sortie \textbf{continue} ($+$ et $-$).
    \item \textbf{Régulateur de tension :} boîtier électronique limitant la tension de sortie à $U_{\text{reg}} \approx \SI{14,4}{\volt}$ en agissant sur $I_{\text{exc}}$.
\end{enumerate}
\textbf{Interprétation :} Le régulateur assure que la batterie (\SI{12}{\volt}) soit rechargée sans surcharge, quel que soit le régime moteur (ralenti ou autoroute). C'est un \textbf{asservissement de tension}.
\end{experience}

\begin{aretenir}[Alternateur automobile]
\begin{itemize}
    \item \textbf{Entrée :} arbre entraîné par courroie ($\omega_{\text{alt}} \approx 3 \times \omega_{\text{moteur}}$).
    \item \textbf{Excitation :} courant continu dans le rotor (électroaimant) $\rightarrow$ champ $B$ variable piloté.
    \item \textbf{Sortie stator :} 3 phases alternatives $\xrightarrow{\text{pont de diodes}}$ tension continue $\approx \SI{14,4}{\volt}$.
    \item \textbf{Puissance :} \SIrange{500}{2000}{\watt} (courant \SIrange{40}{150}{\ampere}).
    \item \textbf{Rôle :} recharger la batterie (démarrage, éclairage, injection, clim, multimédia) \textbf{moteur tournant}.
\end{itemize}
\end{aretenir}

\courssub{Le groupe électrogène : autonomie et délestage (contexte camerounais)}

Au Cameroun, le \textbf{délestage} (coupures programmées ou imprévues du réseau ENEO/SONATREL) rend le groupe électrogène indispensable pour les hôpitaux, les entreprises, les administrations et de nombreux foyers.

\begin{definition}[Groupe électrogène]
Ensemble couplé rigidement : \textbf{Moteur thermique} (essence $\le \SI{10}{\kilo\watt}$, diesel/gasoil $> \SI{10}{\kilo\watt}$) + \textbf{Alternateur synchrone} (généralement triphasé \SI{400}{\volt} ou monophasé \SI{230}{\volt}).
\end{definition}

\begin{methode}[Analyser un groupe électrogène]
\textbf{Étape 1 :} Identifier la source primaire $\rightarrow$ carburant (énergie chimique).
\textbf{Étape 2 :} Identifier la conversion thermique $\rightarrow$ moteur (énergie mécanique, arbre tournant à \SI{1500}{\tr\per\minute} pour \SI{50}{\hertz} ou \SI{1800}{\tr\per\minute} pour \SI{60}{\hertz}).
\textbf{Étape 3 :} Identifier l'alternateur $\rightarrow$ conversion mécanique $\rightarrow$ électrique (régulation tension/fréquence via régulateur AVR \emph{Automatic Voltage Regulator} et régulateur de vitesse moteur).
\textbf{Étape 4 :} Noter les protections : disjoncteur thermique, différentiel, mise à la terre (piquet).
\end{methode}

\begin{exemplebox}[Groupe électrogène domestique typique au Cameroun]
\textbf{Caractéristiques plaque signalétique :}
\begin{itemize}
    \item Moteur : 4 temps, essence, \SI{6,5}{\kilo\watt} ($\approx \SI{8,8}{\ch}$) à \SI{3600}{\tr\per\minute}.
    \item Alternateur : monophasé, \SI{230}{\volt}, \SI{50}{\hertz}, \SI{5}{\kilo\volt\ampere} (facteur de puissance $\cos\varphi \approx 0,8$ $\rightarrow$ $P \approx \SI{4}{\kilo\watt}$ utiles).
    \item Réservoir : \SI{15}{\liter} $\rightarrow$ autonomie $\approx \SIrange{8}{10}{\hour}$ à $50\%$ charge.
    \item Rendement global (chimique $\rightarrow$ électrique) : $\eta \approx 20\text{--}25\%$.
\end{itemize}
\textbf{Calcul :} Consommation horaire $\approx \SI{1,5}{\liter\per\hour}$ à mi-charge. Coût horaire $\approx 1,5 \times \SI{850}{\franccfa} \approx \SI{1275}{\franccfa\per\hour}$ (prix essence 2024). Coût du kWh produit $\approx \SI{320}{\franccfa\per\kilo\watt\hour}$ (vs $\approx \SIrange{50}{80}{\franccfa\per\kilo\watt\hour}$ réseau ENEO). \emph{D'où l'intérêt de l'optimisation énergétique (LED, onduleur, solaire).}
\end{exemplebox}

\courssub{Alternateurs de forte puissance : production centralisée (Centrales thermiques et hydroélectriques)}

C'est le cœur du système électrique national (SEN) géré par SONATREL (transport) et ENEO (distribution). La puissance unitaire d'un alternateur de centrale atteint plusieurs centaines de \textbf{Mégawatts (MW)}.

\courssub{Turboalternateurs des centrales thermiques (gaz, fuel, charbon)}

\begin{propriete}[Centrale thermique à turbine (cycle simple ou combiné)]
\begin{itemize}
    \item \textbf{Source :} combustion gaz naturel (Logbaba, Kribi), fuel lourd (Dibamba, Limbe), charbon (projet).
    \item \textbf{Conversion :} Chaudière $\rightarrow$ vapeur surchauffée ($T \approx \SI{540}{\celsius}, P \approx \SI{150}{\bar}$) $\rightarrow$ \textbf{Turbine à vapeur} (dilatation) $\rightarrow$ Arbre alternateur.
    \item \textbf{Alternateur :} Synchrone, triphasé, refroidi à l'hydrogène (haute conductivité thermique) ou à l'eau (stator). Vitesse de rotation synchronisée sur la fréquence réseau : \SI{3000}{\tr\per\minute} (\SI{50}{\hertz}, $p=1$ paire de pôles) ou \SI{1500}{\tr\per\minute} ($p=2$).
    \item \textbf{Puissance unitaire :} \SIrange{50}{300}{\mega\watt}.
    \item \textbf{Rendement :} Cycle simple $\approx 35\text{--}40\%$ ; Cycle combiné (gaz + vapeur) $\approx 55\text{--}60\%$.
\end{itemize}
\end{propriete}

\courssub{Alternateurs hydroélectriques : le pilier camerounais}

Le Cameroun possède le \textbf{deuxième potentiel hydroélectrique d'Afrique} (après la RDC), estimé à $\approx \SI{23}{\giga\watt}$ de puissance installable, pour une production annuelle théorique $> \SI{100}{\tera\watt\hour}$. Seule une fraction est exploitée.

\begin{savaistu}[Le géant endormi : le potentiel hydroélectrique camerounais]
Avec ses grands fleuves (Sanaga, Benoué, Logone, Nyong) et son relief, le Cameroun pourrait théoriquement alimenter toute l'Afrique centrale. Pourtant, la puissance installée hydroélectrique effective n'est que d'environ \SI{1,5}{\giga\watt} ($< 7\%$ du potentiel). Les barrages de \textbf{Songloulou}, \textbf{Édéa}, \textbf{Lagdo} et \textbf{Lom Pangar} sont les piliers actuels. Le projet \textbf{Nachtigal} (\SI{420}{\mega\watt}) et le futur \textbf{Grand Eweng} (\SI{1000}{\mega\watt}) visent à combler le déficit structurel et à réduire la part thermique (coûteuse en devises pour l'import de fuel/gaz).
\end{savaistu}

\begin{exemplebox}[Centrales de Songloulou et d'Édéa : fierté nationale]
Sur la \textbf{Sanaga}, deux centrales historiques équipées d'alternateurs synchrones couplés à des turbines Francis (actionnées par l'eau sous pression) :
\begin{itemize}
    \item \textbf{Songloulou (mis en service 1980/1988) :} $4$ groupes $\times \SI{96}{\mega\watt}$ = \textbf{\SI{384}{\mega\watt}}. Chute nette $\approx \SI{60}{\meter}$. Débit turbiné $\approx \SI{700}{\meter^3\per\second}$.
    \item \textbf{Édéa (1$^{\text{re}}$ centrale, 1974, réhabilitée) :} $4$ groupes $\times \SI{66}{\mega\watt}$ = \textbf{\SI{264}{\mega\watt}}. Chute nette $\approx \SI{18}{\meter}$ (basse chute, gros débit $\approx \SI{2000}{\meter^3\per\second}$). Turbulences cavitationnelles surveillées de près.
\end{itemize}
Ces alternateurs tournent à \textbf{\SI{150}{\tr\per\minute}} ($p=20$ paires de pôles pour \SI{50}{\hertz}) : ce sont des machines \textbf{lentes, massives} (rotor à pôles saillants, diamètre $\sim \SI{10}{\meter}$, masse $\sim \SI{200}{\tonne}$), refroidies à l'air/eau. Leur inertie énorme assure la \textbf{stabilité de fréquence} du réseau (fréquence $f = \frac{p \cdot n}{60}$).
\end{exemplebox}

\begin{popfigure}
\centering
\begin{tikzpicture}[scale=0.9, every node/.style={font=\small}]
% Couleurs
\colorlet{water}{popblue!40}
\colorlet{rock}{popmuted!50}
\colorlet{pipe}{popdark}
\colorlet{machine}{poporange!80}
\colorlet{elec}{poppurple!80}

% Barrage
\draw[fill=rock, draw=popdark, thick] (0,0) -- (0,5) -- (1,5.5) -- (4,5.5) -- (5,5) -- (5,0) -- cycle;
\draw[fill=water, draw=popblue, thick] (0,5) -- (1,5.5) -- (4,5.5) -- (5,5) -- (5,3) -- (0,3) -- cycle;
\node[align=center, popdark] at (2.5, 5.8) {\textbf{Barrage}\\(Retenue d'eau)};
\node[align=center, popblue] at (2.5, 4.2) {Niveau amont\\$H_{\text{amont}}$};

% Conduite forcée
\draw[fill=pipe, draw=popdark, thick] (4.5, 5.5) -- (4.5, 0.5) -- (6.5, 0.5) -- (6.5, 5.5) -- cycle;
\draw[->, thick, popblue] (5.5, 4.5) -- (5.5, 2.5) node[midway, right] {$Q, \Delta P$};
\node[align=center, popdark] at (5.5, 5.9) {\textbf{Conduite forcée}\\(Pénstock)};

% Salle des machines
\draw[fill=popcream, draw=popdark, thick] (6.5, -0.5) rectangle (12, 5.5);
% Turbine
\draw[fill=machine, draw=poporange, thick] (8, 0.5) circle (1.2);
\node[align=center, poporange] at (8, 0.5) {\textbf{Turbine}\\(Francis/Kaplan)};
\draw[->, thick, poporange] (9.2, 0.5) -- (10, 0.5) node[midway, above, align=center]{Arbre\\$P_{\text{meca}}$};
% Alternateur
\draw[fill=elec, draw=poppurple, thick] (10.5, 0.5) circle (1.2);
\node[align=center, poppurple] at (10.5, 0.5) {\textbf{Alternateur}\\(Synchrone)};
% Transformateur
\draw[fill=popgreen!60, draw=popgreen, thick] (10.5, 3) rectangle (11.5, 4.5);
\node[align=center, popgreen] at (11, 3.75) {\textbf{Transformateur}\\(Élévateur)};
\draw[->, thick, poppurple] (10.5, 1.7) -- (10.5, 3) node[midway, right] {$U_{\text{gen}} \sim \SIrange{10}{15}{\kilo\volt}$};
\draw[->, thick, popgreen] (11, 4.5) -- (11, 5.5) node[midway, right] {$U_{\text{HT}} \sim \SIrange{225}{400}{\kilo\volt}$};

% Ligne HT
\draw[thick, popline] (11, 5.5) -- (13, 5.5);
\draw[thick, popline] (11, 5.5) -- (11, 6.5) -- (13, 6.5) -- (13, 5.5);
\node[align=center, popdark] at (13.5, 6) {\textbf{Ligne HT}\\(Transport SONATREL)};

% Aval
\draw[fill=water, draw=popblue, thick] (6.5, -0.5) -- (12, -0.5) -- (12, 0) -- (6.5, 0) -- cycle;
\node[align=center, popblue] at (9.25, -0.7) {Canal de fuite\\$H_{\text{aval}}$};

% Légende énergie
\node[align=left, font=\tiny, draw=popdark, fill=popcream, rounded corners, inner sep=3pt] at (14.5, 2.5) {
\textbf{Chaîne énergétique :}\\
$\text{Potentielle } (mgh)$\\
$\downarrow$ \textit{Conduite forcée}\\
$\text{Cinétique/Pression}$\\
$\downarrow$ \textit{Turbine}\\
$\text{Mécanique } (\omega, \Gamma)$\\
$\downarrow$ \textit{Alternateur}\\
$\text{Électrique } (U, I, f)$\\
$\downarrow$ \textit{Transformateur}\\
$\text{HT } (U \uparrow, I \downarrow)$
};
\end{tikzpicture}
\legende{Schéma de principe d'une centrale hydroélectrique au fil de l'eau ou à retenue (type Songloulou/Édéa). L'alternateur est couplé directement à la turbine (arbre horizontal ou vertical). Le transformateur élève la tension pour le transport en Haute Tension (HT).}
\end{popfigure}

\courssub{Les éoliennes : alternateurs à vitesse variable}

L'énergie cinétique du vent fait tourner une hélice (rotor aérodynamique). La puissance éolienne disponible est $P_{\text{vent}} = \frac{1}{2} \rho S v^3$. La vitesse de rotation $\omega_{\text{hélice}}$ varie avec le vent ($v \approx \SIrange{4}{25}{\meter\per\second}$). Or, le réseau impose $f = \SI{50}{\hertz}$ constante.

\begin{propriete}[Chaîne de conversion éolienne moderne]
Deux architectures principales :
\begin{enumerate}
    \item \textbf{À vitesse fixe (ancien) :} Multiplicateur de vitesse (boîte de vitesses) + Alternateur synchrone (ou asynchrone) raccordé directement au réseau. $\omega_{\text{alt}} = \text{cste}$. Perte dans le multiplicateur, contrainte mécanique forte.
    \item \textbf{À vitesse variable (standard actuel) :} Multiplicateur (parfois supprimé \emph{direct drive}) + \textbf{Alternateur asynchrone à double alimentation (DFIG)} ou \textbf{Alternateur synchrone à aimants permanents (PMSG)} + \textbf{Convertisseur de puissance (électronique de puissance)}. L'alternateur tourne à $\omega_{\text{opt}}(v)$ pour maximiser $C_p$ (coefficient de puissance), le convertisseur redresse et synchronise sur le réseau (\SI{50}{\hertz}, $U_{\text{réseau}}$).
\end{enumerate}
\textbf{Puissance unitaire :} \SIrange{2}{8}{\mega\watt} (terrestre), \SIrange{10}{15}{\mega\watt} (offshore).
\end{propriete}

\begin{attention}[Alternateur $\neq$ Moteur électrique -- Piège classique au Probatoire]
\begin{itemize}
    \item \textbf{Alternateur (Générateur) :} \textbf{Énergie mécanique $\rightarrow$ Énergie électrique}. Couple électromagnétique \textbf{résistant} (s'oppose à la rotation). Puissance mécanique $P_{\text{meca}} = \Gamma \omega > 0$ (fourni par la turbine/moteur). Puissance électrique $P_{\text{elec}} = U I \cos\varphi > 0$ (sortie).
    \item \textbf{Moteur électrique :} \textbf{Énergie électrique $\rightarrow$ Énergie mécanique}. Couple électromagnétique \textbf{moteur} (entraîne la rotation). $P_{\text{elec}}$ absorbée, $P_{\text{meca}}$ fournie à la charge.
    \item \textbf{Réversibilité :} La \textbf{même machine} peut fonctionner dans les deux sens (ex: démarreur-alternateur hybride, freinage régénératif voiture électrique/TGV). Mais dans une centrale, l'alternateur \textbf{ne fonctionne qu'en générateur}.
\end{itemize}
\textbf{Astuce mnémotechnique :} \textbf{AL}ternateur = \textbf{AL}imentation (il alimente le réseau). \textbf{M}oteur = \textbf{M}ouvement (il crée le mouvement).
\end{attention}

\courssub{Tableau comparatif synthétique}

\begin{popfigure}
\centering
\begin{tikzpicture}
\node[inner sep=0, align=center]{
\begin{tabularx}{\linewidth}{|>{\bfseries}l|X|X|X|X|}
\hline
\rowcolor{popblueL}
\textbf{Type d'alternateur} & \textbf{Source d'énergie mécanique} & \textbf{Puissance typique} & \textbf{Particularités techniques} & \textbf{Exemple camerounais} \\
\hline
Dynamo vélo & Musculaire (cycliste) & \SIrange{1}{10}{\watt} & Aimant permanent, mono/non régulé, \SI{6}{\volt} & Éclairage vélo (Yaoundé, Douala) \\
\hline
Alternateur auto & Moteur essence/diesel (courroie) & \SIrange{0,5}{2}{\kilo\watt} & Excitation pilotée, redresseur + régulateur \SI{14,4}{\volt} CC & Parc automobile ($> 1\,\text{M}$ véhicules) \\
\hline
Groupe électrogène & Moteur diesel/essence (arbre direct) & \SIrange{2}{2000}{\kilo\watt} & Synchrone, AVR, régulateur vitesse, autonomie carburant & Hôpitaux, usines, domiciles (délestage) \\
\hline
Turboalternateur thermique & Turbine vapeur/gaz (combustion) & \SIrange{50}{300}{\mega\watt} & Synchrone, refroidi H$_2$/eau, \SI{3000}{\tr\per\minute} ($p=1$) & Centrale de Kribi (gaz, \SI{216}{\mega\watt}), Dibamba (fuel, \SI{88}{\mega\watt}) \\
\hline
Alternateur hydroélectrique & Turbine hydraulique (eau sous pression) & \SIrange{50}{400}{\mega\watt}/groupe & Synchrone lent (\SIrange{150}{300}{\tr\per\minute}), pôles saillants, inertie forte & \textbf{Songloulou (\SI{384}{\mega\watt})}, \textbf{Édéa (\SI{264}{\mega\watt})}, Lagdo (\SI{72}{\mega\watt}), Lom Pangar (\SI{30}{\mega\watt}), Nachtigal (\SI{420}{\mega\watt}) \\
\hline
Alternateur éolien & Hélice (vent) + Multiplicateur/Convertisseur & \SIrange{2}{15}{\mega\watt} & Vitesse variable, PMSG ou DFIG + électronique puissance & Projets Nord (Mokolo, Figuil) -- potentiel $\sim \SI{1}{\giga\watt}$ \\
\hline
\end{tabularx}
};
\end{tikzpicture}
\legende{Tableau comparatif des principaux types d'alternateurs. La puissance s'échelonne sur 8 ordres de grandeur. Au Cameroun, l'hydroélectricité domine le mix de production (environ 70\% de l'énergie produite), le thermique (gaz/fuel) assure le complément et la pointe, les groupes électrogènes pallient les défaillances du réseau.}
\end{popfigure}

\courssub{Complément Probatoire : Le transport de l'énergie électrique en Haute Tension (HT)}

Quelle que soit la puissance de l'alternateur (de \SIrange{10}{25}{\kilo\volt} typiquement en sortie de machine), l'énergie doit être transportée sur des centaines de kilomètres (ex: Songloulou $\rightarrow$ Yaoundé $\approx \SI{200}{\kilo\meter}$ ; Nachtigal $\rightarrow$ Yaoundé $\approx \SI{50}{\kilo\meter}$) jusqu'aux postes de distribution.

\begin{propriete}[Nécessité de la Haute Tension (HT / THT)]
Soit une puissance active $P$ à transporter sur une ligne de résistance $R_{\text{ligne}}$ (par phase). Le courant de ligne est $I = \frac{P}{\sqrt{3} U \cos\varphi}$ (trois phases). Les pertes par effet Joule dans la ligne (trois phases) sont :
\[ P_{\text{pertes}} = 3 R_{\text{ligne}} I^2 = 3 R_{\text{ligne}} \frac{P^2}{3 U^2 \cos^2\varphi} = R_{\text{ligne}} \frac{P^2}{U^2 \cos^2\varphi} \]
\textbf{Pour une puissance $P$ donnée, les pertes sont inversement proportionnelles au carré de la tension $U$ de transport.}
\end{propriete}

\begin{exoresolu}[Calcul des pertes sur une ligne de transport]
\textbf{Énoncé :} La centrale de Songloulou injecte $P = \SI{300}{\mega\watt}$ sur le réseau \SI{225}{\kilo\volt} (tension composée efficace). La ligne vers Yaoundé a une résistance équivalente par phase $R_{\text{ligne}} = \SI{15}{\ohm}$. Facteur de puissance $\cos\varphi = 0,95$. Calculer le pourcentage de pertes. Que deviendrait-il si on transportait à \SI{30}{\kilo\volt} (tension générateur) ?

\tcblower

\textbf{Solution :}
\begin{enumerate}
    \item \textbf{Courant à 225 kV :} $I_{225} = \frac{P}{\sqrt{3} U_{225} \cos\varphi} = \frac{300 \times 10^6}{\sqrt{3} \times 225 \times 10^3 \times 0,95} \approx \frac{300 \times 10^6}{369,8 \times 10^3} \approx \SI{811}{\ampere}$.
    \item \textbf{Pertes à 225 kV :} $P_{\text{pertes},225} = 3 R_{\text{ligne}} I_{225}^2 = 3 \times 15 \times (811)^2 \approx 45 \times 657\,721 \approx \SI{29,6}{\mega\watt}$.
    \item \textbf{Pourcentage :} $\frac{29,6}{300} \times 100 \approx \mathbf{9,9\%}$. (Réel : lignes doubles, $R$ plus faible $\rightarrow \sim 3\text{--}5\%$).
    \item \textbf{Scénario à 30 kV :} $I_{30} = I_{225} \times \frac{225}{30} = 811 \times 7,5 \approx \SI{6083}{\ampere}$.
    \item \textbf{Pertes à 30 kV :} $P_{\text{pertes},30} = P_{\text{pertes},225} \times \left(\frac{225}{30}\right)^2 = 29,6 \times (7,5)^2 \approx 29,6 \times 56,25 \approx \mathbf{\SI{1665}{\mega\watt}}$.
\end{enumerate}
\textbf{Conclusion :} À \SI{30}{\kilo\volt}, les pertes (\SI{1665}{\mega\watt}) \textbf{dépassent largement la puissance produite} (\SI{300}{\mega\watt}) : transport impossible. L'élévation de tension par transformateur (sortie alternateur \SI{15}{\kilo\volt} $\rightarrow$ \SI{225}{\kilo\volt} ou \SI{400}{\kilo\volt}) est \textbf{indispensable} pour rendre le transport rentable. C'est la raison historique de la victoire du courant alternatif (Tesla/Westinghouse) sur le courant continu (Edison) à la fin du XIX$^{\text{e}}$ siècle : le transformateur ne fonctionne qu'en alternatif.
\end{exoresolu}

\begin{aretenir}[Bilan global : de la centrale à la prise]
\begin{enumerate}
    \item \textbf{Production :} Alternateur synchrone ($U_{\text{gen}} \approx \SIrange{10}{20}{\kilo\volt}$, \SI{50}{\hertz}).
    \item \textbf{Élévation :} Transformateur abaisseur/élévateur ($U \uparrow$, $I \downarrow$) $\rightarrow$ \textbf{HT/THT} (\SI{90}{\kilo\volt}, \SI{225}{\kilo\volt}, \SI{400}{\kilo\volt} au Cameroun).
    \item \textbf{Transport :} Lignes HT (SONATREL), pertes Joule limitées ($\sim 3\text{--}5\%$).
    \item \textbf{Abaissement :} Postes source ($225/30\,\text{kV}$ ou $225/15\,\text{kV}$).
    \item \textbf{Distribution :} Réseau MT (\SIrange{15}{30}{\kilo\volt}) puis BT (\SI{230}{\volt}/\SI{400}{\volt}) vers l'usager (ENEO).
\end{enumerate}
L'alternateur est la \textbf{porte d'entrée} de cette chaîne : sa synchronisation rigoureuse ($f=\SI{50}{\hertz}$, $\varphi$ calé sur le réseau) est la condition \emph{sine qua non} de la stabilité du système électrique national.
\end{aretenir}

\begin{exercice}[Analyse d'une installation de production]
On considère une petite centrale hydroélectrique de type « au fil de l'eau » installée sur un affluent de la Sanaga. La turbine Francis entraîne directement un alternateur synchrone triphasé.
Données : Chute nette $H = \SI{25}{\meter}$ ; Débit turbiné $Q = \SI{15}{\meter^3\per\second}$ ; Rendement turbine $\eta_t = 0,90$ ; Rendement alternateur $\eta_a = 0,96$ ; Tension de sortie alternateur $U = \SI{15}{\kilo\volt}$ (composée) ; $\cos\varphi = 0,9$.
$\rho_{\text{eau}} = \SI{1000}{\kilogram\per\meter^3}$ ; $g = \SI{9,81}{\newton\per\kilogram}$.
\begin{enumerate}
    \item Calculer la puissance hydraulique disponible $P_{\text{hyd}}$.
    \item Calculer la puissance mécanique sur l'arbre $P_{\text{meca}}$.
    \item Calculer la puissance électrique active $P_{\text{elec}}$ produite.
    \item Calculer le courant de ligne $I$ délivré par l'alternateur.
    \item Si cette centrale doit injecter sur le réseau \SI{225}{\kilo\volt} via un transformateur idéal, quel sera le courant dans la ligne HT ?
\end{enumerate}
\end{exercice}

\begin{corrige}[Exercice : Analyse d'une installation de production]
\begin{enumerate}
    \item $P_{\text{hyd}} = \rho g Q H = 1000 \times 9,81 \times 15 \times 25 = 3\,678\,750\,\text{W} \approx \mathbf{\SI{3,68}{\mega\watt}}$.
    \item $P_{\text{meca}} = \eta_t \times P_{\text{hyd}} = 0,90 \times 3,68 = \mathbf{\SI{3,31}{\mega\watt}}$.
    \item $P_{\text{elec}} = \eta_a \times P_{\text{meca}} = 0,96 \times 3,31 = \mathbf{\SI{3,18}{\mega\watt}}$.
    \item $P_{\text{elec}} = \sqrt{3} U I \cos\varphi \Rightarrow I = \frac{P_{\text{elec}}}{\sqrt{3} U \cos\varphi} = \frac{3,18 \times 10^6}{\sqrt{3} \times 15 \times 10^3 \times 0,9} \approx \frac{3,18 \times 10^6}{23,38 \times 10^3} \approx \mathbf{\SI{136}{\ampere}}$.
    \item Transformateur idéal : $P_{\text{HT}} = P_{\text{elec}}$ (pertes négligées). $U_{\text{HT}} = \SI{225}{\kilo\volt}$.
    $I_{\text{HT}} = \frac{P_{\text{elec}}}{\sqrt{3} U_{\text{HT}} \cos\varphi} = \frac{3,18 \times 10^6}{\sqrt{3} \times 225 \times 10^3 \times 0,9} \approx \frac{3,18 \times 10^6}{350,7 \times 10^3} \approx \mathbf{\SI{9,07}{\ampere}}$.
    \textbf{Observation :} Le courant est divisé par $15$ (ratio des tensions), les pertes Joule divisées par $225$ ($15^2$).
\end{enumerate}
\end{corrige}

\begin{methode}[Réaliser un dispositif produisant un courant alternatif (SAVOIR-FAIRE)]
\textbf{Objectif :} Construire un alternateur rudimentaire (type dynamo vélo ou manivelle) et visualiser la tension alternative à l'oscilloscope.
\begin{enumerate}
    \item \textbf{Matériel :} Aimant néodyme puissant (disque ou bloc), fil émaillé (\SIrange{0,2}{0,3}{\milli\meter}), support rotatif (bouchon liège, bouchon plastique, rotor imprimé 3D), oscilloscope (ou carte son + logiciel Audacity/Visual Analyser), diodes LED (pour redresser).
    \item \textbf{Bobinage :} Enrouler $N \approx 500\text{--}1000$ spires sur un noyau (air ou fer doux - clous). Deux bobines décalées pour biphasé (optionnel).
    \item \textbf{Assemblage :} Fixer l'aimant sur la partie tournante, la bobine sur le stator fixe. Entrefer $< \SI{2}{\milli\meter}$.
    \item \textbf{Test :} Faire tourner manuellement (manivelle) ou par moteur DC faible vitesse. Observer $u(t)$ à l'oscilloscope : sinusoïde, $f = \frac{p \cdot n}{60}$ ($p=1$ paire de pôles ici). Mesurer $U_{\text{max}}$.
    \item \textbf{Variations :} Vitesse $\uparrow \Rightarrow f \uparrow, U_{\text{max}} \uparrow$. Nombre de spires $\uparrow \Rightarrow U_{\text{max}} \uparrow$. Aimant plus fort $\Rightarrow U_{\text{max}} \uparrow$.
    \item \textbf{Redressement :} Ajouter une diode (ou pont de 4 diodes) + condensateur $\rightarrow$ tension continue pour allumer une LED.
\end{enumerate}
\textbf{Compétence évaluée :} Identifier les parties (stator/rotor), expliquer le rôle de l'entrefer, relier $f$ à la vitesse de rotation, observer l'induction $e = -N \frac{\mathrm{d}\Phi}{\mathrm{d}t}$.
\end{methode}

\coursec{Réalisation d'un alternateur simple et synthèse}

Après avoir étudié les alternateurs industriels qui alimentent nos villes, depuis le barrage de \textbf{Lom Pangar} jusqu'aux groupes électrogènes des zones rurales non raccordées au réseau \textbf{ENEO}, il est temps de passer à la pratique. Comprendre la physique, c'est aussi savoir la mettre en œuvre. Dans cette dernière section, nous allons réaliser nous-mêmes un mini-alternateur, identifier ses composants, tracer son schéma de principe et synthétiser l'ensemble de la leçon. C'est l'aboutissement concret de la compétence « \emph{Réaliser un dispositif pouvant produire un courant alternatif} » attendue au Probatoire.

\courssub{Objectif et matériel nécessaire}

L'objectif est de \textbf{construire un dispositif expérimental simple mettant en évidence la production d'une tension alternative par induction électromagnétique}. Ce montage, bien que rudimentaire, reprend exactement le principe des grands alternateurs : la variation du flux magnétique à travers une bobine par mouvement relatif entre un aimant et un circuit conducteur.

\begin{experience}[Matériel pour un mini-alternateur de laboratoire]
\textbf{Liste du matériel (par binôme) :}
\begin{itemize}
    \item \textbf{Bobine :} fil de cuivre émaillé (diamètre \SI{0,3}{mm} à \SI{0,5}{mm}), enroulé sur un mandrin cylindrique (stylo, tube PVC) pour former \textbf{300 à 500 spires} serrées et régulières. Les extrémités sont dénudées sur \SI{2}{cm} (papier de verre fin).
    \item \textbf{Aimant :} aimant néodyme (NdFeB) puissant, de type disque (\SI{20}{mm} $\times$ \SI{5}{mm}) ou parallélépipédique. \emph{Attention : manipuler avec précaution (risque de pincement, éloigner des cartes magnétiques, téléphones, pacemakers).}
    \item \textbf{Axe de rotation :} un crayon en bois, une tige filetée \SI{6}{mm} ou un axe de moteur récupéré, traversant l'aimant (perçage central pour disque) ou collé solidement sur l'aimant (cyanoacrylate).
    \item \textbf{Manivelle :} une manivelle de récupération ou une pince crocodile vissée sur l'axe pour actionner la rotation manuellement.
    \item \textbf{Charge / Détecteur :} une \textbf{LED rouge ou jaune} (seuil \SI{\approx 1,8}{V}) \emph{en antiparallèle} (deux LED tête-bêche) pour visualiser l'alternance, \textbf{ou} un galvanomètre à aiguille (zéro central) sensible (\SI{\approx 1}{mA}), \textbf{ou} un multimètre en mode tension alternative (V~).
    \item \textbf{Support :} planchette en bois, étau ou pince pour fixer la bobine immobile face à l'aimant tournant.
\end{itemize}
\end{experience}

\begin{popfigure}
\begin{tikzpicture}[scale=0.9, every node/.style={font=\small}]
    % Base
    \fill[popcream] (-3.5,-1.5) rectangle (3.5,-1.3);
    % Bobine fixe (Stator)
    \draw[popblue, very thick] (-1.5,-0.5) -- (-1.5,1.5) -- (1.5,1.5) -- (1.5,-0.5);
    \foreach \y in {-0.3,0.1,0.5,0.9,1.3} {
        \draw[popblue!70!black, thick] (-1.7,\y) -- (-1.3,\y);
        \draw[popblue!70!black, thick] (1.3,\y) -- (1.7,\y);
    }
    \node[popblue, align=center] at (0,2.2) {\textbf{Bobine (Stator)}\\(300--500 spires)};
    
    % Fils de sortie
    \draw[popdark, thick] (-1.5,-0.5) -- (-2.5,-1.0) node[left] {Fil 1};
    \draw[popdark, thick] (1.5,-0.5) -- (2.5,-1.0) node[right] {Fil 2};
    
    % LED antiparallèle
    \begin{scope}[shift={(-2.5,-1.5)}]
        \draw[poporange, thick] (0,0) -- (0.5,0);
        \draw[poporange!80!black, fill=poporange!20] (0.5,-0.3) rectangle (1.0,0.3);
        \node at (0.75,0) {LED};
        \draw[poporange, thick] (1.0,0) -- (1.5,0);
        % Symbole LED antiparallèle simplifié
        \draw[poporange, thick] (0.75,-0.5) -- (0.75,-0.8);
        \draw[poporange, thick] (0.75,-0.5) -- (0.5,-0.5) -- (0.5,-0.8);
        \draw[poporange, thick] (0.75,-0.5) -- (1.0,-0.5) -- (1.0,-0.8);
        \node[poporange, align=center] at (0.75,-1.1) {2 LED\\tête-bêche};
    \end{scope}
    
    % Aimant tournant (Rotor)
    \draw[poppurple!80!black, thick, fill=poppurple!10] (-1.2,-0.6) rectangle (1.2,0.6);
    \node[poppurple!80!black] at (0,0) {\textbf{N} $\uparrow$ \textbf{S}};
    % Axe
    \draw[popmuted, thick] (0,0.6) -- (0,1.8);
    \draw[popmuted, thick] (0,-0.6) -- (0,-1.2);
    \node[poppurple, align=center] at (0,-2.0) {\textbf{Aimant (Rotor)}\\(Néodyme)};
    
    % Manivelle
    \draw[popdark, thick] (0,1.8) -- (-0.8,2.6);
    \draw[popdark, thick] (-0.8,2.6) -- (-0.8,3.0);
    \draw[popdark, thick] (-0.8,2.6) -- (-1.3,2.6);
    \node[popdark] at (-1.5,2.8) {Manivelle};
    
    % Flèches de rotation
    \draw[popgreen, thick, ->] (2.0,1.0) arc (0:180:2.0 and 1.0);
    \node[popgreen] at (2.5,1.5) {Rotation $\omega$};
    
    % Champ B
    \draw[poppurple, ->, thick] (0,0.6) -- (0,1.3) node[above] {$\vec{B}$};
    
    % Cotation distance
    \draw[<->, popmuted] (1.5,0.5) -- (1.2,0.5);
    \node[popmuted] at (1.5,0.5) {entrefer $e$};
\end{tikzpicture}
\legende{Dispositif expérimental : mini-alternateur à aimant tournant (rotor) face à une bobine fixe (stator). La LED antiparallèle clignote au rythme de la rotation.}
\end{popfigure}

\courssub{Protocole de réalisation : méthode en quatre étapes}

Voici la démarche rigoureuse pour mener à bien cette réalisation, exactement comme on l'attend dans un compte-rendu de TP au Probatoire.

\begin{methode}[Réaliser un alternateur simple en 4 étapes]
\textbf{Étape 1 : Construction du stator (partie fixe).}
Enroule soigneusement le fil émaillé sur le mandrin (\SI{300}{--}\SI{500}{ spires}). Plus le nombre de spires $N$ est grand, plus la f.é.m. induite $e = -N \frac{d\Phi}{dt}$ sera intense. Dénude proprement les deux extrémités sur \SI{2}{cm}. Fixe la bobine verticalement sur le support (planche + étau), face libre vers l'avant.

\textbf{Étape 2 : Construction du rotor (partie tournante).}
Fixe l'aimant néodyme sur l'axe de rotation (perçage central ou collage robuste). Vérifie que l'aimant tourne librement sans toucher la bobine. L'entrefer $e$ (distance aimant/bobine) doit être minimal (\SI{\approx 2}{--}\SI{5}{mm}) pour maximiser le flux $\Phi = \iint \vec{B} \cdot d\vec{S}$ sans frottement.

\textbf{Étape 3 : Raccordement de la charge.}
Connecte les deux fils de la bobine à la charge :
\begin{itemize}
    \item \textbf{LED antiparallèle :} soude deux LED en parallèle inverse (anode 1 sur cathode 2, cathode 1 sur anode 2). Cela permet de visualiser les deux alternances.
    \item \textbf{Galvanomètre :} branche en série (zéro central).
    \item \textbf{Multimètre :} mode V~ (mesure efficace $U_{\text{eff}}$).
\end{itemize}

\textbf{Étape 4 : Essai et observations.}
Actionne la manivelle pour faire tourner l'aimant.
\begin{enumerate}
    \item \textbf{Vitesse lente :} la LED clignote distinctement (rouge/vert ou allumée/éteinte selon montage), l'aiguille oscille de part et d'autre du zéro.
    \item \textbf{Vitesse rapide :} la LED semble allumée en continu (persistance rétinienne), l'aiguille vibre autour de zéro, le multimètre affiche une tension efficace croissante.
\end{enumerate}
\end{methode}

\begin{attention}[Pièges classiques à éviter]
\begin{itemize}
    \item \textbf{Bobine à trop peu de spires ($N < 100$) :} la f.é.m. de pointe $E_{\text{max}} = N B S \omega$ reste inférieure au seuil de la LED (\SI{\approx 1,8}{V}). \emph{Remède :} viser $N \ge 300$.
    \item \textbf{Aimant trop faible (ferrite de haut-parleur) :} $B$ trop faible ($\sim \SI{0,1}{T}$ vs $\sim \SI{1,2}{T}$ pour le néodyme). \emph{Remède :} utiliser impérativement un aimant néodyme (terres rares).
    \item \textbf{Entrefer trop grand :} le flux $\Phi$ chute comme $1/e^3$ (dipôle). \emph{Remède :} rapprocher au maximum sans contact mécanique.
    \item \textbf{Fil non dénudé :} l'émail isole parfaitement. \emph{Remède :} vérifier la continuité à l'ohmmètre avant montage.
    \item \textbf{LED unique (non antiparallèle) :} elle ne conduit qu'une alternance sur deux, clignote à $f/2$ et risque d'être détruite en inverse (tension de claquage \SI{\approx 5}{V}). \emph{Remède :} toujours deux LED tête-bêche ou une diode de roue libre.
\end{itemize}
\end{attention}

\courssub{Observations attendues et interprétation physique}

Lorsque l'aimant tourne devant la bobine, le flux magnétique $\Phi(t)$ traversant la bobine varie sinusoïdalement (approximation dipolaire) :
\[
\Phi(t) = B S N \cos(\omega t) \quad \text{(choix d'origine du temps)}
\]
La loi de Faraday-Lenz donne la f.é.m. induite :
\[
e(t) = -N \frac{d\Phi}{dt} = N B S \omega \sin(\omega t) = E_{\text{max}} \sin(\omega t)
\]
avec $E_{\text{max}} = N B S \omega$.

\begin{exemplebox}[Estimation d'ordre de grandeur de la tension obtenue]
Prenons des valeurs réalistes pour notre montage :
\begin{itemize}
    \item $N = 400$ spires
    \item Aimant néodyme : $B \approx \SI{1,0}{T}$ au niveau de la bobine
    \item Section de la bobine : diamètre \SI{2}{cm} $\Rightarrow S = \pi (0,01)^2 \approx \SI{3,14e-4}{m^2}$
    \item Vitesse de rotation manuelle : $f = 3$ tours/s $\Rightarrow \omega = 2\pi f \approx \SI{18,8}{rad/s}$
\end{itemize}
Calculons l'amplitude :
\[
E_{\text{max}} = 400 \times 1,0 \times 3,14 \times 10^{-4} \times 18,8 \approx \SI{2,36}{V}
\]
La tension efficace est $U_{\text{eff}} = \frac{E_{\text{max}}}{\sqrt{2}} \approx \SI{1,67}{V}$.
\textbf{Conclusion :} Cette valeur est juste au seuil d'une LED rouge (\SI{1,8}{V}). Cela explique pourquoi il faut tourner \emph{vite} ($f > 3$ Hz) ou augmenter $N$ pour voir une lumière franche. C'est un excellent exercice d'analyse dimensionnelle : $[N B S \omega] = 1 \cdot T \cdot m^2 \cdot s^{-1} = \frac{N}{A\cdot m} \cdot m^2 \cdot s^{-1} = \frac{N\cdot m}{A\cdot s} = \frac{J}{C} = V$.
\end{exemplebox}

\begin{popfigure}
\begin{tikzpicture}
\begin{axis}[
    width=\linewidth, height=6cm,
    axis lines=middle,
    xlabel={$t$ (s)}, ylabel={$e(t)$ (V)},
    domain=0:1.33, samples=200,
    xtick={0,0.33,0.66,1.0,1.33},
    xticklabels={$0$, $T/4$, $T/2$, $3T/4$, $T$},
    ytick={-2.36,0,2.36},
    yticklabels={$-E_{\text{max}}$, $0$, $+E_{\text{max}}$},
    ymin=-2.8, ymax=2.8,
    xmin=0, xmax=1.35,
    grid=major, grid style={popline},
    popcurve
]
\addplot {2.36*sin(deg(2*pi*3*x))};
\addplot[poporange, dashed, thick, domain=0:1.33] {1.8};
\addplot[poporange, dashed, thick, domain=0:1.33] {-1.8};
\node[poporange, anchor=south west] at (axis cs:0.1,1.8) {Seuil LED ($\approx 1,8$ V)};
\end{axis}
\end{tikzpicture}
\legende{Forme d'onde de la f.é.m. induite $e(t)$ pour $f=3$ Hz. La tension dépasse le seuil de la LED seulement près des crêtes.}
\end{popfigure}

\courssub{Variante : le mouvement relatif}

Une question fondamentale de l'induction électromagnétique : \emph{qu'est-ce qui tourne vraiment ?}

\begin{propriete}[Invariance par mouvement relatif]
\textbf{L'effet d'induction ne dépend que du mouvement relatif entre l'aimant et la bobine.}
Faire tourner l'aimant devant une bobine fixe \textbf{ou} faire tourner la bobine dans le champ d'un aimant fixe produit \textbf{exactement la même f.é.m. alternative} (même amplitude, même fréquence, même forme d'onde), pourvu que la vitesse angulaire relative $\omega$ soit identique.
\end{propriete}

\begin{experience}[Variante : Bobine tournante (Rotor) / Aimant fixe (Stator)]
\textbf{Montage :} Fixe l'aimant sur le support. Monte la bobine sur l'axe de rotation (rotor). Le raccordement de la bobine tournante vers la charge fixe (LED) nécessite des \textbf{anneaux collecteurs (bagues) et balais} (ou fils torsadés acceptant quelques tours).
\textbf{Observation :} Même comportement : LED clignotante, tension alternative.
\textbf{Intérêt pédagogique :} Cette variante montre que dans un \textbf{alternateur synchrone industriel}, c'est souvent l'induit (bobines) qui est fixe (stator) et l'inducteur (aimants/électroaimants) qui tourne (rotor), \emph{mais} l'inverse existe (alternateurs à rotor bobiné, stator à aimants permanents dans certaines éoliennes).
\end{experience}

\courssub{Schéma de principe du mini-alternateur réalisé}

C'est un savoir-faire exigible : « \emph{Faire le schéma de principe d'un alternateur} ». Le schéma de principe ne représente pas l'aspect réaliste (manivelle, support) mais la structure fonctionnelle : stator, rotor, entrefer, connexions.

\begin{popfigure}
\begin{tikzpicture}[scale=1.1, every node/.style={font=\small}]
    % Stator (Bobine) - Représenté en coupe
    \draw[popblue, very thick, fill=popblue!5] (-2.5,-1.5) rectangle (-0.5,1.5);
    \draw[popblue, very thick, fill=popblue!5] (0.5,-1.5) rectangle (2.5,1.5);
    % Spires symboliques
    \foreach \y in {-1.0,-0.5,0,0.5,1.0} {
        \draw[popblue!70!black] (-2.5,\y) -- (-0.5,\y);
        \draw[popblue!70!black] (0.5,\y) -- (2.5,\y);
    }
    \node[popblue, align=center] at (-1.5,2.0) {\textbf{Stator}\\(Bobine $N$ spires)};
    
    % Entrefer
    \draw[<->, popmuted] (-0.5,0) -- (0.5,0);
    \node[popmuted] at (0,0.4) {Entrefer $e$};
    
    % Rotor (Aimant) - Disque
    \draw[poppurple!80!black, thick, fill=poppurple!10] (0,-0.8) circle (1.2);
    \node[poppurple!80!black] at (0,-0.8) {\textbf{N} $\uparrow$ \textbf{S}};
    \node[poppurple, align=center] at (0,-2.3) {\textbf{Rotor}\\(Aimant permanent)};
    
    % Axe
    \draw[popmuted, thick] (0,0.4) -- (0,1.5);
    \draw[popmuted, thick] (0,-2.0) -- (0,-2.5);
    
    % Flux B
    \draw[poppurple, ->, thick] (0,-0.8) -- (0,0.2);
    \draw[poppurple, ->, thick] (0,-0.8) -- (-1.0,0.5);
    \draw[poppurple, ->, thick] (0,-0.8) -- (1.0,0.5);
    \node[poppurple] at (-1.5,0.8) {$\vec{B}$};
    
    % Sorties
    \draw[popdark, thick] (-2.5,-1.5) -- (-3.5,-1.5) node[left] {Bornes de sortie};
    \draw[popdark, thick] (2.5,-1.5) -- (3.5,-1.5);
    
    % Charge (LED antiparallèle)
    \begin{scope}[shift={(-3.5,-2.0)}]
        \draw[poporange, thick] (0,0) -- (1,0);
        \draw[poporange!80!black, fill=poporange!20] (1,-0.25) rectangle (1.5,0.25);
        \node at (1.25,0) {LED};
        \draw[poporange, thick] (1.5,0) -- (2,0);
        \draw[poporange, thick] (1.25,-0.5) -- (1.25,-0.8);
        \draw[poporange, thick] (1.25,-0.5) -- (1.0,-0.5) -- (1.0,-0.8);
        \draw[poporange, thick] (1.25,-0.5) -- (1.5,-0.5) -- (1.5,-0.8);
        \node[poporange, align=center] at (1.25,-1.0) {Charge};
    \end{scope}
    
    % Légende rotation
    \draw[popgreen, thick, ->] (2.5,1.0) arc (0:180:1.5 and 1.0);
    \node[popgreen] at (3.5,1.5) {Rotation $\vec{\omega}$};
\end{tikzpicture}
\legende{Schéma de principe du mini-alternateur à aimant tournant (rotor) et bobine fixe (stator). À reproduire soigneusement : stator (bobines), rotor (aimant N/S), entrefer, flux $\vec{B}$, sens de rotation $\vec{\omega}$, bornes de sortie.}
\end{popfigure}

\courssub{Synthèse de la leçon : L'alternateur}

Nous avons maintenant tous les éléments pour dresser le bilan complet. Cette synthèse constitue ta fiche de révision « tout-en-un » pour le Probatoire.

\begin{aretenir}[Synthèse : L'alternateur]
\textbf{1. Définition et Rôle}
Un \textbf{alternateur} est une machine électromécanique qui \textbf{transforme de l'énergie mécanique de rotation en énergie électrique alternative} (tension et courant sinusoïdaux). C'est la source quasi unique de l'électricité industrielle (barrages, centrales thermiques, nucléaires, éoliennes, groupes électrogènes).

\textbf{2. Principe de fonctionnement : Induction électromagnétique}
\begin{itemize}
    \item Un aimant (ou électroaimant) tourne $\rightarrow$ variation du flux magnétique $\Phi(t)$ à travers une bobine.
    \item Loi de Faraday-Lenz : $e(t) = -N \frac{d\Phi}{dt}$.
    \item Pour une rotation uniforme $\omega$ dans un champ radial uniforme : $e(t) = E_{\text{max}} \sin(\omega t)$ avec $E_{\text{max}} = N B S \omega$.
    \item Fréquence $f = \frac{\omega}{2\pi} = \frac{p \cdot n}{60}$ ($p$ paires de pôles, $n$ tr/min).
\end{itemize}

\textbf{3. Parties essentielles et rôles}
\begin{center}
\begin{tabularx}{\linewidth}{|l|X|X|}\hline
\textbf{Partie} & \textbf{Nom technique} & \textbf{Rôle} \\ \hline
\textbf{Rotor} & Inducteur (tournant) & Crée le champ magnétique tournant $\vec{B}$ (aimants permanents ou électroaimant alimenté en courant continu via bagues/balais). \\ \hline
\textbf{Stator} & Induit (fixe) & Reçoit le flux variable ; bobinages (en général 3 phases décalées de $120^\circ$) où s'induit la f.é.m. alternative $e(t)$. \\ \hline
\textbf{Entrefer} & Zone d'air & Couplage magnétique rotor/stator ; doit être mince et régulier. \\ \hline
\textbf{Bagues/balais} & (Si rotor bobiné) & Alimentent l'électroaimant du rotor en courant continu (excitation) sans transmettre la puissance principale. \\ \hline
\end{tabularx}
\end{center}

\textbf{4. Chaîne de transformations d'énergie}
\[
\boxed{\text{\'Energie primaire}} \xrightarrow{\text{Turbine/Moteur}} \boxed{\text{\'Energie mécanique (rotation)}} \xrightarrow{\text{Alternateur (Induction)}} \boxed{\text{\'Energie électrique alternative}}
\]
\begin{itemize}
    \item \textbf{Barrage (Lom Pangar, Nachtigal, Edéa) :} Pesanteur $\rightarrow$ Cinétique eau $\rightarrow$ Mécanique turbine $\rightarrow$ Électrique alternateur.
    \item \textbf{Centrale thermique/Gaz :} Chimique (combustion) $\rightarrow$ Thermique vapeur $\rightarrow$ Mécanique turbine $\rightarrow$ Électrique.
    \item \textbf{Éolienne :} Cinétique vent $\rightarrow$ Mécanique pale/rotor $\rightarrow$ Électrique (souvent via multiplicateur ou direct-drive).
    \item \textbf{Groupe électrogène (quartier/hôpital) :} Chimique carburant $\rightarrow$ Mécanique moteur Diesel $\rightarrow$ Électrique alternateur.
\end{itemize}

\textbf{5. Exemples et paramètres clés}
\begin{itemize}
    \item \textbf{Alternateur monophasé :} 1 bobine (ou plusieurs en série) $\rightarrow$ 1 tension $e(t)$. Usage : groupes portables, soudure, petite éolienne.
    \item \textbf{Alternateur triphasé (standard industriel) :} 3 jeux de bobines décalés de $120^\circ$ (stator) $\rightarrow$ 3 tensions $e_1, e_2, e_3$. Transport (SONATREL/ENEO) : 90 kV, 225 kV. Distribution : 230 V / 400 V, $f = \SI{50}{Hz}$.
    \item \textbf{Alternateur à aimants permanents :} Rotor = aimants néodyme. Compact, rendement élevé, pas d'excitation. Usage : éoliennes, véhicules hybrides, micro-alternateurs.
    \item \textbf{Alternateur à rotor bobiné (synchrone) :} Rotor = électroaimant. Régulation de tension par courant d'excitation $I_{\text{exc}}$. Usage : toutes grosses centrales, groupes électrogènes industriels.
\end{itemize}
\end{aretenir}

\begin{savaistu}[Histoire et contexte camerounais]
\textbf{1866 :} Werner von Siemens construit la première \textbf{machine dynamo-électrique industrielle} à induction (alternateur à rotor bobiné). C'est le point de départ de l'électrification mondiale.
\textbf{Années 1890 :} Guerre des courants (Edison DC vs Tesla/Westinghouse AC). L'alternateur triphasé de Tesla (1888) l'emporte grâce au transformateur (haute tension = faibles pertes).
\textbf{Au Cameroun :} La première centrale hydroélectrique, \textbf{Edéa (Sanaga)}, mise en service en \textbf{1947} (alors 12 MW, aujourd'hui \SI{\approx 263}{MW}), utilise des alternateurs synchrones entraînés par turbines Francis. \textbf{Song Loulou} (1988, \SI{384}{MW}), \textbf{Lom Pangar} (2017, \SI{30}{MW} + régulation), \textbf{Nachtigal} (2023/24, \SI{420}{MW}) sont les piliers actuels. Chaque MW injecté sur le réseau \textbf{SONATREL} sort des enroulements d'un stator d'alternateur tournant à \SI{1500}{tr/min} (pour $f=50$ Hz, $p=1$) ou \SI{1000}{tr/min} ($p=2$), etc. L'électricité qui charge ton téléphone MTN/Orange ou alimente l'éclairage du stade d'Olembé a parcouru ce chemin : \emph{Eau $\rightarrow$ Turbine $\rightarrow$ Rotor $\rightarrow$ Flux $\rightarrow$ Stator $\rightarrow$ Réseau $\rightarrow$ Toi}.
\end{savaistu}

\courssub{Exercice résolu : Caractéristiques d'un mini-alternateur de TP}

\begin{exoresolu}[Analyse quantitative du montage réalisé]
\textbf{Énoncé :} Un élève réalise le mini-alternateur de la section 6.1. La bobine a $N = 450$ spires, section $S = \SI{3,5e-4}{m^2}$. L'aimant néodyme crée un champ radial $B = \SI{1,1}{T}$ dans l'entrefer. L'élève tourne la manivelle à une fréquence de rotation $f = \SI{4,0}{Hz}$.
\begin{enumerate}
    \item Calculer la pulsation $\omega$ et la période $T$ de la f.é.m. induite.
    \item Calculer l'amplitude $E_{\text{max}}$ et la valeur efficace $U_{\text{eff}}$ de la tension aux bornes de la bobine (circuit ouvert).
    \item La charge est une LED rouge (seuil $V_{\text{seuil}} = \SI{1,85}{V}$) montée en antiparallèle. La résistance interne de la bobine est $r = \SI{15}{\ohm}$. La LED conduit-elle en permanence ? Justifier.
    \item Si l'élève double la vitesse de rotation ($f' = \SI{8,0}{Hz}$), que deviennent $E_{\text{max}}$ et $U_{\text{eff}}$ ? La LED brille-t-elle plus fort ?
\end{enumerate}

\tcblower
\textbf{Solution détaillée :}

\textbf{1. Pulsation et période}
\[
\omega = 2\pi f = 2\pi \times 4,0 = \SI{8\pi}{\radian/\second} \approx \SI{25,1}{\radian/\second}
\]
\[
T = \frac{1}{f} = \frac{1}{4,0} = \SI{0,25}{s}
\]

\textbf{2. Tensions à vide}
Amplitude :
\[
E_{\text{max}} = N B S \omega = 450 \times 1,1 \times 3,5 \times 10^{-4} \times 25,1
\]
\[
E_{\text{max}} \approx \SI{4,35}{V}
\]
Valeur efficace (rms) :
\[
U_{\text{eff}} = \frac{E_{\text{max}}}{\sqrt{2}} \approx \frac{4,35}{1,414} \approx \SI{3,08}{V}
\]

\textbf{3. Comportement avec la LED (circuit fermé)}
La f.é.m. est $e(t) = 4,35 \sin(25,1 t)$.
La LED conduit seulement si $|e(t)| > V_{\text{seuil}} = \SI{1,85}{V}$.
Résolution : $4,35 |\sin(\omega t)| > 1,85 \Rightarrow |\sin(\omega t)| > 0,425$.
Cela correspond à des angles $\omega t \in [25^\circ, 155^\circ] \cup [205^\circ, 335^\circ]$ modulo $360^\circ$.
La LED conduit pendant $\frac{130^\circ \times 2}{360^\circ} \approx 72\%$ de la période.
\textbf{Non, elle ne conduit pas en permanence.} Elle clignote à la fréquence $f = \SI{4}{Hz}$ (allumée $\approx \SI{180}{ms}$, éteinte $\approx \SI{70}{ms}$ par alternance). L'œil perçoit un scintillement.
Courant de crête dans la LED (approx, négligeant $r$ pour l'ordre de grandeur) : $I_{\text{crête}} \approx \frac{E_{\text{max}} - V_{\text{seuil}}}{r} = \frac{4,35 - 1,85}{15} \approx \SI{0,167}{A} = \SI{167}{mA}$. \emph{Attention :} cela dépasse souvent le courant nominal d'une LED standard (\SI{20}{mA}). Une résistance de limitation série serait nécessaire en vrai TP, mais ici la résistance interne $r$ limite un peu.

\textbf{4. Vitesse doublée ($f' = \SI{8}{Hz}$)}
$\omega' = 2\omega$.
$E'_{\text{max}} = N B S \omega' = 2 E_{\text{max}} \approx \SI{8,70}{V}$.
$U'_{\text{eff}} = 2 U_{\text{eff}} \approx \SI{6,16}{V}$.
La LED conduit pour $|e'(t)| > 1,85 \Rightarrow 8,70 |\sin| > 1,85 \Rightarrow |\sin| > 0,21$.
Angle de conduction $\approx 78^\circ \times 2 / 360^\circ \approx 43\%$ par alternance ? Non, $\arcsin(0,21) \approx 12^\circ$. Conduction de $12^\circ$ à $168^\circ$ ($156^\circ$) $\times 2 = 312^\circ / 360^\circ \approx \mathbf{87\%}$ du temps.
\textbf{Oui, elle brille beaucoup plus fort} (courant moyen plus élevé, crête $\approx \frac{8,7-1,85}{15} \approx \SI{450}{mA}$ -- danger pour la LED sans résistance série !). La luminosité perçue augmente fortement (fréquence \SI{8}{Hz} $\rightarrow$ clignotement moins visible, courant moyen $\times 2$ environ).
\end{exoresolu}

\courssub{Exercices d'entraînement (Probatoire)}

\begin{exercice}[Exercice 1 : Identification et schéma]
On présente à un élève un alternateur monophasé de groupe électrogène (rotor à aimants permanents, stator bobiné).
\begin{enumerate}
    \item Nomme les parties 1 (tournant), 2 (fixe), 3 (espace entre 1 et 2), 4 (fils de sortie).
    \item Dessine le schéma de principe de cet alternateur en indiquant le sens de rotation $\vec{\omega}$, le champ $\vec{B}$ dans l'entrefer, et la polarité des bornes de sortie à l'instant où le pôle Nord passe devant le côté gauche de la bobine (utiliser la règle du flux ou Lenz).
    \item Quelle est la nature du courant délivré ? Quelle est sa fréquence si le rotor tourne à \SI{3000}{tr/min} et possède 2 pôles ?
\end{enumerate}
\end{exercice}

\begin{exercice}[Exercice 2 : Chaîne énergétique et rendement]
Un groupe électrogène de secours d'un hôpital de district (alimentation bloc opératoire) est constitué d'un moteur Diesel entraînant un alternateur triphasé \SI{400}{V} / \SI{50}{Hz}.
\begin{enumerate}
    \item Complète la chaîne d'énergie : \_\_\_\_\_\_ (carburant) $\rightarrow$ \_\_\_\_\_\_ (moteur) $\rightarrow$ \_\_\_\_\_\_ (arbre) $\rightarrow$ \_\_\_\_\_\_ (alternateur) $\rightarrow$ \_\_\_\_\_\_ (réseau).
    \item La puissance électrique active fournie est $P_e = \SI{50}{kW}$. Le rendement global (moteur + alternateur) est $\eta = 0,35$. Calculer la puissance thermique fournie par le carburant (pouvoir calorifique massique $PCI = \SI{43}{MJ/kg}$) et le débit massique de carburant en $\SI{}{kg/h}$.
    \item Pourquoi l'alternateur est-il triphasé et non monophasé pour cette puissance ? (2 arguments techniques).
\end{enumerate}
\end{exercice}

\begin{corrige}[Exercice 1]
\begin{enumerate}
    \item 1 = Rotor (Inducteur) ; 2 = Stator (Induit) ; 3 = Entrefer ; 4 = Bornes de sortie (phasé + neutre).
    \item Schéma : Cercle rotor (N/S), rectangle stator (bobine), flèche $\vec{\omega}$, flèches $\vec{B}$ radial entrant/sortant. À l'instant où N passe à gauche, le flux $\Phi$ vers la droite est max et décroît $\rightarrow$ courant induit s'oppose à la décroissance $\rightarrow$ champ créé par bobine vers la droite $\rightarrow$ courant sortant par le haut (pôle Nord créé à gauche). Borne du haut = $+$ (convention générateur).
    \item Courant alternatif sinusoïdal. $f = \frac{p n}{60} = \frac{1 \times 3000}{60} = \SI{50}{Hz}$.
\end{enumerate}
\end{corrige}

\begin{corrige}[Exercice 2]
\begin{enumerate}
    \item Chimique $\rightarrow$ Thermique (combustion) $\rightarrow$ Mécanique (vilebrequin) $\rightarrow$ Électrique (alternateur) $\rightarrow$ Électrique (réseau triphasé).
    \item $P_{\text{thermique}} = \frac{P_e}{\eta} = \frac{50}{0,35} \approx \SI{142,9}{kW} = \SI{142,9}{kJ/s}$.
    $\dot{m} = \frac{P_{\text{thermique}}}{PCI} = \frac{142,9 \times 10^3}{43 \times 10^6} \approx \SI{3,32e-3}{kg/s} = \SI{11,96}{kg/h} \approx \SI{12}{kg/h}$.
    \item 1) Pour une même puissance, le courant par phase est divisé par $\sqrt{3}$ par rapport au monophasé $\rightarrow$ sections de câbles réduites, pertes Joule divisées par 3. 2) Production d'un champ tournant parfait pour moteurs asynchrones (pompes, ventilateurs, compresseurs médicaux) sans dispositif de démarrage complexe. 3) Puissance instantanée constante (pas de pulsation à $2f$), couple moteur régulier, moins de vibrations mécaniques sur l'arbre.
\end{enumerate}
\end{corrige}

\begin{aretenir}[Check-list finale pour le Probatoire - Leçon 17]
\begin{itemize}
    \item[\checkmark] \textbf{Définir} l'alternateur : machine méca $\rightarrow$ élec alternative.
    \item[\checkmark] \textbf{Expliquer} le principe : Induction (Faraday-Lenz), variation flux $\Phi$ par rotation.
    \item[\checkmark] \textbf{Nommer et donner le rôle} : Rotor (inducteur, champ B), Stator (induit, f.é.m.), Entrefer, Bagues/Balais (si excitation).
    \item[\checkmark] \textbf{Tracer} le schéma de principe (stator, rotor, entrefer, $\vec{B}$, $\vec{\omega}$, bornes).
    \item[\checkmark] \textbf{Écrire} la chaîne d'énergie complète selon la source (barrage, thermique, éolien, groupe).
    \item[\checkmark] \textbf{Calculer} $f = p n / 60$, $E_{\text{max}} = N B S \omega$, $U_{\text{eff}}$.
    \item[\checkmark] \textbf{Réaliser/Décrire} le TP mini-alternateur (bobine, aimant, rotation, LED antiparallèle, observations).
    \item[\checkmark] \textbf{Savoir} que le mouvement relatif seul compte (rotor/stator interchangeable pour le principe).
\end{itemize}
\end{aretenir}

\textbf{Félicitations !} Tu viens de boucler le Module 4 « Aspects énergétiques des circuits électriques ». Tu maîtrises maintenant la physique qui se cache derrière chaque prise de courant, chaque lampe allumée à Yaoundé ou à Douala, chaque moteur qui tourne. L'alternateur est le cœur battant de notre monde électrique. Révise cette fiche de synthèse, refais les calculs de l'exercice résolu, et tu seras prêt pour l'épreuve. \emph{Courage et rigueur !}

\newpage

\coursec{Activité d'intégration}

\coursec{L'alternateur}

\courssub{Objectifs de l'activité}

À la fin de cette activité, tu seras capable de :
\begin{itemize}
\item mobiliser le \cle{rôle} de l'alternateur (conversion d'énergie mécanique en énergie électrique) dans une situation concrète de production d'électricité ;
\item identifier les \cle{parties essentielles} d'un alternateur (\cle{rotor} et \cle{stator}) et donner le rôle de chacune ;
\item faire le \cle{schéma de principe} d'un alternateur ;
\item établir la \cle{chaîne complète des transformations d'énergie} depuis le carburant jusqu'à l'énergie électrique ;
\item exploiter le \cle{rendement} d'une conversion pour dimensionner une installation ;
\item justifier, par un raisonnement argumenté, le \cle{choix d'un générateur électromécanique} par rapport à d'autres solutions.
\end{itemize}

\courssub{Situation}

\begin{experience}[Le délestage du quartier Plateau à Garoua]
Il est 19~h~15 à Garoua, chef-lieu de la région du Nord. La chaleur est encore écrasante (34~\textdegree C) et, comme souvent en saison sèche, la société de distribution d'électricité annonce un \emph{délestage} : le quartier Plateau sera privé de courant de 19~h à 23~h, soit quatre heures consécutives.

Dans ce quartier, la coopérative des habitants \og Wuro Sallé \fg{} décide d'acheter un \cle{groupe électrogène} pour alimenter la salle commune (études des enfants, veillées, conservation des vaccins du poste de santé de proximité). Le responsable technique, M.~Hamadou, recense les appareils à faire fonctionner \emph{simultanément} :
\begin{itemize}
\item 20 lampes de puissance $P_1 = \SI{60}{W}$ chacune ;
\item 5 ventilateurs de puissance $P_2 = \SI{75}{W}$ chacun ;
\item 2 réfrigérateurs de puissance $P_3 = \SI{150}{W}$ chacun.
\end{itemize}

Le vendeur de matériel électrique de Douala propose un groupe électrogène constitué d'un \emph{moteur thermique} (fonctionnant au gasoil) accouplé à un \emph{alternateur} dont le rendement est $\eta = \SI{90}{\%} = \num{0,90}$. Il affirme aussi que \og des batteries suffiraient largement \fg{}.

M.~Hamadou, qui se souvient de ses cours de physique de Première, hésite. Il sollicite ta classe pour l'aider à \textbf{dimensionner le groupe électrogène}, à \textbf{comprendre le fonctionnement de son alternateur} et à \textbf{trancher le débat} groupe électrogène contre batteries.
\end{experience}

\courssub{Consignes}

Tu répondras aux questions suivantes, en rédigeant soigneusement tes raisonnements et en justifiant chaque résultat numérique par un calcul littéral puis une application numérique avec unité.

\begin{enumerate}
\item Calcule la \cle{puissance électrique totale} $P_{\text{él}}$ nécessaire pour alimenter simultanément tous les appareils du quartier.
\item En tenant compte du rendement $\eta = \num{0,90}$ de l'alternateur, détermine la \cle{puissance mécanique} $P_{\text{méc}}$ que le moteur thermique doit fournir à l'arbre de l'alternateur. Exprime le résultat en watts puis en kilowatts.
\item Trace le \cle{schéma de principe} de l'alternateur du groupe électrogène. Tu y feras apparaître le \cle{rotor} (partie tournante) et le \cle{stator} (partie fixe), et tu indiqueras pour chacun son rôle ainsi que le sens de la conversion d'énergie.
\item Établis la \cle{chaîne complète des transformations d'énergie} dans le groupe électrogène, depuis l'énergie chimique du gasoil jusqu'à l'énergie électrique consommée par les lampes. Tu préciseras à chaque étape l'organe qui réalise la conversion.
\item M.~Hamadou hésite entre le groupe électrogène et un parc de batteries. En t'appuyant sur le rôle de l'alternateur (production \emph{continue} d'énergie électrique tant qu'on lui fournit de l'énergie mécanique) et sur la durée du délestage (4~h, renouvelée chaque jour), rédige un court paragraphe (5 à 8 lignes) justifiant le choix du groupe électrogène.
\end{enumerate}

\courssub{Ressources à mobiliser}

\begin{aretenir}[Boîte à outils de l'activité]
\begin{itemize}
\item \textbf{Puissances en parallèle :} des appareils branchés simultanément sur la même installation consomment des puissances qui s'ajoutent : $P_{\text{totale}} = \sum_i P_i$.
\item \textbf{Rendement de l'alternateur :} $\eta = \dfrac{P_{\text{électrique fournie}}}{P_{\text{mécanique reçue}}}$, donc $P_{\text{méc}} = \dfrac{P_{\text{él}}}{\eta}$. Le rendement est un nombre sans dimension, toujours inférieur à 1.
\item \textbf{Rôle de l'alternateur :} convertir de l'\cle{énergie mécanique} (de rotation) en \cle{énergie électrique} sous forme de tension et de courant \emph{alternatifs}.
\item \textbf{Parties essentielles :} le \cle{rotor} (inducteur, généralement un aimant ou un électroaimant tournant) crée un champ magnétique en rotation ; le \cle{stator} (induit, bobines fixes) est le siège de la tension induite. C'est le mouvement relatif rotor/stator qui produit la tension alternative.
\item \textbf{Chaîne énergétique du groupe électrogène :} énergie chimique (gasoil) $\rightarrow$ énergie thermique (combustion) $\rightarrow$ énergie mécanique (moteur thermique) $\rightarrow$ énergie électrique (alternateur) $\rightarrow$ énergie lumineuse, mécanique, thermique (récepteurs).
\end{itemize}
\end{aretenir}

\courssub{Démarche et corrigé commenté}

\begin{exoresolu}[Dimensionnement du groupe électrogène de Garoua]
Énoncé : reprendre les cinq consignes de la situation \og Wuro Sallé \fg{} (20 lampes de \SI{60}{W}, 5 ventilateurs de \SI{75}{W}, 2 réfrigérateurs de \SI{150}{W}, rendement de l'alternateur $\eta = \num{0,90}$, délestage de 4~h par jour).
\tcblower
\textbf{Question 1 : puissance électrique totale.}

Les appareils fonctionnent en même temps : leurs puissances s'ajoutent.
\begin{align*}
P_{\text{él}} &= 20 \times P_1 + 5 \times P_2 + 2 \times P_3 \\
&= 20 \times 60 + 5 \times 75 + 2 \times 150 \\
&= 1200 + 375 + 300 \\
&= \SI{1875}{W} \approx \SI{1,9}{kW}.
\end{align*}
L'alternateur devra donc fournir une puissance électrique d'au moins $\boxed{P_{\text{él}} = \SI{1875}{W}}$.

\medskip
\textbf{Question 2 : puissance mécanique du moteur.}

Le rendement de l'alternateur s'écrit $\eta = \dfrac{P_{\text{él}}}{P_{\text{méc}}}$, d'où :
\[
P_{\text{méc}} = \dfrac{P_{\text{él}}}{\eta} = \dfrac{1875}{\num{0,90}} \approx \SI{2083}{W} \approx \SI{2,1}{kW}.
\]
\emph{Vérification dimensionnelle :} un rapport de puissances est sans dimension, donc diviser des watts par un nombre sans dimension redonne des watts : le résultat est homogène.

Le moteur thermique doit fournir environ $\boxed{P_{\text{méc}} \approx \SI{2,1}{kW}}$ à l'arbre de l'alternateur. En pratique, M.~Hamadou choisira un groupe de puissance légèrement supérieure (par exemple \SI{2,5}{kW}) pour disposer d'une marge de sécurité. On remarque que $\SI{208}{W}$ ($2083 - 1875$) sont perdus dans l'alternateur, essentiellement sous forme de chaleur (effet Joule dans les bobinages, frottements) : c'est le prix de la conversion.

\medskip
\textbf{Question 3 : schéma de principe de l'alternateur.}

\begin{center}
\begin{tikzpicture}[scale=1, >=Stealth]
% Stator (cadre fixe)
\draw[very thick, popblue, fill=popblueL] (-3,-1.6) rectangle (3,1.6);
\node[popblue, font=\bfseries] at (0,1.95) {STATOR (fixe) : bobines (induit)};
% Rotor (aimant tournant)
\draw[very thick, poporange, fill=poporangeL] (-0.4,-1.1) rectangle (0.4,1.1);
\node[popdark, font=\bfseries] at (0,0.55) {N};
\node[popdark, font=\bfseries] at (0,-0.55) {S};
\node[poporange, font=\bfseries] at (1.6,0) {ROTOR};
% Axe de rotation
\draw[dashed, popdark] (0,-2.1) -- (0,2.1);
\draw[->, thick, popgreen] (1.1,0.9) arc[start angle=40, end angle=140, radius=1.4];
\node[popgreen] at (0,1.35) {\footnotesize rotation $\omega$};
% Moteur
\draw[thick, poppurple, fill=poppurpleL] (-6.2,-0.7) rectangle (-3.6,0.7);
\node[poppurple, align=center] at (-4.9,0) {\footnotesize Moteur\\ \footnotesize thermique};
\draw[very thick, popdark] (-3.6,0) -- (-0.4,0);
\node[popmuted, font=\footnotesize] at (-2,-0.35) {arbre (énergie mécanique)};
% Sortie électrique
\draw[very thick, poppink] (3,0.5) -- (4.6,0.5);
\draw[very thick, poppink] (3,-0.5) -- (4.6,-0.5);
\node[poppink, align=center] at (5.9,0) {\footnotesize tension alternative\\ \footnotesize $u(t)$ vers le quartier};
\end{tikzpicture}
\end{center}

\emph{Rôles :} le \textbf{rotor} (aimant ou électroaimant) est entraîné en rotation par le moteur thermique : c'est l'\emph{inducteur}, il crée un champ magnétique tournant. Le \textbf{stator}, fixe, porte les bobines : c'est l'\emph{induit}, siège de la tension alternative induite par la variation du flux magnétique. La conversion \og énergie mécanique $\rightarrow$ énergie électrique \fg{} a lieu dans l'interaction rotor--stator.

\medskip
\textbf{Question 4 : chaîne des transformations d'énergie.}

\begin{center}
\begin{tikzpicture}[node distance=0.35cm, >=Stealth,
box/.style={draw=popblue, very thick, rounded corners, fill=popblueL, minimum height=0.9cm, align=center, font=\footnotesize},
org/.style={draw=poporange, thick, rounded corners, fill=poporangeL, font=\scriptsize, align=center}]
\node[box, align=center] (e1){Énergie\\ chimique\\ (gasoil)};
\node[box, right=of e1, align=center] (e2){Énergie\\ thermique\\ (combustion)};
\node[box, right=of e2, align=center] (e3){Énergie\\ mécanique\\ (rotation)};
\node[box, right=of e3, align=center] (e4){Énergie\\ électrique\\ (alternative)};
\node[box, right=of e4, align=center] (e5){Énergie\\ lumineuse,\\ mécanique,\\ thermique};
\draw[->, very thick, popdark] (e1) -- (e2);
\draw[->, very thick, popdark] (e2) -- (e3);
\draw[->, very thick, popdark] (e3) -- (e4);
\draw[->, very thick, popdark] (e4) -- (e5);
\node[org, below=0.25cm of e1] {réservoir};
\node[org, below=0.25cm of e2] {moteur thermique};
\node[org, below=0.25cm of e3] {arbre + rotor};
\node[org, below=0.25cm of e4] {alternateur (stator)};
\node[org, below=0.25cm of e5] {lampes, ventilateurs, réfrigérateurs};
\end{tikzpicture}
\end{center}

Chaque flèche correspond à une conversion réalisée par un organe précis : la combustion dans le moteur, l'entraînement de l'arbre, l'alternateur pour l'étape décisive mécanique $\rightarrow$ électrique, puis les récepteurs du quartier. À chaque étape, une partie de l'énergie est dégradée en chaleur : c'est pourquoi les rendements se multiplient et jamais ne dépassent 1.

\medskip
\textbf{Question 5 : groupe électrogène ou batteries ?}

Une batterie ne \emph{produit} pas d'énergie : elle ne fait que \emph{restituer} une énergie électrique stockée au préalable, en quantité limitée. Pour alimenter \SI{1875}{W} pendant 4~h, il faudrait stocker $E = P_{\text{él}} \times t = 1875 \times 4 = \SI{7500}{Wh} = \SI{7,5}{kWh}$ chaque jour, soit un parc de batteries volumineux, coûteux et à recharger... pendant que le réseau est disponible. L'alternateur, lui, produit de l'énergie électrique \emph{en continu}, tant que le moteur lui fournit de l'énergie mécanique, c'est-à-dire tant qu'il y a du gasoil : quelques litres de carburant suffisent pour tenir toute la nuit, et le délestage peut durer autant que nécessaire. Le groupe électrogène est donc la solution adaptée à une alimentation de plusieurs heures, renouvelable chaque jour.
\end{exoresolu}

\courssub{Bilan de l'activité}

\begin{aretenir}[Ce que cette activité t'a appris]
Cette situation de délestage à Garoua t'a permis de mettre en œuvre l'ensemble des savoir-faire de la leçon :
\begin{itemize}
\item tu as \textbf{dimensionné} une installation en additionnant des puissances et en exploitant le rendement $\eta = \dfrac{P_{\text{él}}}{P_{\text{méc}}}$ : le moteur doit toujours fournir \emph{plus} que ce que l'alternateur restitue ;
\item tu as \textbf{schématisé} un alternateur en distinguant rotor (inducteur tournant) et stator (induit fixe) et en précisant le rôle de chacun ;
\item tu as \textbf{tracé la chaîne énergétique complète} : énergie chimique $\rightarrow$ thermique $\rightarrow$ mécanique $\rightarrow$ électrique, en identifiant l'organe responsable de chaque conversion ;
\item tu as \textbf{argumenté un choix technologique} : l'alternateur, convertisseur d'énergie mécanique en énergie électrique, produit en continu tant qu'on l'alimente, contrairement à une batterie qui ne fait que restituer un stock limité.
\end{itemize}
C'est exactement ce raisonnement qui guide les ingénieurs d'ENEO ou les techniciens des barrages de Lom Pangar et de Nachtigal lorsqu'ils choisissent et dimensionnent les alternateurs qui éclairent le Cameroun.
\end{aretenir}

\newpage

\coursec{Méthodes et exercices résolus}

\courssub{Méthodes et exercices résolus}

\begin{methode}[Identifier les parties essentielles d'un alternateur et leur rôle]
    \textbf{Objectif :} Nommer les composants principaux sur une vue éclatée ou un schéma de principe et associer à chacun sa fonction.
    \begin{enumerate}
        \item \textbf{Repérer le stator (partie fixe) :} C'est l'armature induite (enroulements de cuivre logés dans les encoches) où se produit la f.é.m. alternative triphasée. Le champ magnétique tournant y induit les tensions de sortie.
        \item \textbf{Repérer le rotor (partie mobile) :} C'est l'inducteur (pôles saillants ou lisses) portant l'enroulement d'excitation alimenté en courant continu. Il crée le champ magnétique $\vec{B}$ qui tourne avec l'arbre.
        \item \textbf{Identifier les organes de liaison :} Anneaux collecteurs (bagues) solidaires de l'arbre et balais (graphite) fixes. Ils assurent la continuité électrique du circuit d'excitation (courant continu) vers le rotor tournant.
        \item \textbf{Associer les rôles :}
        \begin{itemize}
            \item \emph{Inducteur (pôles du rotor) :} Crée le champ magnétique $\vec{B}$ dans l'entrefer. Alimenté en courant continu $I_e$ via les anneaux.
            \item \emph{Induit (enroulements du stator) :} Reçoit le flux magnétique variable $\Phi(B)$ et fournit la tension alternative $u(t)$ aux bornes de la charge (puissance forte).
            \item \emph{Arbre / Turbine :} Transmet l'énergie mécanique de rotation (vitesse angulaire $\omega$, fréquence $f$).
            \item \emph{Anneaux/Balais :} Liaison électrique tournante pour l'excitation (courant continu), sans commutation (anneaux lisses).
        \end{itemize}
    \end{enumerate}
    \textbf{Réflexe Probatoire :} Ne pas confondre \emph{inducteur} (source de $\vec{B}$, ici rotor) et \emph{induit} (siège de la f.é.m. induite, ici stator). En centrale, l'induit est au stator pour éviter les frottements balais/anneaux sur les forts courants de sortie.
\end{methode}

\begin{exoresolu}[Reconnaissance des constituants d'un alternateur synchrone]
Sur la figure ci-après représentant un alternateur synchrone à pôles tournants (vue en coupe axiale), on a repéré les éléments A, B, C, D.

\begin{popfigure}
\begin{tikzpicture}[scale=0.9, every node/.style={font=\small}]
    % Stator (carcasse + encoches)
    \draw[thick, popblue] (0,0) circle (3.5);
    \draw[thick, popblue] (0,0) circle (2.5);
    \foreach \a in {0,30,60,90,120,150,180,210,240,270,300,330} {
        \draw[popblue, fill=popblue!15] (\a:2.5) -- (\a:3.5) -- (\a+15:3.5) -- (\a+15:2.5) -- cycle;
    }
    \foreach \a in {15,45,75,105,135,165,195,225,255,285,315,345} {
        \node[popblue, rotate=\a+90] at (\a:3.0) {\tiny Enroulement};
    }
    \node[popblue, align=center] at (0, -4.0) {A : Carcasse + Enroulements statoriques (Induit)};
    
    % Rotor (pôles saillants 4 pôles = 2 paires)
    \draw[thick, poporange] (0,0) circle (2.0);
    \foreach \a in {0,90,180,270} {
        \draw[poporange, fill=poporange!30] (\a:1.2) rectangle (\a+45:2.0);
        \node[poporange, rotate=\a] at (\a+22.5:1.6) {\tiny N/S};
    }
    \node[poporange, align=center] at (0, 2.5) {B : Pôles de l'inducteur (Rotor)};
    
    % Arbre
    \draw[thick, popdark] (-0.3,-2.0) -- (-0.3,-3.5);
    \draw[thick, popdark] (0.3,-2.0) -- (0.3,-3.5);
    \draw[thick, popdark] (-0.3,-3.5) -- (0.3,-3.5);
    \node[popdark, align=center] at (0, -3.8) {C : Arbre};
    
    % Anneaux et balais (excitation)
    \draw[thick, poppurple] (-1.5, -1.5) rectangle (-1.0, -1.0);
    \draw[thick, poppurple] (1.0, -1.5) rectangle (1.5, -1.0);
    \draw[poppurple, fill=poppurple!30] (-1.5, -1.5) rectangle (-1.0, -1.0);
    \draw[poppurple, fill=poppurple!30] (1.0, -1.5) rectangle (1.5, -1.0);
    % Balais
    \draw[thick, poppurple!80!black] (-1.7, -1.5) rectangle (-0.8, -1.8);
    \draw[thick, poppurple!80!black] (0.8, -1.5) rectangle (1.7, -1.8);
    \node[poppurple, align=center] at (-1.25, -2.1) {D : Anneaux};
    \node[poppurple, align=center] at (1.25, -2.1) {D' : Anneaux};
    \node[poppurple!80!black, align=center] at (-1.25, -2.5) {Balai -};
    \node[poppurple!80!black, align=center] at (1.25, -2.5) {Balai +};
    
    % Flèches rotation
    \draw[->, thick, popgreen] (3.0, 1.0) arc (0:180:3.0 and 1.0);
    \node[popgreen] at (3.5, 0) {Sens de rotation $\vec{\omega}$};
    
    % Légende globale
    \node[align=center, font=\footnotesize] at (0, -5.2) {Vue en coupe axiale simplifiée (alternateur synchrone 4 pôles, $p=2$)};
\end{tikzpicture}
\legende{Structure d'un alternateur synchrone à pôles tournants (excitation courant continu via anneaux).}
\end{popfigure}

\textbf{Questions :}
\begin{enumerate}
    \item Nommer chaque élément repéré A, B, C, D (anneaux).
    \item Préciser pour chacun s'il appartient au stator ou au rotor.
    \item Donner le rôle physique de l'élément B et de l'élément D.
    \item Dans quel élément (A ou B) la tension alternative est-elle générée ? Justifier.
\end{enumerate}
\tcblower
\textbf{Solution détaillée :}
\begin{enumerate}
    \item \textbf{A :} Carcasse (stator) supportant l'armature induite (enroulements triphasés dans les encoches). \quad \textbf{B :} Pôles de l'inducteur (électroaimants excités en courant continu). \quad \textbf{C :} Arbre de rotation (liaison mécanique avec la turbine hydraulique). \quad \textbf{D :} Anneaux collecteurs (bagues) solidaires de l'arbre.
    \item \textbf{Stator (fixe) :} A (carcasse + enroulements induits). \quad \textbf{Rotor (tournant) :} B (pôles), C (arbre), D (anneaux).
    \item \textbf{Rôle de B (Inducteur) :} Créer un champ magnétique $\vec{B}$ radial dans l'entrefer. Alimenté en courant continu $I_e$ via les anneaux D, il produit un flux magnétique $\Phi$ qui tourne avec le rotor à la vitesse $\omega$. \quad \textbf{Rôle de D (Anneaux) :} Assurer la liaison électrique tournante entre le circuit d'excitation extérieur (courant continu stabilisé par l'AVR) et l'enroulement d'excitation du rotor (bobines sur les pôles B). Contrairement au collecteur d'une machine à courant continu, ils ne commutent pas le courant (segments isolés), ils le transmettent simplement sans inversion.
    \item La tension alternative est générée dans l'\textbf{enroulement induit logé dans le stator (élément A)}. \textbf{Justification :} D'après la loi de Faraday-Lenz ($e = -\frac{d\Phi}{dt}$), une f.é.m. apparaît dans un conducteur soumis à un flux magnétique variable. Ici, le champ $\vec{B}$ tourne (rotor B) devant les bobines fixes du stator (A). Le flux $\Phi$ traversant chaque spire du stator varie sinusoïdalement, induisant une tension alternative triphasée aux bornes des enroulements du stator. C'est l'avantage majeur des alternateurs synchrones de puissance : la puissance sort sur des pièces fixes, évitant les frottements balais/anneaux pour les forts courants.
\end{enumerate}
\par\medskip\noindent\textbf{Conclusion :} Dans un alternateur synchrone de centrale (Lom Pangar, Nachtigal), l'inducteur est au rotor (pôles tournants excités en continu) et l'induit est au stator (sortie puissance HT triphasée).
\end{exoresolu}

\begin{methode}[Tracer le schéma de principe d'un alternateur monophasé simple]
    \textbf{Objectif :} Représenter schématiquement le dispositif de base : une spire (ou bobine) tournante dans un champ magnétique uniforme, avec les anneaux collecteurs et la charge.
    \begin{enumerate}
        \item \textbf{Champ magnétique :} Dessiner deux pôles plats N et S (aimant en U ou électroaimant) face à face. Tracer les lignes de champ $\vec{B}$ horizontales, uniformes, de N vers S. Légender $\vec{B}$.
        \item \textbf{Bobine tournante (Induit) :} Au centre, dessiner un rectangle (spire) vu de profil (côtés actifs verticaux). Indiquer le sens de rotation $\vec{\omega}$ (flèche courbe autour de l'axe horizontal). Repérer les deux côtés actifs (fils verticaux) coupant les lignes de champ.
        \item \textbf{Liaison électrique tournante :} À chaque extrémité de la spire, souder un anneau métallique (bague) solidaire de l'axe. Dessiner deux balais (graphite) fixes appuyant sur les anneaux.
        \item \textbf{Circuit extérieur :} Relier les balais à une charge $R$ (récepteur) ou un voltmètre.
        \item \textbf{Repères et conventions :} Définir le vecteur normal $\vec{n}$ à la spire (règle de la main droite sur le sens de parcours du courant). Noter l'angle $\theta = \omega t$ entre $\vec{B}$ et $\vec{n}$. Indiquer le sens de la f.é.m. $e(t)$ (convention générateur : $e$ sort du générateur par la borne +).
    \end{enumerate}
    \textbf{Formule clé :} $e(t) = -\frac{d\Phi}{dt} = N B S \omega \sin(\omega t)$ (pour $\theta_0=0$ au flux maximal, $\Phi = NBS\cos\omega t$). Fréquence $f = \frac{\omega}{2\pi} = \frac{N_{\text{tr/min}}}{60}$ (bipolaire $p=1$).
    \textbf{Attention :} Ne pas oublier les anneaux (courant alternatif) \emph{vs} collecteur à segments (courant continu). Le schéma doit montrer la continuité des anneaux (cercles pleins).
\end{methode}

\begin{exoresolu}[Schéma de principe et expression de la f.é.m.]
On souhaite réaliser le schéma de principe d'un alternateur monophasé élémentaire constitué d'une spire rectangulaire de surface $S = \SI{100}{cm^2}$, à $N = 200$ spires, tournant à $\SI{3000}{tr/min}$ dans un champ magnétique uniforme $B = \SI{0,50}{T}$.
\begin{enumerate}
    \item Réaliser le schéma de principe annoté (champ, spire, anneaux, balais, charge, repères $\vec{B}$, $\vec{n}$, $\omega$, $e$).
    \item Déterminer la fréquence $f$ de la tension alternative produite.
    \item Établir l'expression littérale de la f.é.m. instantanée $e(t)$ en choisissant l'origine des temps ($t=0$) lorsque le plan de la spire est perpendiculaire au champ (flux maximal).
    \item Calculer la valeur efficace $U_{\text{eff}}$ de la tension aux bornes de la charge.
\end{enumerate}
\tcblower
\textbf{Solution détaillée :}
\begin{enumerate}
    \item \textbf{Schéma de principe (vue de dessus, axe de rotation vertical) :}
    \begin{popfigure}
    \begin{tikzpicture}[scale=1.1, every node/.style={font=\small}]
        % Aimant (Stator) - Vue de dessus : N à gauche, S à droite
        \draw[thick, popblue] (-3, 1.5) -- (-3, -1.5) node[midway, left=2mm] {\textbf{N}};
        \draw[thick, popred] (3, 1.5) -- (3, -1.5) node[midway, right=2mm] {\textbf{S}};
        % Lignes de champ B horizontal
        \foreach \y in {-1.0, -0.3, 0.3, 1.0} {
            \draw[->, popblue!70, thick] (-2.5, \y) -- (2.5, \y);
        }
        \node[popblue] at (0, 1.8) {$\vec{B}$ uniforme};
        
        % Spire (Rotor) - Position t=0 : flux max => plan spire perpendiculaire à B => vue de dessus = trait vertical
        % On dessine les deux côtés actifs (fils verticaux vus de dessus = points) et les connexions
        % Pour la clarté, on représente la spire par son diamètre (vecteur normal) et l'axe.
        \draw[thick, popdark] (0, -1.5) -- (0, 1.5); % Côtés actifs projetés (superposés au centre)
        \draw[popdark, fill=popdark!20] (-0.3, -1.5) rectangle (0.3, 1.5); % Représentation symbolique de la spire (épaisseur)
        % Axe de rotation (sortant du plan)
        \draw[thick, popgreen] (0,0) circle (0.15);
        \node[popgreen] at (0, -0.5) {Axe $\uparrow$};
        % Anneaux (sur l'axe)
        \draw[fill=poporange!50] (-0.6, 0) circle (0.15);
        \draw[fill=poporange!50] (0.6, 0) circle (0.15);
        % Balais
        \draw[thick, poppurple] (-0.6, -2.0) rectangle (-0.3, -2.3);
        \draw[thick, poppurple] (0.3, -2.0) rectangle (0.6, -2.3);
        \node[poppurple] at (-0.45, -2.5) {Balai -};
        \node[poppurple] at (0.45, -2.5) {Balai +};
        % Charge
        \draw[thick] (-0.45, -2.0) -- (-0.45, -3.0) to[R=$R$, l_={$u(t)$}] (0.45, -3.0) -- (0.45, -2.0);
        % Vecteurs
        \draw[->, thick, popgreen] (0,0) -- (1.5, 0) node[right] {$\vec{n}$};
        \draw[->, thick, popblue] (-2, 1.8) -- (2, 1.8) node[above] {$\vec{B}$};
        % Angle theta = 0
        \node[popgreen] at (0.8, 0.3) {$\theta = 0$};
        % Rotation omega (sortant de page = sens trigonométrique)
        \draw[->, thick, popgreen] (1.0, 0.5) arc (0:90:1.0 and 0.5);
        \node[popgreen] at (1.5, 0.8) {$\vec{\omega}$};
    \end{tikzpicture}
    \legende{Schéma de principe (vue de dessus) : spire tournante dans champ uniforme horizontal, anneaux collecteurs, charge résistive. À $t=0$, $\vec{n}\parallel\vec{B}$ (flux max).}
    \end{popfigure}
    
    \item \textbf{Fréquence :}
    $N = \SI{3000}{tr/min} = \frac{3000}{60} = \SI{50}{tr/s} = \SI{50}{Hz}$.
    \textit{Rappel :} Un tour mécanique = une période électrique pour un alternateur bipolaire ($p=1$ paire de pôles). $f = \frac{N}{60} = \SI{50}{Hz}$.
    
    \item \textbf{Expression de $e(t)$ :}
    Flux traversant la spire : $\Phi(t) = N B S \cos\theta(t)$ avec $\theta(t) = \omega t$ (origine $t=0$ : $\theta=0$, flux max).
    $\omega = 2\pi f = \SI{100\pi}{rad.s^{-1}}$.
    Loi de Faraday (convention générateur) : $e(t) = -\frac{d\Phi}{dt} = -N B S (-\omega \sin(\omega t)) = N B S \omega \sin(\omega t)$.
    \textbf{Valeurs numériques :}
    $S = \SI{100}{cm^2} = \SI{0,01}{m^2}$.
    Amplitude $E_{\text{max}} = N B S \omega = 200 \times 0,50 \times 0,01 \times 100\pi = 100\pi \approx \SI{314}{V}$.
    $e(t) = \SI{314}{\volt} \sin(100\pi t)$ ($t$ en secondes).
    
    \item \textbf{Tension efficace :}
    $U_{\text{eff}} = \frac{E_{\text{max}}}{\sqrt{2}} = \frac{314}{\sqrt{2}} \approx \SI{222}{V}$.
\end{enumerate}
\par\medskip\noindent\textbf{Conclusion :} Cet alternateur élémentaire produit une tension sinusoïdale de $\SI{50}{Hz}$ et $\SI{222}{V}$ efficaces, compatible avec le réseau domestique camerounais (monophasé $\SI{230}{V}$).
\end{exoresolu}

\begin{methode}[Expliquer le principe de fonctionnement par induction électromagnétique]
    \textbf{Objectif :} Justifier l'apparition d'une tension alternative aux bornes de l'induit à partir de la loi de Faraday-Lenz et du mouvement relatif inducteur/induit.
    \begin{enumerate}
        \item \textbf{Rappeler la loi de Faraday :} $e = -\frac{d\Phi}{dt}$ (convention générateur) ou $u = \frac{d\Phi}{dt}$ (convention récepteur). Préciser que $\Phi = \iint \vec{B}\cdot\vec{n}\,dS = B S \cos\theta$ pour une spire plane dans champ uniforme.
        \item \textbf{Analyser le mouvement :} L'inducteur (pôles) crée $\vec{B}$. L'induit (bobines) tourne (ou $\vec{B}$ tourne devant l'induit fixe). L'angle $\theta$ entre $\vec{B}$ et la normale $\vec{n}$ aux spires varie : $\theta = \omega t + \varphi_0$.
        \item \textbf{Calculer la dérivée :} $\Phi(t) = \Phi_{\text{max}}\cos(\omega t + \varphi_0) \Rightarrow e(t) = \Phi_{\text{max}}\omega \sin(\omega t + \varphi_0)$.
        \item \textbf{Conclure sur la nature du signal :} C'est une tension \emph{sinusoïdale} de pulsation $\omega$ (même que la rotation mécanique pour $p=1$), d'amplitude proportionnelle à $B$, $S$, $N$, $\omega$.
        \item \textbf{Cas triphasé (usuel) :} 3 enroulements décalés de $120^\circ$ (ou $2\pi/3$ rad) dans l'espace $\Rightarrow$ 3 tensions décalées de $2\pi/3$ dans le temps.
    \end{enumerate}
    \textbf{Mots-clés pour la rédaction :} \cle{Flux magnétique}, \cle{Variation}, \cle{Dérivée}, \cle{Sinusoïdale}, \cle{Fréquence}, \cle{Paires de pôles}.
\end{methode}

\begin{exoresolu}[Fonctionnement et fréquence d'un alternateur de centrale]
L'alternateur de la centrale hydroélectrique de Nachtigal (Cameroun) est une machine synchrone à $p = 12$ paires de pôles. La turbine Francis entraîne l'alternateur à une vitesse de rotation $N = \SI{250}{tr/min}$.
\begin{enumerate}
    \item Expliquer physiquement pourquoi la tension aux bornes du stator est alternative. Utiliser la loi de Faraday.
    \item Calculer la fréquence $f$ de la tension produite. Vérifier qu'elle correspond au standard du réseau interconnecté national (SONATREL/ENEO).
    \item Si l'on veut doubler la fréquence sans changer la vitesse de la turbine, que faut-il modifier sur la machine ? Est-ce réaliste ?
\end{enumerate}
\tcblower
\textbf{Solution détaillée :}
\begin{enumerate}
    \item \textbf{Explication physique :}
    Le rotor (inducteur) porte $2p = 24$ pôles magnétiques alternés N/S. Il est mis en rotation par la turbine à la vitesse angulaire $\Omega = 2\pi N/60$.
    Le stator (induit) possède des enroulements fixes répartis sur sa circonférence.
    Lorsqu'un pôle Nord passe devant une bobine du stator, le flux magnétique $\Phi$ traversant cette bobine est maximal et positif. Lorsque le pôle Sud suivant passe, le flux devient maximal et négatif. Entre les deux, le flux passe par zéro.
    Comme le rotor tourne de manière continue et uniforme, le flux $\Phi(t)$ traversant chaque enroulement statorique varie \emph{périodiquement} et \emph{sinusoïdalement}.
    D'après la \textbf{loi de Faraday-Lenz} ($e = -d\Phi/dt$), toute variation de flux magnétique à travers un circuit conducteur engendre une force électromotrice (f.é.m.) aux bornes de ce circuit. Comme la variation de flux est sinusoïdale, la f.é.m. $e(t)$ est elle aussi sinusoïdale : c'est une \textbf{tension alternative}.
    
    \item \textbf{Calcul de la fréquence :}
    Relation fondamentale : $f = p \times \frac{N}{60}$ (avec $N$ en tr/min).
    $p = 12$ paires de pôles. $N = \SI{250}{tr/min}$.
    $f = 12 \times \frac{250}{60} = 12 \times \frac{25}{6} = 2 \times 25 = \SI{50}{Hz}$.
    \textbf{Vérification :} Le réseau électrique camerounais (SONATREL/ENEO) fonctionne en $\SI{50}{Hz}$. La machine est bien synchronisée sur le réseau.
    
    \item \textbf{Modification pour doubler $f$ :}
    La formule $f = p \cdot N/60$ montre que pour doubler $f$ à $N$ constant, il faut \textbf{doubler le nombre de paires de pôles $p$} (passer de 12 à 24 paires, soit 48 pôles).
    \textbf{Réalisme :} C'est \emph{théoriquement possible} mais \emph{peu réaliste} sur une machine existante : cela impose de reconstruire totalement le rotor (nouvelles empreintes, nouveaux pôles) et le stator (nouveaux encoches, nouveaux enroulements adaptés au nouveau pas de pôle). En pratique, on choisit $p$ et $N$ dès la conception pour obtenir $\SI{50}{Hz}$ (ou $\SI{60}{Hz}$). Changer $f$ impose de changer de machine ou d'utiliser un convertisseur de fréquence (électronique de puissance).
\end{enumerate}
\par\medskip\noindent\textbf{Conclusion :} La fréquence électrique est verrouillée sur la vitesse mécanique par le nombre de paires de pôles : $f = p \cdot f_{\text{méc}}$. C'est la base de la synchronisation sur le réseau.
\end{exoresolu}

\begin{methode}[Établir le bilan des transformations d'énergie]
    \textbf{Objectif :} Retracer la chaîne énergétique depuis la source primaire jusqu'à l'énergie électrique distribuée, en identifiant les convertisseurs et les rendements.
    \begin{enumerate}
        \item \textbf{Identifier la source primaire :} Énergie potentielle de pesanteur (barrage), énergie thermique (centrale thermique à gaz/fuel/charbon), énergie cinétique (éolien), énergie nucléaire.
        \item \textbf{Lister les convertisseurs successifs :}
        \begin{itemize}
            \item \emph{Turbine (hydraulique, à gaz, à vapeur, éolienne) :} Énergie primaire $\to$ Énergie mécanique de rotation (arbre). Rendement $\eta_t$.
            \item \emph{Alternateur (machine synchrone) :} Énergie mécanique $\to$ Énergie électrique alternative (triphasée). Rendement $\eta_a$ (très élevé, $95-99\%$).
            \item \emph{Transformateur élévateur :} Tension moyenne ($\SI{10}{kV}$-$\SI{20}{kV}$) $\to$ Haute tension ($\SI{90}{kV}$-$\SI{225}{kV}$) pour le transport. Rendement $\eta_{tr} \approx 99\%$.
        \end{itemize}
        \item \textbf{Écrire la chaîne :} $E_{\text{primaire}} \xrightarrow{\text{Turbine}} E_{\text{méc}} \xrightarrow{\text{Alternateur}} E_{\text{élect (BT/MT)}} \xrightarrow{\text{Transfo}} E_{\text{élect (HT)}}$.
        \item \textbf{Calculer le rendement global :} $\eta_{\text{global}} = \eta_t \times \eta_a \times \eta_{tr} \times \dots$
        \item \textbf{Noter les pertes :} Joule (cuivre/fer), mécaniques (frottements, ventilation), magnétiques (hystérésis, courants de Foucault).
    \end{enumerate}
    \textbf{Exigence Probatoire :} Savoir nommer les dispositifs (Turbine, Alternateur, Transformateur) et le sens de conversion. Ne pas confondre "production" et "transformation".
\end{methode}

\begin{exoresolu}[Chaîne énergétique du barrage de Lom Pangar]
La centrale hydroélectrique de Lom Pangar (Cameroun) a une puissance installée de $\SI{30}{MW}$ (phase 1). La chute d'eau brute est $H = \SI{45}{m}$. Le débit turbiné maximal est $Q = \SI{85}{m^3/s}$. La masse volumique de l'eau est $\rho = \SI{1000}{kg/m^3}$, $g = \SI{9,81}{N/kg}$.
\begin{enumerate}
    \item Calculer la puissance hydraulique brute $P_{\text{hyd}}$ disponible.
    \item Sachant que le rendement de la turbine Francis est $\eta_t = 0,92$ et celui de l'alternateur $\eta_a = 0,97$, calculer la puissance électrique active $P_{\text{élec}}$ injectée sur le réseau (avant transformateur).
    \item Déduire le rendement global de la chaîne "eau $\to$ électricité" (hors transformateur).
    \item Compléter le diagramme des transferts d'énergie (puissances) en indiquant les pertes à chaque étage.
\end{enumerate}
\tcblower
\textbf{Solution détaillée :}
\begin{enumerate}
    \item \textbf{Puissance hydraulique :}
    $P_{\text{hyd}} = \rho \cdot g \cdot Q \cdot H$
    $P_{\text{hyd}} = 1000 \times 9,81 \times 85 \times 45$
    $P_{\text{hyd}} = \SI{37523250}{W} \approx \SI{37,5}{MW}$.
    \textit{Analyse dimensionnelle :} $[\rho] = \SI{}{kg/m^3}$, $[g] = \SI{}{m/s^2}$, $[Q] = \SI{}{m^3/s}$, $[H] = \SI{}{m} \Rightarrow \SI{}{kg.m^2/s^3} = \SI{}{W}$. Correct.
    
    \item \textbf{Puissance mécanique sur l'arbre :}
    $P_{\text{méc}} = \eta_t \times P_{\text{hyd}} = 0,92 \times 37,5 = \SI{34,5}{MW}$.
    \textbf{Puissance électrique (sortie alternateur) :}
    $P_{\text{élec}} = \eta_a \times P_{\text{méc}} = 0,97 \times 34,5 = \SI{33,5}{MW}$.
    (Note : La puissance installée annoncée $\SI{30}{MW}$ est une valeur nominale garantie, légèrement inférieure au max théorique pour des marges de sécurité).
    
    \item \textbf{Rendement global :}
    $\eta_{\text{global}} = \frac{P_{\text{élec}}}{P_{\text{hyd}}} = \eta_t \times \eta_a = 0,92 \times 0,97 = 0,8924 \approx \mathbf{89,2\%}$.
    
    \item \textbf{Diagramme des puissances (Bilan) :}
    \begin{center}
    \begin{tabularx}{\linewidth}{|l|X|X|}\hline
    \rowcolor{popblueL} \textbf{Étage} & \textbf{Puissance d'entrée} & \textbf{Puissance de sortie / Pertes} \\ \hline
    Ressource (Eau) & $P_{\text{hyd}} = \SI{37,5}{MW}$ & -- \\ \hline
    Turbine Francis & $\SI{37,5}{MW}$ & $P_{\text{méc}} = \SI{34,5}{MW}$ \newline Pertes $P_{\text{pertes},t} = \SI{3,0}{MW}$ (frottements, turbulences) \\ \hline
    Alternateur Synchrone & $\SI{34,5}{MW}$ & $P_{\text{élec}} = \SI{33,5}{MW}$ \newline Pertes $P_{\text{pertes},a} = \SI{1,0}{MW}$ (Joule stator/rotor, fer, mécanique) \\ \hline
    Transformateur Élévateur & $\SI{33,5}{MW}$ & $P_{\text{réseau}} \approx \SI{33,2}{MW}$ \newline Pertes $\approx \SI{0,3}{MW}$ \\ \hline
    \end{tabularx}
    \end{center}
\end{enumerate}
\par\medskip\noindent\textbf{Conclusion :} Sur $\SI{37,5}{MW}$ d'énergie hydraulique, $\SI{33,5}{MW}$ sont convertis en électricité utilisable (rendement global $\approx 89\%$). Les $\SI{4}{MW}$ restants sont dissipés en chaleur (pertes mécaniques, Joule, magnétiques). L'alternateur est le maillon le plus performant ($\eta > 97\%$).
\end{exoresolu}

\begin{methode}[Réaliser / Décrire un dispositif produisant un courant alternatif (TP / Expérience)]
    \textbf{Objectif :} Concevoir ou analyser un montage de laboratoire mettant en évidence la production d'une f.é.m. alternative par induction.
    \begin{enumerate}
        \item \textbf{Matériel minimal :} Aimant en U (ou néodyme fort) + bobine (solénoïde) sur support rotatif (manivelle) OU aimant tournant devant bobine fixe + galvanomètre / oscilloscope / LED + résistance.
        \item \textbf{Montage "Alternateur manuel" (statique/rotatif) :}
        \begin{itemize}
            \item Fixer une bobine (ex: $N=500$ spires) sur un châssis.
            \item Monter un aimant puissant sur un axe avec manivelle (rotor).
            \item Relier les bornes de la bobine à un oscilloscope (mode AC) ou un multimètre AC / LED (avec résistance de limitation).
        \end{itemize}
        \item \textbf{Protocole :}
        \begin{enumerate}
            \item Tourner la manivelle à vitesse constante. Observer la tension sur l'oscilloscope : signal sinusoïdal.
            \item Varier la vitesse de rotation : constater que la fréquence $f$ et l'amplitude $U_{\text{max}}$ augmentent avec la vitesse.
            \item Inverser le sens de rotation : le signal est déphasé de $\pi$ (inversé).
            \item Approcher/éloigner l'aimant (sans rotation) : impulsion transitoire (courant induit pulsé), pas d'alternatif soutenu.
        \end{enumerate}
        \item \textbf{Interprétation :} $e(t) = -N \frac{d\Phi}{dt}$. Rotation $\Rightarrow \Phi(t) = \Phi_{\text{max}}\cos(\omega t) \Rightarrow e(t) \propto \omega \sin(\omega t)$.
        \item \textbf{Compte-rendu :} Schéma du montage, courbe $u(t)$ relevée, mesures $f$, $U_{\text{max}}$ pour 2 vitesses, conclusion liant $\omega$, $f$, $U_{\text{max}}$.
    \end{enumerate}
    \textbf{Sécurité :} Aimants puissants (risque pincement, pacemaker), pièces tournantes (cheveux, vêtements).
\end{methode}

\begin{exoresolu}[Analyse d'un TP : Alternateur à aimant permanent tournant]
Au laboratoire, un élève réalise un alternateur élémentaire : un aimant cylindrique néodyme ($B_{\text{surf}} \approx \SI{0,4}{T}$) est fixé sur un moteur pas-à-pas utilisé comme générateur (vitesse imposée). Le stator est une bobine de $N = 1000$ spires, section $S = \SI{2}{cm^2}$, résistance interne $r = \SI{50}{\ohm}$. L'ensemble est relié à un oscilloscope (entrée $\SI{1}{M\ohm}$).
On impose une rotation à $f_{\text{rot}} = \SI{10}{Hz}$ ($\SI{600}{tr/min}$).
\begin{enumerate}
    \item Préciser le rôle du moteur pas-à-pas dans ce montage.
    \item Calculer l'amplitude théorique $E_{\text{max}}$ de la f.é.m. (on suppose que le flux traverse la bobine de façon optimale : $\Phi_{\text{max}} = B S$).
    \item L'oscilloscope affiche une sinusoïde de période $T = \SI{100}{ms}$ et d'amplitude crête $U_{\text{max}} = \SI{4,5}{V}$. Commenter la différence avec le calcul.
    \item Si on court-circuite la bobine (fil direct aux bornes), quel courant maximal traverse-t-elle ? Pourquoi cette expérience est-elle déconseillée ?
\end{enumerate}
\tcblower
\textbf{Solution détaillée :}
\begin{enumerate}
    \item Le moteur pas-à-pas sert ici de \textbf{moteur d'entraînement (premier moteur)}. Il impose une rotation \emph{rigoureusement constante et connue} au rotor (aimant), simulant la turbine. C'est une source d'énergie mécanique à vitesse imposée.
    \item \textbf{Calcul théorique :}
    $\omega = 2\pi f = \SI{20\pi}{rad/s}$.
    $S = \SI{2}{cm^2} = \SI{2e-4}{m^2}$.
    $\Phi_{\text{max}} = B S = 0,4 \times 2 \times 10^{-4} = \SI{8e-5}{Wb}$.
    $E_{\text{max}} = N \omega \Phi_{\text{max}} = 1000 \times 20\pi \times 8 \times 10^{-5} = 1,6\pi \approx \SI{5,03}{V}$.
    
    \item \textbf{Commentaire :}
    Mesure : $U_{\text{max}} = \SI{4,5}{V}$. Calcul : $\SI{5,03}{V}$. Écart relatif $\approx 10\%$.
    \textbf{Causes probables :}
    \begin{itemize}
        \item Le champ $B$ n'est pas uniforme sur toute la section $S$ (champ radial d'un aimant cylindrique, lignes courbes). Flux réel $< B_{\text{surf}} \times S$.
        \item Couplage magnétique imparfait (entrefer, fuites de flux).
        \item Chargement de l'oscilloscope ($\SI{1}{M\ohm}$) négligeable par rapport à $r=\SI{50}{\ohm}$, donc pas de chute de tension interne significative ici.
    \end{itemize}
    La période mesurée $T=\SI{100}{ms}$ correspond bien à $f=\SI{10}{Hz}$ (1 paire de pôles sur l'aimant cylindrique $\Rightarrow$ 1 alternance par tour). Cohérence validée.
    
    \item \textbf{Court-circuit :}
    $I_{\text{max}} = \frac{E_{\text{max}}}{r} = \frac{5,03}{50} \approx \SI{0,10}{A}$ ($\SI{100}{mA}$).
    \textbf{Danger / Déconseillé :} Ce courant ($\SI{100}{mA}$) dans une bobine de fil fin ($r=\SI{50}{\ohm}$) provoque un échauffement Joule $P = r I^2_{\text{eff}} \approx 50 \times (0,07)^2 \approx \SI{0,25}{W}$ concentré dans un petit volume. Risque de \textbf{surchauffe locale}, dégradation de l'isolant (vernis), court-circuit inter-spires, destruction de la bobine. De plus, la réaction magnétique (freinage électromagnétique / loi de Lenz) bloque brutalement le rotor, risquant d'endommager le moteur d'entraînement ou l'axe.
\end{enumerate}
\par\medskip\noindent\textbf{Conclusion :} Le montage valide la loi $E_{\text{max}} \propto N B S \omega$. L'écart expérience/théorie illustre la géométrie réelle du champ. Le court-circuit est destructeur (thermique + mécanique).
\end{exoresolu}

\begin{methode}[Caractériser un alternateur : Fréquence, Tension, Puissance]
    \textbf{Objectif :} Utiliser les formelles liaisons $f = p \frac{N}{60}$, $E_{\text{max}} = N_s B S \omega$ (ou $K \Phi \omega$), $P = \sqrt{3} U I \cos\varphi$ (triphasé) pour dimensionner ou vérifier une machine.
    \begin{enumerate}
        \item \textbf{Identifier le type :} Monophasé / Triphasé (usuel réseau). Nombre de paires de pôles $p$.
        \item \textbf{Lier vitesse et fréquence :} $f = p \cdot n$ ($n$ en tr/s) ou $f = p \cdot \frac{N}{60}$ ($N$ en tr/min). \cle{Formule clé Probatoire}.
        \item \textbf{Lier tension et excitation/vitesse :} $U \propto \Phi \cdot f \propto I_{\text{exc}} \cdot f$ (si non saturé). Régulation de tension par variation du courant d'excitation $I_{\text{exc}}$ (via AVR - Automatic Voltage Regulator).
        \item \textbf{Puissance apparente :} $S_n = \sqrt{3} U_n I_n$ (triphasé). Puissance active $P = S \cos\varphi$.
        \item \textbf{Rendement :} $\eta = \frac{P_{\text{sortie}}}{P_{\text{méc, entrée}}} = \frac{P_{\text{élec}}}{P_{\text{élec}} + P_{\text{pertes}}}$. Pertes = $P_{\text{Joule}} (stator+rotor) + P_{\text{fer}} + P_{\text{méc}}$.
    \end{enumerate}
    \textbf{Unités :} $f$ en Hz, $N$ en tr/min, $\omega$ en rad/s, $U,I$ en V,A (valeurs efficaces réseau), $P,S$ en W, VA.
\end{methode}

\begin{exoresolu}[Dimensionnement d'un groupe électrogène de secours]
Un hôpital de district au Cameroun installe un groupe électrogène diesel de secours (marque Kirloskar / Cummins type). L'alternateur est triphasé, $\SI{400}{V}$ entre phases (régime étoile), $\SI{50}{Hz}$, puissance apparente nominale $S_n = \SI{100}{kVA}$, $\cos\varphi_n = 0,8$. Le moteur diesel tourne à $N = \SI{1500}{tr/min}$.
\begin{enumerate}
    \item Déterminer le nombre de paires de pôles $p$ de l'alternateur.
    \item Calculer le courant nominal $I_n$ par phase.
    \item La charge réelle de l'hôpital en secours est de $P = \SI{60}{kW}$ avec $\cos\varphi = 0,9$. Calculer la puissance apparente $S$ demandée et le courant $I$ par phase. L'alternateur est-il surchargé en courant ? En puissance active ?
    \item Si le moteur diesel ralentit à $\SI{1450}{tr/min}$ (sous-charge soudaine), quelle devient la fréquence ? Quel risque pour les équipements médicaux (respirateurs, scanners) ?
\end{enumerate}
\tcblower
\textbf{Solution détaillée :}
\begin{enumerate}
    \item \textbf{Nombre de paires de pôles :}
    $f = p \frac{N}{60} \Rightarrow p = \frac{f \times 60}{N} = \frac{50 \times 60}{1500} = \frac{3000}{1500} = \mathbf{2}$ paires de pôles (4 pôles).
    \textit{Note :} C'est la standard pour les machines $\SI{1500}{tr/min}$ / $\SI{50}{Hz}$ (ou $\SI{1800}{tr/min}$ / $\SI{60}{Hz}$).
    
    \item \textbf{Courant nominal :}
    Triphasé étoile : $U_{\text{phase}} = \frac{U_{\text{composée}}}{\sqrt{3}} = \frac{400}{\sqrt{3}} \approx \SI{231}{V}$.
    $S_n = \sqrt{3} U_n I_n \Rightarrow I_n = \frac{S_n}{\sqrt{3} U_n} = \frac{100 \times 10^3}{\sqrt{3} \times 400} = \frac{100000}{692,8} \approx \mathbf{\SI{144,3}{A}}$.
    
    \item \textbf{Charge réelle :}
    $P = \SI{60}{kW}$, $\cos\varphi = 0,9$.
    $S = \frac{P}{\cos\varphi} = \frac{60}{0,9} \approx \mathbf{\SI{66,7}{kVA}}$.
    $I = \frac{S}{\sqrt{3} U_n} = \frac{66700}{692,8} \approx \mathbf{\SI{96,3}{A}}$.
    \textbf{Comparaison :}
    \begin{itemize}
        \item Courant : $\SI{96,3}{A} < \SI{144,3}{A}$ ($\approx 67\% I_n$). \textbf{Pas de surcharge courant.}
        \item Puissance active : $P = \SI{60}{kW}$. Puissance active nominale $P_n = S_n \cos\varphi_n = 100 \times 0,8 = \SI{80}{kW}$. $60 < 80$. \textbf{Pas de surcharge active.}
        \item Puissance apparente : $\SI{66,7}{kVA} < \SI{100}{kVA}$. \textbf{Marge confortable.}
    \end{itemize}
    L'alternateur fonctionne dans de bonnes conditions (charge partielle).
    
    \item \textbf{Variation de fréquence :}
    $f' = p \frac{N'}{60} = 2 \times \frac{1450}{60} = \frac{2900}{60} \approx \mathbf{\SI{48,3}{Hz}}$.
    \textbf{Risques :}
    \begin{itemize}
        \item \textbf{Moteurs asynchrones (ventilation, pompes, compresseurs frigos) :} Vitesse $\propto f$. Ralentissement $\Rightarrow$ échauffement (ventilation propre réduite), couple réduit $\propto f^2$, risque de blocage/déclenchement thermique.
        \item \textbf{Transformateurs / Ballasts :} Flux $\Phi \propto U/f$. Si $U$ constant et $f$ baisse $\Rightarrow$ saturation magnétique $\Rightarrow$ courant d'aimantation excessif, surchauffe, bruit.
        \item \textbf{Électronique médicale (alimentations à découpage) :} Généralement tolèrent $\SI{47}{-}\SI{63}{Hz}$, mais horloges internes, synchronisation réseau, peuvent dysfonctionner.
        \item \textbf{Régulation AVR :} L'alternateur doit maintenir $U=\SI{400}{V}$ malgré la baisse de $f$ (régulation $U/f$ constante souvent requise pour machines tournantes, mais ici charge statique $\Rightarrow$ tension maintenue par AVR).
    \end{itemize}
    \textbf{Conclusion :} Une chute à $\SI{48,3}{Hz}$ est \textbf{critique} pour les équipements à induction. Le régulateur de vitesse (gouvernor) du diesel doit réagir en $< \SI{1}{-}\SI{2}{s}$ pour ramener $N=\SI{1500}{tr/min}$.
\end{enumerate}
\par\medskip\noindent\textbf{Conclusion :} Le dimensionnement ($p=2$, $S_n=\SI{100}{kVA}$) est cohérent. La stabilité de la fréquence (gouvernor) est vitale pour la qualité de l'énergie en site isolé (hôpital, usine).
\end{exoresolu}

\begin{attention}[Les 5 erreurs qui coûtent des points au Probatoire]
    \begin{enumerate}
        \item \textbf{Confondre Inducteur et Induit :} Dire "l'induit tourne" dans un alternateur de centrale (stator = induit fixe !). \textbf{Réflexe :} L'induit est \emph{toujours} là où sort la puissance principale (courants forts). En centrale $\rightarrow$ Stator.
        \item \textbf{Oublier les paires de pôles dans la formule $f = p N/60$ :} Utiliser $f = N/60$ (valable seulement pour $p=1$, bipolaire). \textbf{Sanction :} Résultat faux par facteur $p$ (ex: $\SI{250}{tr/min} \to \SI{4,16}{Hz}$ au lieu de $\SI{50}{Hz}$ pour $p=12$).
        \item \textbf{Mélanger valeurs efficaces et valeurs maximales (crête) :} Appliquer $U = 4,44 f N \Phi$ (formule efficace) avec $\Phi_{\text{max}}$, ou calculer $I = P/U$ avec $U_{\text{max}}$. \textbf{Règle :} En régime sinusoïdal établi, \emph{toutes} les grandeurs de circuit ($U, I, P, S, Q$) sont des \textbf{valeurs efficaces (RMS)} sauf mention explicite "amplitude", "crête", "instantané".
        \item \textbf{Schémas de principe incomplets :} Oublier les \textbf{anneaux collecteurs} (dessiner un collecteur à segments = machine continue), ne pas représenter le \textbf{champ magnétique $\vec{B}$}, ne pas indiquer le \textbf{sens de rotation $\omega$} ni le \textbf{repère de la f.é.m. $e$} (convention générateur).
        \item \textbf{Bilan d'énergie sans pertes / Rendement $> 1$ :} Écrire $P_{\text{élec}} = P_{\text{méc}}$ ou $\eta = P_{\text{méc}}/P_{\text{élec}} > 1$. \textbf{Rappel :} $\eta = \frac{\text{Énergie utile fournie}}{\text{Energie reçue}} < 1$. Toujours détailler : $P_{\text{entr}} = P_{\text{sort}} + \Sigma P_{\text{pertes}}$.
    \end{enumerate}
\end{attention}

\newpage

\coursec{Exercices}

\courssub{Je vérifie mes connaissances}

\begin{exercice}[QCM : Rôle et principe]
Parmi les affirmations suivantes, une seule est correcte. Laquelle ?
\begin{enumerate}
    \item L'alternateur transforme l'énergie électrique en énergie mécanique.
    \item L'alternateur transforme l'énergie mécanique en énergie électrique alternative.
    \item L'alternateur transforme l'énergie chimique en énergie électrique continue.
    \item L'alternateur ne fonctionne que s'il est alimenté par une pile.
\end{enumerate}
\end{exercice}

\begin{exercice}[Vrai ou Faux : Parties et fonctionnement]
Pour chaque affirmation, répondre Vrai ou Faux et justifier brièvement.
\begin{enumerate}
    \item La partie tournante d'un alternateur s'appelle le stator.
    \item La tension alternative est induite dans les bobines du rotor.
    \item L'excitation magnétique peut être assurée par un aimant permanent ou par un électroaimant alimenté en courant continu.
    \item La fréquence de la tension de sortie dépend uniquement du nombre de paires de pôles, pas de la vitesse de rotation.
\end{enumerate}
\end{exercice}

\begin{exercice}[Définitions et unités]
Donner la définition précise et l'unité SI (si applicable) des grandeurs suivantes :
\begin{enumerate}
    \item La fréquence $f$ d'un signal alternatif.
    \item La pulsation $\omega$ (lien avec $f$).
    \item Le rendement $\eta$ d'un alternateur.
    \item La f.é.m. induite $e(t)$ dans une bobine tournante dans un champ magnétique uniforme (expression en fonction de $B$, $S$, $N$, $\omega$, $t$).
\end{enumerate}
\end{exercice}

\begin{exercice}[Calculs directs : Fréquence et période]
Un rotor d'alternateur tourne à une vitesse de rotation $n = 3000~\text{tr/min}$. Il possède $p = 2$ paires de pôles.
\begin{enumerate}
    \item Calculer la fréquence $f$ de la tension alternative produite (en Hz).
    \item En déduire la période $T$ (en ms).
    \item Si l'on veut obtenir une fréquence de $50~\text{Hz}$ avec ce même alternateur, quelle doit être la vitesse de rotation en tr/min ?
\end{enumerate}
\end{exercice}

\courssub{Je m'entraîne}

\begin{exercice}[Schéma de principe et identification]
On considère un alternateur monophasé simple : un aimant permanent (rotor) tourne à l'intérieur d'une bobine fixe (stator) à vitesse constante.
\begin{enumerate}
    \item Faire un schéma de principe légendé de cet alternateur en vue de face (coupe transversale). On y fera apparaître : le rotor (aimant N/S), le stator (bobine), l'arbre d'entraînement, le sens de rotation, les bornes de sortie $A$ et $B$.
    \item Indiquer sur le schéma, par une flèche, le sens du champ magnétique $\vec{B}$ dans l'entrefer au moment où le pôle Nord passe devant la bobine.
    \item Préciser le rôle de chaque partie : rotor, stator, arbre, bagues collectrices (s'il y en a) ou fils de sortie directs.
\end{enumerate}
\end{exercice}

\begin{exercice}[Chaîne énergétique : Microcentrale villageoise]
Une microcentrale hydroélectrique villageoise sur la rivière Méfou (région du Centre) utilise une turbine Pelton entraînée par une chute d'eau de hauteur $h = 50~\text{m}$. Le débit d'eau est $Q = 0,2~\text{m}^3/\text{s}$. La turbine entraîne un alternateur synchrone triphasé.
Données : $\rho_{\text{eau}} = 1000~\text{kg/m}^3$, $g = 9,81~\text{N/kg}$, rendement turbine $\eta_t = 0,85$, rendement alternateur $\eta_a = 0,92$.
\begin{enumerate}
    \item Calculer la puissance hydraulique $P_h$ disponible (puissance mécanique théorique de l'eau).
    \item Calculer la puissance mécanique $P_m$ transmise à l'arbre de l'alternateur (sortie turbine).
    \item Calculer la puissance électrique active $P_e$ fournie par l'alternateur (sortie bornes).
    \item Dresser la chaîne complète des transformations d'énergie depuis l'eau en amont jusqu'au réseau électrique du village, en indiquant les rendements à chaque étape.
\end{enumerate}
\end{exercice}

\begin{exercice}[Alternateur de voiture vs Groupe électrogène]
\begin{enumerate}
    \item Un alternateur de voiture (type Lundell) tourne à $n = 2000~\text{tr/min}$ (régime moteur). Il a $p = 6$ paires de pôles. Calculer la fréquence de la tension alternative produite avant redressement.
    \item Pourquoi cette tension doit-elle être redressée et régulée avant de charger la batterie de $12~\text{V}$ ?
    \item Un groupe électrogène de chantier (marque courante au Cameroun : Elemax, Honda, ou Kipor) doit délivrer du $230~\text{V} / 50~\text{Hz}$ monophasé. Son moteur thermique tourne à $3000~\text{tr/min}$. Combien de paires de pôles doit avoir l'alternateur synchrone couplé ?
    \item Comparer la source d'énergie mécanique primaire dans les deux cas.
\end{enumerate}
\end{exercice}

\begin{exercice}[Réalisation : Alternateur de démonstration]
On souhaite réaliser en classe un petit alternateur pour allumer une LED rouge ($U_{\text{seuil}} \approx 1,8~\text{V}$) à l'aide d'un aimant néodyme et d'une bobine en fil émaillé.
Matériel : aimant cylindrique $B \approx 0,5~\text{T}$, diamètre $d = 10~\text{mm}$ ; fil émaillé $\varnothing 0,2~\text{mm}$ ; manivelle ; support imprimé 3D ou bois.
\begin{enumerate}
    \item Proposer un montage (schéma descriptif) : bobine fixe (stator) et aimant tournant (rotor) ou l'inverse. Justifier le choix pour la simplicité de réalisation.
    \item Estimer le nombre de spires $N$ nécessaire pour obtenir une tension de crête $\hat{u} \ge 3~\text{V}$ si l'on tourne la manivelle à $n = 120~\text{tr/min}$ ($2~\text{tr/s}$). On suppose que le flux traverse la bobine de manière sinusoïdale et que la surface utile est celle de l'aimant ($S = \pi d^2/4$).
    \item Expliquer pourquoi une LED seule s'éteint la moitié du temps et proposer un composant simple à ajouter pour qu'elle reste allumée en continu (sans condensateur).
\end{enumerate}
\end{exercice}

\courssub{Je me prépare au Probatoire}

\begin{exercice}[Probatoire 2022 type : Turboalternateur du barrage de Lom Pangar]
\label{exo:probatoire_lom_pangar}
Le barrage hydroélectrique de Lom Pangar (région de l'Est) est équipé de groupes turboalternateurs Francis. Un des groupes fonctionne aux conditions nominales suivantes :
\begin{itemize}
    \item Puissance mécanique sur l'arbre : $P_m = 30~\text{MW}$.
    \item Tension nominale triphasée aux bornes du stator : $U_n = 13,8~\text{kV}$.
    \item Courant nominal par phase : $I_n = 1250~\text{A}$.
    \item Facteur de puissance (cos $\varphi$) : $0,9$.
    \item Vitesse de rotation synchrone : $n_s = 428,6~\text{tr/min}$.
    \item Nombre de pôles de l'alternateur : $2p = 14$.
\end{itemize}
\begin{enumerate}
    \item Vérifier que la fréquence du réseau électrique camerounais ($f = 50~\text{Hz}$) est cohérente avec la vitesse de rotation et le nombre de pôles donnés.
    \item Calculer la puissance électrique apparente $S_n$ (en MVA) et la puissance active $P_e$ (en MW) délivrées par l'alternateur.
    \item En déduire le rendement global $\eta$ de l'alternateur à la charge nominale.
    \item Calculer l'énergie électrique active $W_e$ (en kWh) produite en une journée de fonctionnement continu à la puissance nominale.
    \item L'excitation du rotor est assurée par un système sans balais (brushless). Expliquer brièvement le principe : comment le courant continu d'excitation est-il généré et amené au rotor sans bagues collectrices ni balais ?
    \item \textbf{Savoir-être :} En tant qu'ingénieur de maintenance à ENEO, vous observez une surchauffe anormale du stator. Citer trois causes possibles liées à l'exploitation (surcharge, déséquilibre, ventilation) et une conséquence grave si rien n'est fait.
\end{enumerate}
\end{exercice}

\begin{exercice}[Probatoire 2021 type : Choix d'un générateur pour une zone rurale isolée]
\label{exo:probatoire_zone_rurale}
Un village isolé de l'Adamaoua (hors réseau interconnecté SONATREL) doit être électrifié. Deux solutions sont étudiées :
\begin{description}
    \item[Solution A :] Groupe électrogène diesel monophasé $10~\text{kVA} / 230~\text{V} / 50~\text{Hz}$, moteur $1500~\text{tr/min}$.
    \item[Solution B :] Microcentrale hydroélectrique au fil de l'eau sur un petit affluent de la Vina. Chute $h = 15~\text{m}$, débit $Q = 0,1~\text{m}^3/\text{s}$. Turbine Banki-Michell couplée à un alternateur synchrone monophasé. Rendement global turbine-alternateur $\eta_g = 0,75$.
\end{description}
\begin{enumerate}
    \item Pour la solution A : Déterminer le nombre de paires de pôles $p$ de l'alternateur pour obtenir $50~\text{Hz}$ à $1500~\text{tr/min}$.
    \item Pour la solution B : Calculer la puissance hydraulique $P_h$, puis la puissance électrique $P_e$ disponible. Comparer aux $10~\text{kVA}$ du groupe diesel (supposer $\cos\varphi = 1$ pour l'alternateur hydro).
    \item Dresser le bilan énergétique annuel (en kWh/an) pour chaque solution, en supposant un fonctionnement $8000~\text{h/an}$ pour le diesel et $6000~\text{h/an}$ pour l'hydro (étiage).
    \item Coût du carburant diesel au village : $850~\text{FCFA/L}$. Consommation spécifique : $0,35~\text{L/kWh}$. Calculer le coût annuel du carburant pour la solution A.
    \item \textbf{Analyse multicritères :} Compléter le tableau comparatif ci-dessous (Avantages / Inconvénients) pour aider le conseil municipal à choisir.
    \begin{center}
    \begin{tabularx}{\linewidth}{|l|X|X|}\hline
    Critère & Solution A (Diesel) & Solution B (Hydro) \\ \hline
    Investissement initial (CAPEX) & & \\ \hline
    Coût d'exploitation (OPEX) & & \\ \hline
    Disponibilité / Fiabilité & & \\ \hline
    Impact environnemental & & \\ \hline
    Maintenance locale & & \\ \hline
    \end{tabularx}
    \end{center}
    \item Conclusion argumentée : quelle solution recommandez-vous pour un horizon de 10 ans ? Justifier.
\end{enumerate}
\end{exercice}

\begin{exercice}[Probatoire 2023 type : Étude complète d'un alternateur synchrone à aimants permanents (PMSG) pour éolienne]
\label{exo:probatoire_eolienne}
On étudie un générateur synchrone à aimants permanents (PMSG) destiné à une petite éolienne à axe horizontal de $5~\text{kW}$ (type installé dans le Grand Nord pour le pompage d'eau).
Données constructeur :
\begin{itemize}
    \item Tension de sortie triphasée (rms, ligne-ligne) à vitesse nominale : $U_n = 400~\text{V}$.
    \item Fréquence nominale : $f_n = 50~\text{Hz}$.
    \item Vitesse de rotation nominale : $n_n = 300~\text{tr/min}$.
    \item Nombre de paires de pôles : $p = 10$.
    \item Résistance statorique par phase (à froid) : $R_s = 0,8~\Omega$.
    \item Flux d'aimantation par pôle (rotor) : $\Phi_f = 0,15~\text{Wb}$.
    \item Nombre de spires en série par phase : $N_s = 120$.
    \item Facteur de bobinage / pas / distribution global : $k_w = 0,95$.
\end{itemize}
Le générateur alimente un redresseur triphasé + onduleur pour injection réseau $230~\text{V} / 50~\text{Hz}$.
\begin{enumerate}
    \item Vérifier la cohérence entre $n_n$, $p$ et $f_n$.
    \item Calculer la f.é.m. de crête $\hat{E}$ induite par phase à vide à la vitesse nominale (formule : $\hat{E} = \omega \cdot N_s \cdot k_w \cdot \Phi_f$ avec $\omega = 2\pi f$). Donner la valeur efficace $E$.
    \item À charge nominale ($P_n = 5~\text{kW}$, $\cos\varphi = 0,9$), le courant de phase est $I_n$. Calculer $I_n$.
    \item En négligeant la réactance synchrone (machine non saillante, $X_s \approx 0$ pour cette estimation), écrire l'équation vectorielle de phase : $\vec{U} = \vec{E} - R_s \vec{I}$. Faire le diagramme de Fresnel (schéma vectoriel) pour un fonctionnement générateur ($\vec{I}$ en opposition de phase avec $\vec{U}$ pour $\varphi=0$, décalé pour $\cos\varphi=0,9$).
    \item Calculer la chute de tension ohmique $R_s I_n$ et comparer à $E$. Commenter la régulation de tension inhérente à cette machine.
    \item \textbf{Chaîne énergétique complète :} Depuis le vent ($v = 10~\text{m/s}$, $\rho_{\text{air}} = 1,2~\text{kg/m}^3$, rotor éolienne $D = 5~\text{m}$, $C_p = 0,4$) jusqu'au réseau $230~\text{V}$. Calculer la puissance éolienne $P_{\text{vent}}$, la puissance mécanique $P_m$, la puissance électrique $P_e$ sortie onduleur (rendement électronique $\eta_{\text{élec}} = 0,95$). Dresser la chaîne avec rendements.
    \item Pourquoi un PMSG est-il préféré à une machine asynchrone (cage d'écureuil) pour cette application éolienne à vitesse variable ? (Indices : pas d'excitation réactive, rendement, couplage direct sans multiplicateur).
\end{enumerate}
\end{exercice}

\newpage

\coursec{Corrigés des exercices}

\begin{corrige}[Exercice 1 — QCM : Rôle et principe]
La bonne réponse est la \textbf{b)}.

Un alternateur est un \cle{convertisseur électromécanique} : il reçoit de l'\cle{énergie mécanique} de rotation sur son arbre (fournie par une turbine, un moteur thermique, des pales d'éolienne...) et la convertit, par \cle{induction électromagnétique}, en \cle{énergie électrique alternative}.

\begin{itemize}
\item a) Faux : c'est la description d'un \textbf{moteur} électrique (conversion électrique $\to$ mécanique), c'est-à-dire le fonctionnement inverse.
\item c) Faux : la conversion chimique $\to$ électrique continue est celle d'une \textbf{pile} ou d'un \textbf{accumulateur}.
\item d) Faux : l'alternateur n'a pas besoin de pile pour fonctionner ; au contraire, c'est lui qui \emph{produit} l'énergie électrique. (Seul le circuit d'excitation de certains alternateurs utilise une source continue auxiliaire, mais ce n'est pas une condition générale de fonctionnement.)
\end{itemize}

\textbf{Point de vigilance :} ne jamais confondre le sens de conversion de l'alternateur (mécanique $\to$ électrique) avec celui du moteur (électrique $\to$ mécanique). C'est la question de cours la plus fréquente au Probatoire sur ce chapitre.
\end{corrige}

\begin{corrige}[Exercice 2 — Vrai ou Faux : Parties et fonctionnement]
\begin{enumerate}
\item \textbf{Faux.} La partie tournante s'appelle le \cle{rotor} ; le \cle{stator} est la partie \emph{fixe} de la machine. Moyen mnémotechnique : \emph{stator} = \emph{statique}.
\item \textbf{Faux.} Dans la grande majorité des alternateurs, la tension alternative est induite dans les bobines du \textbf{stator} (l'induit), car le rotor crée un champ magnétique \emph{tournant} qui balaie les conducteurs fixes. C'est la variation du flux à travers les bobines statoriques qui engendre la f.é.m. induite (loi de Faraday). Placer l'induit au stator évite d'avoir à extraire de fortes puissances par des contacts glissants.
\item \textbf{Vrai.} L'\cle{excitation} (création du champ magnétique inducteur) peut être assurée :
\begin{itemize}
\item par un \textbf{aimant permanent} (petits alternateurs : éoliennes, alternateurs de vélo, microcentrales) ;
\item par un \textbf{électroaimant} alimenté en \textbf{courant continu} (gros turboalternateurs et hydroalternateurs), ce qui permet de \emph{régler} l'intensité du champ, donc la tension de sortie.
\end{itemize}
\item \textbf{Faux.} La fréquence dépend des \emph{deux} grandeurs : $f = p \times n_s$, où $p$ est le nombre de paires de pôles et $n_s$ la vitesse de rotation (en tr/s). À $p$ fixé, si la vitesse de rotation varie, la fréquence varie proportionnellement. C'est pourquoi les groupes électrogènes possèdent un \emph{régulateur de vitesse} pour maintenir $f = \SI{50}{Hz}$.
\end{enumerate}
\end{corrige}

\begin{corrige}[Exercice 3 — Définitions et unités]
\begin{enumerate}
\item \textbf{Fréquence $f$ :} nombre de périodes (cycles complets) du signal par seconde. Unité SI : le \textbf{hertz} (\si{Hz}), avec $\SI{1}{Hz} = \SI{1}{s^{-1}}$. Le réseau camerounais fonctionne à $f = \SI{50}{Hz}$.
\item \textbf{Pulsation $\omega$ :} vitesse angulaire associée au signal sinusoïdal ; elle s'exprime en \textbf{radians par seconde} (\si{rad/s}). Lien avec la fréquence :
\[ \omega = 2\pi f = \frac{2\pi}{T}. \]
\item \textbf{Rendement $\eta$ :} rapport de la puissance électrique utile fournie aux bornes à la puissance mécanique absorbée sur l'arbre :
\[ \eta = \frac{P_{\text{électrique}}}{P_{\text{mécanique}}} \quad (0 < \eta < 1, \text{ grandeur sans unité, souvent exprimée en }\%). \]
\item \textbf{F.é.m. induite :} pour une bobine de $N$ spires de surface $S$ tournant à la pulsation $\omega$ dans un champ magnétique uniforme $B$, le flux magnétique à travers une spire est $\phi(t) = BS\cos(\omega t)$. Le flux enchainé par la bobine est $\Psi(t) = N\phi(t) = NBS\cos(\omega t)$. La loi de Faraday $e = -\dfrac{d\Psi}{dt}$ donne :
\[ e(t) = NBS\omega\,\sin(\omega t) = \hat{E}\,\sin(\omega t), \quad \text{avec } \hat{E} = NBS\omega. \]
Unités : $B$ en tesla (\si{T}), $S$ en \si{m^2}, $\omega$ en \si{rad/s}, $e$ en volt (\si{V}). Vérification dimensionnelle : $\si{T}\cdot\si{m^2}\cdot\si{s^{-1}} = \si{Wb/s} = \si{V}$. \checkmark
\end{enumerate}
\end{corrige}

\begin{corrige}[Exercice 4 — Calculs directs : Fréquence et période]
\textbf{Données :} $n = \SI{3000}{tr/min}$, $p = 2$ paires de pôles.

\begin{enumerate}
\item Conversion de la vitesse : $n_s = \dfrac{3000}{60} = \SI{50}{tr/s}$. La fréquence de la tension produite par un alternateur synchrone est :
\[ f = p \times n_s = 2 \times 50 = \boxed{\SI{100}{Hz}}. \]
\item Période :
\[ T = \frac{1}{f} = \frac{1}{100} = \SI{0,010}{s} = \boxed{\SI{10}{ms}}. \]
\item Pour obtenir $f' = \SI{50}{Hz}$ avec le même alternateur ($p = 2$) :
\[ n'_s = \frac{f'}{p} = \frac{50}{2} = \SI{25}{tr/s} \quad \Longrightarrow \quad n' = 25 \times 60 = \boxed{\SI{1500}{tr/min}}. \]
\end{enumerate}

\textbf{Point de vigilance :} la formule $f = p\,n_s$ exige $n_s$ en \textbf{tours par seconde}. Oublier la conversion tr/min $\to$ tr/s est l'erreur la plus fréquente : elle donne un résultat 60 fois trop grand.
\end{corrige}

\begin{corrige}[Exercice 5 — Schéma de principe et identification]
\begin{enumerate}
\item \textbf{Schéma de principe (vue de face) :}

\begin{popfigure}
\begin{center}
\begin{tikzpicture}[scale=1.0]
% Stator (carcasse)
\draw[thick, popblue] (0,0) circle (2.3);
\draw[thick, popblue] (0,0) circle (2.0);
\node[popblue] at (0,2.55) {\small Stator (bobine fixe)};
% Bobines stator (sections)
\foreach \a in {45,135,225,315}{
  \draw[fill=poporangeL, draw=poporange] (\a:2.15) circle (0.13);
}
% Rotor (aimant)
\draw[thick, fill=popgreenL] (0,0) circle (1.1);
\draw[thick, popdark] (-1.1,0) -- (1.1,0);
\node at (0,0.55) {\textbf{N}};
\node at (0,-0.55) {\textbf{S}};
\node[popdark] at (0,-1.5) {\small Rotor (aimant)};
% Arbre
\draw[fill=popmuted] (0,0) circle (0.18);
\node at (0.25, -0.25) {\small arbre};
% Sens de rotation
\draw[-{Stealth}, thick, poppurple] (1.7,0.9) arc (27:75:1.95);
\node[poppurple] at (1.75,1.35) {\small $\omega$};
% Champ B
\draw[-{Stealth}, very thick, poppink] (0,-0.35) -- (0,0.95);
\node[poppink, right] at (0.08,0.3) {$\vec{B}$};
% Bornes de sortie
\draw[thick] (2.15,0.6) -- (3.0,0.6) node[right] {$A$};
\draw[thick] (2.15,-0.6) -- (3.0,-0.6) node[right] {$B$};
\node at (3.6,0) {\small $u_{AB}(t)$};
\end{tikzpicture}
\end{center}
\legende{Alternateur monophasé à aimant tournant : le champ $\vec{B}$ du rotor balaie les bobines du stator et y induit une tension alternative $u_{AB}$.}
\end{popfigure}

\item \textbf{Sens du champ $\vec{B}$ :} à l'intérieur de l'aimant, le champ va du pôle Sud vers le pôle Nord ; dans l'\cle{entrefer} (espace entre rotor et stator), les lignes de champ vont du \textbf{Nord vers le Sud}. Au moment où le pôle Nord passe devant la bobine, $\vec{B}$ est dirigé du rotor vers le stator, radialement (flèche rose sur le schéma).

\item \textbf{Rôle de chaque partie :}
\begin{itemize}
\item \textbf{Rotor (aimant N/S) :} partie tournante ; il crée le \emph{champ magnétique inducteur} qui, par sa rotation, produit un flux variable à travers la bobine.
\item \textbf{Stator (bobine) :} partie fixe ; c'est l'\emph{induit} : la variation de flux y induit la f.é.m. alternative (loi de Faraday) recueillie entre $A$ et $B$.
\item \textbf{Arbre :} transmet l'\emph{énergie mécanique} de rotation depuis la source primaire (turbine, manivelle, moteur).
\item \textbf{Fils de sortie directs :} ici la bobine est \emph{fixe}, donc les bornes $A$ et $B$ sont reliées directement au circuit extérieur : \textbf{aucune bague collectrice n'est nécessaire}. C'est l'avantage de la configuration « aimant au rotor, bobine au stator » : pas de contact glissant, pas d'usure ni d'étincelles. (Des bagues + balais ne sont requises que si l'on doit alimenter un \emph{rotor bobiné} en courant continu d'excitation.)
\end{itemize}
\end{enumerate}
\end{corrige}

\begin{corrige}[Exercice 6 — Chaîne énergétique : Microcentrale villageoise]
\textbf{Données :} $h = \SI{50}{m}$, $Q = \SI{0,2}{m^3/s}$, $\rho = \SI{1000}{kg/m^3}$, $g = \SI{9,81}{N/kg}$, $\eta_t = 0,85$, $\eta_a = 0,92$.

\begin{enumerate}
\item \textbf{Puissance hydraulique :} la puissance mécanique théorique d'une chute d'eau est $P_h = \rho\, g\, Q\, h$ (énergie potentielle $mgh$ libérée par unité de temps, avec $\dfrac{m}{t} = \rho Q$) :
\[ P_h = 1000 \times 9,81 \times 0,2 \times 50 = \boxed{\SI{98100}{W} \approx \SI{98,1}{kW}}. \]
Vérification dimensionnelle : $\si{kg/m^3}\cdot\si{m/s^2}\cdot\si{m^3/s}\cdot\si{m} = \si{kg.m^2/s^3} = \si{W}$. \checkmark

\item \textbf{Puissance mécanique sur l'arbre} (sortie turbine) :
\[ P_m = \eta_t \times P_h = 0,85 \times 98,1 = \boxed{\SI{83,4}{kW}}. \]

\item \textbf{Puissance électrique aux bornes} (sortie alternateur) :
\[ P_e = \eta_a \times P_m = 0,92 \times 83,4 = \boxed{\SI{76,7}{kW}}. \]
Le rendement global de la centrale est $\eta_g = \eta_t \times \eta_a = 0,85 \times 0,92 = 0,782$, soit environ $78\,\%$.

\item \textbf{Chaîne des transformations d'énergie :}
\begin{center}
\begin{tabularx}{\linewidth}{|X|X|X|X|}\hline
\rowcolor{popblueL} Énergie potentielle de l'eau ($P_h = 98,1$ kW) & Énergie cinétique puis mécanique de rotation — turbine Pelton ($\eta_t = 0,85$) & Énergie mécanique sur l'arbre ($P_m = 83,4$ kW) $\to$ énergie électrique — alternateur ($\eta_a = 0,92$) & Énergie électrique au réseau du village ($P_e = 76,7$ kW) \\ \hline
\end{tabularx}
\end{center}
Les pertes ($21,4$ kW au total) sont dues aux frottements hydrauliques et mécaniques, à l'effet Joule dans les bobinages et aux pertes fer ; elles échauffent la turbine et l'alternateur.
\end{enumerate}
\end{corrige}

\begin{corrige}[Exercice 7 — Alternateur de voiture vs Groupe électrogène]
\begin{enumerate}
\item \textbf{Fréquence de l'alternateur de voiture :} $n_s = \dfrac{2000}{60} = \SI{33,3}{tr/s}$, donc :
\[ f = p \times n_s = 6 \times 33,3 = \boxed{\SI{200}{Hz}}. \]
(On retrouve que les alternateurs automobiles travaillent à des fréquences bien supérieures au $50$ Hz du réseau.)

\item \textbf{Pourquoi redresser et réguler ?} La batterie ($\SI{12}{V}$) est un récepteur à \textbf{courant continu} : elle ne peut pas être chargée par une tension alternative, qui inverserait le sens du courant 200 fois par seconde. On insère donc :
\begin{itemize}
\item un \textbf{pont redresseur} à diodes (triphasé) qui convertit l'alternatif en continu ;
\item un \textbf{régulateur de tension} qui ajuste le courant d'excitation du rotor pour maintenir environ $\SI{14}{V}$ quels que soient le régime moteur et la charge (phares, climatisation...).
\end{itemize}

\item \textbf{Paires de pôles du groupe électrogène :} on veut $f = \SI{50}{Hz}$ avec $n = \SI{3000}{tr/min}$, soit $n_s = \SI{50}{tr/s}$ :
\[ p = \frac{f}{n_s} = \frac{50}{50} = \boxed{1 \text{ paire de pôles} \ (2 \text{ pôles})}. \]

\item \textbf{Comparaison des sources d'énergie mécanique primaire :}
\begin{itemize}
\item \emph{Alternateur de voiture :} entraîné par le \textbf{moteur thermique} du véhicule via la courroie d'accessoires ; l'énergie primaire est l'\emph{énergie chimique du carburant}.
\item \emph{Groupe électrogène :} entraîné par son propre \textbf{moteur thermique} (essence ou diesel) ; l'énergie primaire est également l'énergie chimique du carburant, mais le moteur est \emph{dédié} et régulé en vitesse pour garantir $50$ Hz, alors que l'alternateur de voiture subit les variations du régime moteur (d'où le redressement + régulation).
\end{itemize}
\end{enumerate}
\end{corrige}

\begin{corrige}[Exercice 8 — Réalisation : Alternateur de démonstration]
\begin{enumerate}
\item \textbf{Choix du montage :} on retient la configuration \textbf{aimant tournant (rotor) / bobine fixe (stator)}. L'aimant est monté sur l'axe de la manivelle et tourne devant (ou à l'intérieur de) la bobine fixe, connectée directement à la LED. \emph{Justification :} la bobine étant fixe, on raccorde la LED par de simples fils soudés, sans contacts glissants ; si l'on faisait tourner la bobine, il faudrait des bagues et des balais, délicats à réaliser en classe. De plus, l'aimant cylindrique est facile à centrer et à équilibrer sur un axe.

\begin{popfigure}
\begin{center}
\begin{tikzpicture}[scale=0.95]
% Bobine fixe
\draw[thick, poporange] (-2.2,-1.2) rectangle (-0.4,1.2);
\foreach \y in {-1.0,-0.7,...,1.0}{
  \draw[poporange] (-2.2,\y) -- (-0.4,\y);
}
\node[poporange, below] at (-1.3,-1.3) {\small Bobine fixe ($N$ spires)};
% Aimant tournant
\draw[thick, fill=popgreenL] (1.2,0) circle (0.7);
\draw[thick, popdark] (0.5,0) -- (1.9,0);
\node at (1.2,0.35) {\small \textbf{N}};
\node at (1.2,-0.35) {\small \textbf{S}};
\node[popdark, below] at (1.2,-0.85) {\small Aimant (rotor)};
% Axe + manivelle
\draw[thick] (1.2,0) -- (3.2,0);
\draw[thick] (3.2,0) -- (3.2,0.7) -- (3.7,0.7);
\draw[fill=popmuted] (3.7,0.7) circle (0.09);
\node[right] at (3.75,0.7) {\small manivelle};
\draw[-{Stealth}, thick, poppurple] (2.1,0.75) arc (35:80:0.9);
\node[poppurple] at (2.15,1.15) {\small $\omega$};
% LED
\draw[thick] (-2.2,0.9) -- (-3.2,0.9) -- (-3.2,0.2);
\draw[thick] (-2.2,-0.9) -- (-3.2,-0.9) -- (-3.2,-0.2);
\draw[thick, poppink] (-3.55,0.2) -- (-2.85,0.2);
\draw[thick, poppink] (-3.2,-0.2) -- (-3.2,0.2);
\draw[-{Stealth}, poppink] (-3.35,0.3) -- (-3.55,0.6);
\draw[-{Stealth}, poppink] (-3.1,0.3) -- (-2.9,0.6);
\node[poppink, left] at (-3.6,-0.35) {\small LED};
\end{tikzpicture}
\end{center}
\legende{Alternateur de démonstration : manivelle $\to$ aimant tournant $\to$ flux variable dans la bobine fixe $\to$ tension alternative aux bornes de la LED.}
\end{popfigure}

\item \textbf{Estimation du nombre de spires :} la tension de crête est $\hat{u} = N B S \omega$.
\begin{itemize}
\item Surface utile : $S = \dfrac{\pi d^2}{4} = \dfrac{\pi \times (0,010)^2}{4} = \SI{7,85e-5}{m^2}$.
\item Pulsation : $n = \SI{2}{tr/s}$, donc $\omega = 2\pi n = 4\pi = \SI{12,57}{rad/s}$.
\end{itemize}
On impose $\hat{u} \ge \SI{3}{V}$ :
\[ N \ge \frac{\hat{u}}{B S \omega} = \frac{3}{0,5 \times 7,85\times 10^{-5} \times 12,57} = \frac{3}{4,93\times 10^{-4}} \approx \boxed{6\,100 \text{ spires}}. \]
\emph{Commentaire :} ce nombre élevé montre pourquoi les bobines de démonstration utilisent du fil très fin ($\varnothing\,\SI{0,2}{mm}$) bobiné sur de grandes longueurs. En pratique, on obtient quelques volts avec $2\,000$ à $5\,000$ spires en tournant vite ; un aimant néodyme plus puissant ($B$ jusqu'à $\SI{1}{T}$) divise $N$ par deux.

\item \textbf{LED et alternatif :} une LED ne conduit que dans \emph{un seul sens} (diode) : pendant une alternance sur deux, la tension est inversée et la LED est bloquée — elle s'éteint donc la moitié du temps (clignotement à la fréquence de rotation). \textbf{Solution sans condensateur :} insérer un \textbf{pont de Graëtz} (4 diodes) qui redresse la tension en double alternance : la LED reçoit alors un courant toujours dans le même sens, aux deux alternances, et reste allumée en quasi-continu. (On peut aussi monter deux LED tête-bêche : chacune éclaire sur une alternance.)
\end{enumerate}
\end{corrige}

\begin{corrige}[Exercice 9 — Probatoire 2022 type : Turboalternateur du barrage de Lom Pangar]
\textbf{Barème indicatif (sur 10 pts) :} 1. (1,5 pt) ; 2. (2 pts) ; 3. (1,5 pt) ; 4. (1,5 pt) ; 5. (2 pts) ; 6. (1,5 pt).

\begin{enumerate}
\item \textbf{Cohérence fréquence / vitesse / pôles :} $2p = 14$ pôles donc $p = 7$ paires de pôles ; $n_s = \dfrac{428,6}{60} = \SI{7,143}{tr/s}$.
\[ f = p \times n_s = 7 \times 7,143 = \SI{50,0}{Hz}. \]
La fréquence calculée coïncide avec la fréquence du réseau camerounais : les données sont \textbf{cohérentes}. \checkmark

\item \textbf{Puissances apparente et active :} en triphasé :
\begin{align*}
S_n &= \sqrt{3}\, U_n I_n = \sqrt{3} \times 13\,800 \times 1250 = \SI{2,988e7}{VA} \approx \boxed{\SI{29,9}{MVA}},\\
P_e &= S_n \cos\varphi = 29,9 \times 0,9 \approx \boxed{\SI{26,9}{MW}}.
\end{align*}

\item \textbf{Rendement de l'alternateur :}
\[ \eta = \frac{P_e}{P_m} = \frac{26,9}{30} = 0,897 \approx \boxed{89,7\,\%}. \]
Valeur réaliste pour un hydroalternateur de cette taille (les pertes : Joule au stator, pertes fer, frottements, ventilation).

\item \textbf{Énergie produite en 24 h :}
\[ W_e = P_e \times \Delta t = 26,9~\text{MW} \times 24~\text{h} = 645,6~\text{MWh} = \boxed{6,46\times 10^{5}~\text{kWh}}. \]
(Soit environ $646$ MWh par jour et par groupe.)

\item \textbf{Excitation sans balais (brushless) :} le principe est le suivant : une petite \emph{excitatrice} — un alternateur auxiliaire dont l'induit est \emph{au rotor} — est montée sur le même arbre. Son champ statorique est alimenté en continu réglé ; l'induit tournant de l'excitatrice produit une tension \emph{triphasée alternative}, immédiatement \textbf{redressée par un pont de diodes tournant} (monté sur l'arbre). Le courant continu ainsi obtenu alimente \emph{directement} le bobinage inducteur du rotor principal, par connexions solidaires de l'arbre. Ainsi, aucune énergie ne transite par des contacts glissants : \textbf{ni bagues, ni balais}, donc moins d'usure, pas d'étincelles et une maintenance réduite — un atout décisif pour un site comme Lom Pangar.

\item \textbf{Surchauffe du stator — diagnostic (savoir-être) :} trois causes possibles :
\begin{itemize}
\item \textbf{surcharge} durable ($I > I_n$) : les pertes Joule $R I^2$ croissent comme le carré du courant ;
\item \textbf{déséquilibre des phases} (charge ou connexion défaillante) : circulation de courants homopolaires/inverses échauffant anormalement les enroulements ;
\item \textbf{défaut de ventilation/refroidissement} (filtres encrassés, ventilateur en panne, circuit d'eau obstrué).
\end{itemize}
\textbf{Conséquence grave si rien n'est fait :} la dégradation de l'\emph{isolation} des enroulements statoriques, pouvant conduire à un \textbf{court-circuit entre spires ou à la masse}, à la destruction du groupe et à un arrêt long et coûteux de la production (délestage sur le réseau ENEO).
\end{enumerate}

\textbf{Points de vigilance (rédaction Probatoire) :}
\begin{itemize}
\item Citer la formule $f = p\,n_s$ \emph{avant} l'application numérique, avec $n_s$ en tr/s.
\item En triphasé, ne pas oublier le facteur $\sqrt{3}$ dans $S = \sqrt{3}\,UI$ (et ne pas confondre avec $3\,U_{\text{phase}} I$).
\item Préciser que $\eta < 1$ et commenter physiquement les pertes.
\item Pour l'énergie, vérifier la cohérence des unités : MW $\times$ h = MWh, puis convertir en kWh ($1$ MWh $= 10^3$ kWh).
\end{itemize}
\end{corrige}

\begin{corrige}[Exercice 10 — Probatoire 2021 type : Choix d'un générateur pour une zone rurale isolée]
\textbf{Barème indicatif (sur 12 pts) :} 1. (1 pt) ; 2. (2 pts) ; 3. (2 pts) ; 4. (2 pts) ; 5. (3 pts) ; 6. (2 pts).

\begin{enumerate}
\item \textbf{Solution A — paires de pôles :} $n_s = \dfrac{1500}{60} = \SI{25}{tr/s}$ ;
\[ p = \frac{f}{n_s} = \frac{50}{25} = \boxed{2 \text{ paires de pôles} \ (4 \text{ pôles})}. \]

\item \textbf{Solution B — puissances :}
\[ P_h = \rho g Q h = 1000 \times 9,81 \times 0,1 \times 15 = \SI{14715}{W} \approx \boxed{\SI{14,7}{kW}}. \]
\[ P_e = \eta_g \times P_h = 0,75 \times 14,7 \approx \boxed{\SI{11,0}{kW}}. \]
Comparaison : avec $\cos\varphi = 1$, les $10$ kVA du groupe diesel correspondent à $10$ kW actifs. L'hydro ($11,0$ kW) fournit donc une puissance \emph{légèrement supérieure} au diesel : les deux solutions couvrent le besoin de $10$ kW.

\item \textbf{Bilan énergétique annuel :}
\begin{align*}
W_A &= P_A \times t_A = 10~\text{kW} \times 8000~\text{h} = \boxed{80\,000~\text{kWh/an}},\\
W_B &= P_B \times t_B = 11,0~\text{kW} \times 6000~\text{h} = \boxed{66\,200~\text{kWh/an}}.
\end{align*}
Le diesel produit davantage sur l'année (plus d'heures de fonctionnement), mais au prix d'une consommation de carburant ; l'hydro dépend de la saison (étiage).

\item \textbf{Coût annuel du carburant (solution A) :}
\begin{align*}
\text{Volume} &= 80\,000~\text{kWh} \times 0,35~\text{L/kWh} = 28\,000~\text{L},\\
\text{Coût} &= 28\,000 \times 850 = \boxed{23\,800\,000~\text{FCFA/an}}
\end{align*}
soit environ \textbf{23,8 millions de FCFA par an}, sans compter le transport du gasoil jusqu'au village, la maintenance et les lubrifiants.

\item \textbf{Tableau multicritères complété :}
\begin{center}
\begin{tabularx}{\linewidth}{|l|X|X|}\hline
\rowcolor{popblueL} Critère & Solution A (Diesel) & Solution B (Hydro) \\ \hline
Investissement initial (CAPEX) & Faible à modéré (groupe standard, pose rapide) & Élevé (génie civil : prise d'eau, canal, bâche, local) \\ \hline
Coût d'exploitation (OPEX) & Très élevé : $\approx 23,8$ millions FCFA/an de gasoil + entretien moteur & Très faible : eau gratuite, entretien courant limité \\ \hline
Disponibilité / Fiabilité & Bonne si approvisionnement en carburant assuré (piste difficile en saison des pluies) ; pannes moteur possibles & Dépend du débit de la rivière (baisse à l'étiage) ; machine simple et robuste \\ \hline
Impact environnemental & Émissions de CO$_2$, bruit, risques de fuites d'huile et de gasoil & Quasi nul au fil de l'eau (pas de barrage), pas de combustion \\ \hline
Maintenance locale & Mécanicien diesel disponible localement, mais pièces et carburant à approvisionner & Mécanique simple (turbine Banki), réparable localement ; compétences à former \\ \hline
\end{tabularx}
\end{center}

\item \textbf{Conclusion argumentée (horizon 10 ans) :} je recommande la \textbf{solution B (microcentrale hydro)}. Sur 10 ans, le seul carburant du diesel coûterait $\approx 238$ millions FCFA, somme très supérieure au surcoût d'investissement initial de l'hydro. L'hydro offre une énergie \emph{locale, renouvelable, sans émission}, à coût marginal nul, et sa technologie (turbine Banki-Michell) est robuste et réparable sur place. Ses limites — dépendance à l'étiage et CAPEX élevé — peuvent être levées par un \emph{petit groupe diesel de secours} conservé pour la saison sèche et les pointes : c'est la solution hybride la plus soutenable pour le village.
\end{enumerate}

\textbf{Points de vigilance :} distinguer clairement puissance (kW) et énergie (kWh) ; justifier toute hypothèse ($\cos\varphi = 1$) ; en question de synthèse, structurer la réponse (technique, économique, environnemental) et \emph{conclure explicitement}.
\end{corrige}

\begin{corrige}[Exercice 11 — Probatoire 2023 type : PMSG pour éolienne]
\textbf{Barème indicatif (sur 14 pts) :} 1. (1 pt) ; 2. (2 pts) ; 3. (1,5 pt) ; 4. (2,5 pts) ; 5. (1,5 pt) ; 6. (3 pts) ; 7. (2,5 pts).

\begin{enumerate}
\item \textbf{Cohérence $n_n$, $p$, $f_n$ :} $n_s = \dfrac{300}{60} = \SI{5}{tr/s}$ ;
\[ f = p \times n_s = 10 \times 5 = \SI{50}{Hz} = f_n. \quad \checkmark \]
Les données sont cohérentes. (Le grand nombre de pôles permet une vitesse de rotation \emph{lente} : couplage direct à l'éolienne, sans multiplicateur.)

\item \textbf{F.é.m. induite par phase à vide :} $\omega = 2\pi f = 2\pi \times 50 = \SI{314,2}{rad/s}$ ;
\[ \hat{E} = \omega\, N_s\, k_w\, \Phi_f = 314,2 \times 120 \times 0,95 \times 0,15 = \boxed{\SI{5,37e3}{V} \approx \SI{5,4}{kV} \text{ (crête)}}. \]
Valeur efficace :
\[ E = \frac{\hat{E}}{\sqrt{2}} = \frac{5372}{\sqrt{2}} \approx \boxed{\SI{3,80e3}{V} \approx \SI{3,8}{kV}}. \]
\emph{Commentaire :} cette f.é.m. à vide, calculée avec les données constructeur idéalisées, est très supérieure à la tension nominale $U_n = \SI{400}{V}$ ; en machine réelle, le flux utile embrassé par les bobinages est bien inférieur au flux total par pôle (fuites magnétiques, entrefer) et la formule donne une borne supérieure. On retiendra la \emph{méthode} : $\hat{E} = \omega N_s k_w \Phi_f$, puis $E = \hat{E}/\sqrt{2}$.

\item \textbf{Courant nominal de phase :} en triphasé, $P_n = \sqrt{3}\,U_n I_n \cos\varphi$, d'où :
\[ I_n = \frac{P_n}{\sqrt{3}\,U_n \cos\varphi} = \frac{5000}{\sqrt{3} \times 400 \times 0,9} = \frac{5000}{623,5} \approx \boxed{\SI{8,0}{A}}. \]

\item \textbf{Équation de phase et diagramme de Fresnel :} en convention générateur, pour une phase :
\[ \vec{U} = \vec{E} - R_s\,\vec{I} \qquad \Longleftrightarrow \qquad \vec{E} = \vec{U} + R_s\,\vec{I}. \]
La f.é.m. $\vec{E}$ est la somme vectorielle de la tension aux bornes $\vec{U}$ et de la chute ohmique $R_s\vec{I}$ (colinéaire à $\vec{I}$). Le courant est déphasé de $\varphi = \arccos(0,9) \approx 25,8^\circ$ par rapport à $\vec{U}$.

\begin{popfigure}
\begin{center}
\begin{tikzpicture}[scale=1.1, >=Stealth]
\coordinate (O) at (0,0);
\coordinate (U) at (4,0);
\coordinate (I) at (25.8:1.6);
\coordinate (E) at ($(U)+(I)$);
\draw[->, very thick, popblue] (O) -- (U) node[midway, below] {$\vec{U}$};
\draw[->, very thick, poppink] (O) -- (I) node[midway, above left] {$\vec{I}$};
\draw[->, very thick, poporange] (U) -- (E) node[midway, right] {$R_s\vec{I}$};
\draw[->, very thick, popgreen] (O) -- (E) node[midway, above] {$\vec{E}$};
\draw[dashed, popmuted] (I) -- (E);
\draw[popdark] (0.8,0) arc (0:25.8:0.8);
\node[popdark] at (14:1.15) {$\varphi$};
\end{tikzpicture}
\end{center}
\legende{Diagramme de Fresnel d'une phase du PMSG en générateur : $\vec{E} = \vec{U} + R_s\vec{I}$, avec $\varphi \approx 25,8^\circ$ ($\cos\varphi = 0,9$).}
\end{popfigure}

\item \textbf{Chute de tension ohmique :}
\[ R_s I_n = 0,8 \times 8,0 = \boxed{\SI{6,4}{V}}. \]
Comparée à $U_n = \SI{400}{V}$, elle ne représente que $\dfrac{6,4}{400} \approx 1,6\,\%$ : la chute ohmique est \textbf{négligeable}. La tension de cette machine à aimants permanents est donc naturellement \emph{peu sensible à la charge} (bonne régulation intrinsèque côté résistance) ; en revanche, comme l'excitation n'est \emph{pas réglable} (aimants permanents), la tension varie avec la \emph{vitesse du vent} : c'est l'électronique de puissance (redresseur + onduleur) qui assure la régulation finale à $230$ V / $50$ Hz.

\item \textbf{Chaîne énergétique complète :}
\begin{itemize}
\item Puissance du vent traversant le disque du rotor ($A = \dfrac{\pi D^2}{4} = \dfrac{\pi \times 25}{4} = \SI{19,6}{m^2}$) :
\[ P_{\text{vent}} = \frac{1}{2}\rho A v^3 = 0,5 \times 1,2 \times 19,6 \times 10^3 = \boxed{\SI{11,8}{kW}}. \]
\item Puissance mécanique extraite (coefficient de puissance $C_p = 0,4$, limite de Betz $= 0,593$ non dépassée \checkmark) :
\[ P_m = C_p \times P_{\text{vent}} = 0,4 \times 11,8 = \boxed{\SI{4,71}{kW}}. \]
\item Puissance électrique injectée au réseau (rendement de la chaîne générateur + électronique pris égal à $\eta_{\text{élec}} = 0,95$) :
\[ P_e = \eta_{\text{élec}} \times P_m = 0,95 \times 4,71 = \boxed{\SI{4,48}{kW}}. \]
\end{itemize}
\begin{center}
\begin{tabularx}{\linewidth}{|X|X|X|}\hline
\rowcolor{popblueL} Énergie cinétique du vent $P_{\text{vent}} = 11,8$ kW & Rotor éolien ($C_p = 0,4$) $\to$ énergie mécanique $P_m = 4,71$ kW & PMSG + redresseur + onduleur ($\eta = 0,95$) $\to$ énergie électrique $P_e = 4,48$ kW au réseau $230$ V \\ \hline
\end{tabularx}
\end{center}
Rendement global vent $\to$ réseau : $\eta_g = 0,4 \times 0,95 = 0,38$, soit $38\,\%$ — valeur typique d'une petite éolienne.

\item \textbf{Pourquoi un PMSG plutôt qu'une machine asynchrone ?}
\begin{itemize}
\item \textbf{Pas d'excitation réactive :} la machine asynchrone à cage \emph{consomme} de la puissance réactive pour se magnétiser ; en site isolé ou à vitesse variable, il faut des condensateurs ou une électronique coûteuse. Le PMSG crée son propre champ grâce aux aimants : aucun courant magnétisant n'est nécessaire.
\item \textbf{Rendement supérieur :} pas de courants rotoriques, donc pas de pertes Joule au rotor ; le rendement est meilleur, surtout à charge partielle (vent faible — situation fréquente).
\item \textbf{Couplage direct :} avec $p = 10$ paires de pôles, le $50$ Hz (ou la fréquence variable adaptée) est obtenu à \emph{faible vitesse} ($300$ tr/min) : on supprime le \emph{multiplicateur} mécanique, source de pannes, de bruit et de maintenance — un atout décisif pour une installation isolée du Grand Nord.
\item \textbf{Contrainte acceptée :} la tension et la fréquence variables sont gérées par le redresseur + onduleur, qui découple la machine du réseau.
\end{itemize}
\end{enumerate}

\textbf{Points de vigilance :} dans $P_{\text{vent}} = \tfrac{1}{2}\rho A v^3$, la vitesse est au \emph{cube} (erreur classique : oublier le cube ou utiliser le diamètre au lieu de l'aire) ; vérifier que $C_p \le 0,593$ (limite de Betz) ; en triphasé, bien préciser si la tension est composée (ligne-ligne) ou simple ; pour le diagramme de Fresnel, écrire d'abord l'équation de maille en convention générateur avant de tracer.
\end{corrige}

\newpage

\coursec{Fiche bilan}

\begin{popfigure}
\begin{tikzpicture}[
    mindmap,
    grow cyclic,
    every node/.style={concept, execute at begin node=\hskip0pt},
    concept color=popblue,
    level 1/.append style={level distance=4.5cm, sibling angle=360/6},
    level 2/.append style={level distance=3cm, sibling angle=360/6},
    text=white,
    font=\small\bfseries
]
\node[concept, scale=1.2] {L'alternateur}
    child[concept color=poporange] { node[concept, align=center]{Rôle\\Production $e(t)$} }
    child[concept color=popgreen] { node[concept, align=center]{Principe\\Induction\\électromagnétique} }
    child[concept color=poppurple] { node[concept, align=center]{Parties essentielles\\Stator \textbullet{} Rotor\\Anneaux \textbullet{} Balais} }
    child[concept color=poppink] { node[concept, align=center]{Transformations\\Mécanique $\to$ Électrique} }
    child[concept color=popgold] { node[concept, align=center]{Caractéristiques\\$f = \frac{p n}{60}$\\$e = E_{\max}\sin\omega t$} }
    child[concept color=popdark] { node[concept, align=center]{Exemples\\Barrages (Lom Pangar)\\Groupes électrogènes} };
\end{tikzpicture}
\legende{Carte mentale de la leçon 17 : L'alternateur}
\end{popfigure}

\begin{bilan}

\textbf{Définitions clés}
\begin{itemize}
    \item \cle{Alternateur} : Machine électromécanique transformant l'énergie mécanique de rotation en énergie électrique alternative (courant et tension sinusoïdaux).
    \item \cle{Force électromotrice (f.e.m.) alternative} : Tension $e(t) = E_{\max} \sin(\omega t + \varphi_0)$ produite aux bornes de l'induit (stator).
    \item \cle{Fréquence $f$} : Nombre de périodes par seconde (Hz). Liée à la vitesse de rotation $n$ (tr/min) et au nombre de paires de pôles $p$ par $f = \dfrac{p\, n}{60}$.
    \item \cle{Période $T$} : $T = \dfrac{1}{f}$.
    \item \cle{Pulsation $\omega$} : $\omega = 2\pi f = \dfrac{2\pi}{T}$ (rad$\cdot$s$^{-1}$).
\end{itemize}

\textbf{Formules et lois à connaître par cœur}
\begin{center}
\begin{tabularx}{\linewidth}{|l|X|c|}\hline
\textbf{Grandeur} & \textbf{Relation} & \textbf{Unité SI} \\ \hline
Fréquence & $f = \dfrac{p\, n}{60}$ & Hz \\ \hline
Pulsation & $\omega = 2\pi f$ & rad$\cdot$s$^{-1}$ \\ \hline
f.e.m. instantanée & $e(t) = E_{\max} \sin(\omega t)$ & V \\ \hline
Tension efficace & $E = \dfrac{E_{\max}}{\sqrt{2}}$ & V \\ \hline
Puissance mécanique fournie & $P_{\text{méca}} = \mathcal{C}\, \omega$ & W \\ \hline
Rendement & $\eta = \dfrac{P_{\text{élec}}}{P_{\text{méca}}}$ & -- \\ \hline
\end{tabularx}
\end{center}

\textbf{Méthodes express}
\begin{enumerate}
    \item \textbf{Identifier les parties} : Stator (induit, bobines fixes) \textbullet{} Rotor (inducteur, aimant/électroaimant tournant) \textbullet{} Anneaux collecteurs + balais (liaison tournante).
    \item \textbf{Tracer le schéma de principe} : Représenter le rotor (aimant NS) tournant devant les bobines du stator ; indiquer le sens de rotation, les bornes de sortie $A,B$, la f.e.m. $e(t)$ sinusoïdale.
    \item \textbf{Calculer $f$ et $T$} : Compter les paires de pôles $p$ sur le rotor $\to$ relever $n$ (tr/min) $\to$ appliquer $f = p n / 60$ $\to$ $T = 1/f$.
    \item \textbf{Bilan d'énergie} : Énergie mécanique (moteur thermique, turbine hydraulique) $\to$ Énergie électromagnétique (induction) $\to$ Énergie électrique alternative (réseau/usagers). Pertes : frottements (mécaniques), effet Joule (cuivre), hystérésis/Foucault (fer).
\end{enumerate}

\end{bilan}

\begin{aretenir}[Checklist avant le Probatoire]
\begin{itemize}
    \item $\square$ Définir l'alternateur et donner son rôle principal (production courant alternatif).
    \item $\square$ Citer les 4 parties essentielles : stator, rotor, anneaux, balais ; expliquer le rôle de chacune.
    \item $\square$ Expliquer le principe : induction électromagnétique (loi de Faraday-Lenz), variation du flux $\Phi(B,S,\theta)$.
    \item $\square$ Écrire l'expression de la f.e.m. $e(t) = E_{\max}\sin(\omega t)$ et identifier $E_{\max}$, $\omega$, $f$, $T$.
    \item $\square$ Établir et utiliser la relation $f = \dfrac{p\, n}{60}$ (vérifier homogénéité : tr/min $\to$ s$^{-1}$).
    \item $\square$ Tracer le schéma de principe d'un alternateur monophasé simple (aimant tournant, bobine fixe, anneaux, balais, charge).
    \item $\square$ Décrire la chaîne de transformations d'énergie : Mécanique $\to$ Électrique (citer les pertes principales).
    \item $\square$ Citer 2 exemples concrets au Cameroun : alternateurs des barrages (Lom Pangar, Nachtigal, Edéa) et groupes électrogènes (hôpitaux, zones rurales).
    \item $\square$ Différencier alternateur (courant alternatif) et dynamo (courant continu via collecteur à segments).
    \item $\square$ Savoir réaliser un alternateur simple en TP : aimant tournant (perceuse/manivelle) + bobine + galvanomètre/oscilloscope $\to$ observer signal sinusoïdal.
    \item $\square$ Préciser que la tension aux bornes d'un récepteur connecté n'est pas exactement $e(t)$ cause résistance interne $r$ : $u = e - r i$.
    \item $\square$ Respecter les unités : $n$ en tr/min, $f$ en Hz, $\omega$ en rad/s, $E$ en V, $\mathcal{C}$ en N$\cdot$m.
\end{itemize}
\end{aretenir}

\findecours

\end{document}
